\documentclass[11pt]{article}
\usepackage[margin=1in]{geometry}
\usepackage{amsmath,amssymb,amsthm}
\usepackage[hidelinks]{hyperref}
% These environments also work when the document class is changed to IEEEtran.
\newtheorem{theorem}{Theorem}
\newtheorem{lemma}[theorem]{Lemma}
\newtheorem{corollary}[theorem]{Corollary}
\newtheorem{proposition}[theorem]{Proposition}
\DeclareMathOperator{\tr}{tr}
\DeclareMathOperator{\diag}{diag}
\DeclareMathOperator{\spec}{spec}
\newcommand{\R}{\mathbb R}
\newcommand{\E}{\mathbb E}
\newcommand{\Prob}{\mathbb P}
\newcommand{\ind}{\mathbf 1}
\title{Fractional Group Selection by Node Grounding\\Theory and Oracle Analysis: T1--T12}
\author{Working theory notes}
\date{November 8, 2026}
\begin{document}
\maketitle

\section{Setting and admissible sets}
Let $G=(V,E,w)$ be a finite undirected graph, with positive weights on
its edges, $n=|V|$, Laplacian $L$, and $M=L+\sigma I$, $\sigma\geq0$.
For $0<\beta<\infty$, write
\[
 A_S^{\beta}=M+\beta\diag(\ind_S),\qquad
 A_S^{\infty}=M[V\setminus S,V\setminus S].
\]
The trace of any function of the empty matrix is zero. If $\sigma>0$,
all sets are admissible, including the empty set. If $\sigma=0$, assume
that $G$ is connected and fix a nonempty seed set $S_0$; the admissible
lattice is $\{S:S_0\subseteq S\subseteq V\}$. The corresponding matrices
of positive dimension are positive definite: their quadratic forms are
the edge energy plus the diagonal penalty, and zero energy forces a
constant vector, which the penalty (or the Dirichlet seed) forces to zero.
For the seeded formulation the normalized gain is
$\Phi(S_0)-\Phi(S_0\cup U)$, $U\subseteq V\setminus S_0$.

A completely monotone (CM) function here is finite on $(0,\infty)$ and has
the representation
\begin{equation}\label{eq:cm}
 f(x)=\int_{[0,\infty)}e^{-tx}\,\mu(dt),\qquad \mu\geq0.
\end{equation}
Set $a=\mu(\{0\})=\lim_{x\to\infty}f(x)$; the equality follows by dominated
convergence using $e^{-t x_0}$ for any $x_0>0$.
All objective values $\Phi_f^\beta(S)=\tr f(A_S^\beta)$ are finite on the
admissible lattice. We use the diminishing-drop convention for
supermodularity:
\begin{equation}\label{eq:super}
 \Phi(S)-\Phi(S\cup\{v\})\ \geq\
 \Phi(T)-\Phi(T\cup\{v\}),\quad S\subseteq T,\ v\notin T.
\end{equation}

\section{T1: structural theorem}
\begin{lemma}[One coupling for all grounding sets]\label{lem:coupling}
There is a coupling, independent of the selected set, with killing times
$\zeta_S$ such that $\zeta_{S\cup T}=\min(\zeta_S,\zeta_T)$.
For soft grounding,
\[
 (e^{-tA_S^\beta})_{ij}=\Prob_i(X_t=j,\ t<\zeta_S),\quad t\geq0.
\]
For hard grounding the same identity holds after embedding
$e^{-tA_S^\infty}$ into an $n\times n$ matrix by zeros in rows and
columns in $S$, including at $t=0$.
\end{lemma}
\begin{proof}
Let $X$ be the continuous-time walk on $V$ with jump rate $w_{ij}$ from
$i$ to $j$, holding rate $d_i=\sum_j w_{ij}$, and generator $-L$.
It exists for all time since the state space and all rates are finite;
an isolated state has infinite holding time. Independently sample
$E_0\sim\operatorname{Exp}(\sigma)$ (take $E_0=\infty$ if $\sigma=0$),
and $E_v\sim\operatorname{Exp}(\beta)$ independently for every $v$.
Define occupation times $\ell_v(t)=\int_0^t\ind_{\{X_u=v\}}du$ and
\[
 \kappa_v=\inf\{t:\ell_v(t)\geq E_v\},\qquad
 \zeta_S=E_0\wedge\min_{v\in S}\kappa_v.
\]
These clocks run in local occupation time. In particular an unvisited
vertex cannot kill the walk. The displayed minimum identity holds
pathwise, simultaneously for all sets. Conditioning on the walk gives
\[
 \Prob(t<\zeta_S\mid X)=
 e^{-\sigma t}\exp\left(-\beta\sum_{v\in S}\ell_v(t)\right).
\]
The exponential memoryless property implies that while at $i$ the
additional killing rate is $\sigma+\beta\ind_S(i)$. The transient
transition matrix $K_S(t)$ therefore has initial value $I$ and generator
$Q_{ij}=w_{ij}$ for $i\ne j$, $Q_{ii}=-d_i-\sigma-\beta\ind_S(i)$.
The Markov property and the first-step expansion give $K'_S(t)=K_S(t)Q$;
uniqueness of this finite linear ODE yields $K_S(t)=e^{-tA_S^\beta}$.

For hard grounding replace $\kappa_v$ by the hitting time
$H_v=\inf\{t\geq0:X_t=v\}$ and keep the same $E_0$ and base walk.
On $U=V\setminus S$ the transient rates are $w_{ij}$, $i,j\in U$, $i\ne j$,
and diagonal $-d_i-\sigma$; jumps into $S$ cause absorption.
Thus the transient generator is $-M[U,U]$. A start in $S$ has
$\zeta_S=0$ and contributes zero even at $t=0$. A surviving walk cannot
end in $S$. These facts give precisely the zero embedding and prove
the identity for every $i,j$. All the constructions use one base walk
and the same independent clocks for every set.
\end{proof}
\noindent Numerical check: \texttt{scripts/verify\_t1\_coupling.py}.

\begin{theorem}[CM grounding traces]\label{thm:cm}
For every CM $f$ and $\beta\in(0,\infty]$, $\Phi_f^\beta$ is
nonincreasing and supermodular on the admissible lattice.
For each admissible $S$ and $v\notin S$, the drop is strictly positive
in the soft case exactly when $f$ is nonconstant, and in the hard case
exactly when $f\not\equiv0$.
The soft-to-hard limit is
\begin{equation}\label{eq:hardlimit}
 \lim_{\beta\to\infty}\Phi_f^\beta(S)
 =\Phi_f^\infty(S)+|S|a.
\end{equation}
In particular, the two objectives agree in the limit when $f(\infty)=0$.
\end{theorem}
\begin{proof}
Fix $i,t$ and one coupled outcome. Put
$Z_S=\ind_{\{X_t=i\}}\ind_{\{t<\zeta_S\}}$ for the walk started at $i$.
For $S\subseteq T$, $Z_S\geq Z_T$. Moreover, for the added vertex clock
$\kappa_v$ (or hitting time in the hard case),
\[
 Z_S-Z_{S\cup\{v\}}
 =\ind_{\{X_t=i\}}\ind_{\{t<\zeta_S\}}\ind_{\{\kappa_v\leq t\}}.
\]
Replacing $S$ by $T$ only decreases this product. Taking expectations
and summing over all $i\in V$ proves monotonicity and
\eqref{eq:super} for $h_S(t)=\tr e^{-tA_S^\beta}$, using the zero
embedding in the hard case. This argument includes $t=0$: then the
soft heat trace is $n$ and the hard heat trace is $n-|S|$.

Diagonalizing each positive definite matrix and using Tonelli's theorem
for the nonnegative scalar integrands gives
\[
 \Phi_f^\beta(S)=\int_{[0,\infty)}h_S(t)\,\mu(dt)
 =\sum_i\int_{[0,\infty)}\E_i Z_S\,\mu(dt).
\]
The measure is sigma-finite, since $\mu([0,j])\leq e^j f(1)$.
Every integral is finite by the finite spectral sum. Differences can
therefore be taken, or equivalently Tonelli can be applied directly
to the nonnegative drops. The pointwise drop inequalities integrate
to the claimed inequalities without subtracting infinities.

For strictness, start at $v\notin S$. The event of no jump before $t$
and survival of base killing has probability $e^{-(d_v+\sigma)t}$.
For soft grounding it contributes
$e^{-(d_v+\sigma)t}(1-e^{-\beta t})>0$ to the heat-trace drop for
every $t>0$. Thus any mass of $\mu$ on $(0,\infty)$ gives a strict drop;
if there is no such mass $f=a$ and the soft objective is $na$.
The representation shows that $f$ is constant exactly in this latter
case. For hard grounding the no-jump event contributes
$e^{-(d_v+\sigma)t}>0$, and the drop at $t=0$ is one. Hence every
nonzero $\mu$ gives a strict drop; $\mu=0$ gives the zero objective.

For the limit use common unit exponential variables and set
$E_v^{\beta}=E_v^1/\beta$. For fixed $t>0$, the conditional survival
factor converges to the indicator that no selected vertex has positive
occupation time before $t$. Apart from a null event of a jump exactly
at $t$, this is hard survival; a start in $S$ also has positive
occupation time. Bounded convergence gives
$h_S^\beta(t)\to h_S^\infty(t)$ for $t>0$.
At $t=0$ the limit remains $n$, whereas the hard trace is $n-|S|$.
Choose a fixed $\beta_0>0$. For $\beta\geq\beta_0$,
$0\leq h_S^\beta(t)\leq h_S^{\beta_0}(t)$, an integrable bound by
admissibility. Dominated convergence proves \eqref{eq:hardlimit},
including $S=V$.
\end{proof}
\noindent Numerical check: \texttt{scripts/verify\_t1\_structure.py}.

\begin{corollary}[Powers, heat, and current-flow closeness]\label{cor:instances}
The theorem applies to $f(x)=x^{-s}$ for every $s>0$, to
$f(x)=e^{-\tau x}$ for $\tau\geq0$, and to $(x+c)^{-1}$ for $c\geq0$.
For a connected graph, $\sigma=0$, and $\varnothing\ne S\subsetneq V$,
the group current-flow closeness is
$\operatorname{CFCC}(S)=n/\tr L_{-S}^{-1}$.
Thus the reciprocal-CFCC supermodularity of Li et al.~\cite{li2019}
is the hard-grounding $s=1$ case.
\end{corollary}
\begin{proof}
The positive representations are, respectively,
\[
 x^{-s}=\frac1{\Gamma(s)}\int_0^\infty t^{s-1}e^{-tx}dt,
 \qquad e^{-\tau x}=\int e^{-tx}\delta_\tau(dt),
 \qquad (x+c)^{-1}=\int_0^\infty e^{-tx}e^{-ct}dt.
\]
The first follows by $u=tx$ in the defining Gamma integral, and the
third is an elementary exponential integral. Apply Theorem~\ref{thm:cm}.
For the electrical identity, ground $S$ and inject a unit current at
$u\notin S$. Kirchhoff's equations on $U=V\setminus S$ are
$L[U,U]p=e_u$. Thus the source voltage, equal to its effective
resistance to $S$, is $(L[U,U]^{-1})_{uu}$. Summing these voltages
and using the definition of group current-flow closeness gives the
formula. Multiplication by $1/n$ preserves supermodularity of the
reciprocal objective.
\end{proof}
\noindent Numerical check: \texttt{scripts/verify\_t1\_corollaries.py}.

\begin{proposition}[Edge addition has a different structure]\label{prop:edge}
Even for $f(x)=x^{-1}$, edge addition can violate supermodularity on a
connected five-vertex graph.
\end{proposition}
\begin{proof}
Take the unit-weight path $2-1-0-4-3$, $\sigma=1/10$, and candidate
unit edges $e=\{0,2\}$, $g=\{0,3\}$. Write
$H(B)=\tr(L+\sum_{b\in B}b_b b_b^\top+I/10)^{-1}$, where
$b_{\{u,v\}}=e_u-e_v$. Exact rational inversion gives
\[
 H(\varnothing)=\frac{966250}{72181},\quad
 H(\{e\})=H(\{g\})=\frac{2486050}{197501},\quad
 H(\{e,g\})=\frac{204350}{17391}.
\]
These fractions are certified by rational Gaussian elimination in the
accompanying script, which also checks every entry of $AA^{-1}=I$.
Their mixed difference is exactly
\[
 H(\{e,g\})-H(\{e\})-H(\{g\})+H(\varnothing)
 =-\frac{9866618000}{257984349711}<0.
\]
For grounding the off-diagonal transition rates remain fixed and
only killing changes. Adding an edge instead changes two transition
rates, hence changes the base path and its return event. Additional
jumps cannot be represented by deleting trajectories with independent
vertex killing clocks. Algebraically the edge direction $e_u-e_v$
has both signs, unlike the grounding direction $e_v$; for $R=A^{-1}$,
direct differentiation gives
$\partial_x\partial_y\tr(A+xbb^\top+ycc^\top)^{-1}|_{0}
=2(b^\top Rc)(b^\top R^2c)$, which has no fixed sign.
Indeed $\partial_x\tr R=-b^\top R^2b$ and
$\partial_yR=-Rcc^\top R$ give the displayed derivative by the product
rule. This identifies exactly the unavailable step in the coupling
proof for edge addition.
\end{proof}
\noindent Numerical check: \texttt{scripts/verify\_t1\_edge\_counterexample.py}.

\section{T2: computational hardness}
\begin{lemma}[Cubic vertex-cover gap]\label{lem:cover}
Let $G$ be a connected unit-weight $3$-regular graph, let
$\varnothing\ne S\subsetneq V$, put $r=n-|S|$, and let $e_U$ count the
edges induced by $U=V\setminus S$. For every $s>0$ and $\sigma\geq0$,
\begin{equation}\label{eq:covergap}
 \tr (L_{-S}+\sigma I)^{-s}
 \geq r(3+\sigma)^{-s}+\delta_{s,\sigma}e_U,
 \qquad \delta_{s,\sigma}=\frac{s(s+1)}{(6+\sigma)^{s+2}}>0.
\end{equation}
Equality with the baseline $r(3+\sigma)^{-s}$ holds exactly when
$S$ is a vertex cover. Also $\lambda_{\min}(L_{-S}+\sigma I)\geq n^{-2}$.
\end{lemma}
\begin{proof}
Write $C=(3+\sigma)I-B_U$, where $B_U$ is the induced adjacency
matrix. The eigenvalues of $C$ lie in $(0,6+\sigma]$: positive
definiteness follows from the nonempty Dirichlet boundary, and
$\|B_U\|_2\leq3$. For $f(x)=x^{-s}$, throughout this interval
$f''(x)\geq\delta_{s,\sigma}$. Taylor's formula with integral remainder
therefore gives
\[
 f(x)\geq f(b)+f'(b)(x-b)+\tfrac12\delta_{s,\sigma}(x-b)^2,
 \qquad b=3+\sigma.
\]
Sum this inequality over the eigenvalues. The linear terms vanish
because $\tr C=rb$, and the sum of squared deviations is
$\tr B_U^2=2e_U$, proving \eqref{eq:covergap}. If $e_U=0$, $C=bI$
gives equality; if $e_U\geq1$ the excess is at least
$\delta_{s,\sigma}$. The equivalence with a vertex cover follows from
the definition.

For the lower eigenvalue bound, extend $x\in\R^U$ by zeros on $S$.
For each $u\in U$ choose a simple path from $u$ to $S$, of length at
most $n-1$. Cauchy--Schwarz bounds $|x_u|^2$ by $n$ times the sum of
edge differences squared along that path, hence by $n x^\top Lx$.
Summing over at most $n$ vertices yields
$\|x\|^2\leq n^2x^\top Lx\leq n^2x^\top(L+\sigma I)x$.
The Rayleigh quotient proves the asserted eigenvalue bound.
\end{proof}
\noindent Numerical check: \texttt{scripts/verify\_t2\_cover\_gap.py}.

\begin{lemma}[Uniform finite-penalty approximation]\label{lem:softlimit}
Under the hypotheses of Lemma~\ref{lem:cover}, for $0<s<1$ and $B>0$,
\begin{equation}\label{eq:softbound}
 0\leq \tr(L+\sigma I+B\diag(\ind_S))^{-s}
       -\tr(L_{-S}+\sigma I)^{-s}
 \leq n\left(\frac{1+9n^4}{B}\right)^s.
\end{equation}
\end{lemma}
\begin{proof}
Order the vertices as $U,S$ and write
\[
 L+\sigma I=\begin{pmatrix}C&D\\D^\top&E\end{pmatrix},\quad
 F=E+BI-D^\top C^{-1}D,\quad
 W=\begin{pmatrix}-C^{-1}D\\I\end{pmatrix}.
\]
Because $L+\sigma I\succeq0$ and $C\succ0$, its Schur complement is
positive semidefinite. For completeness, its quadratic form at $y$
is the quadratic form of $L+\sigma I$ at $(-C^{-1}Dy,y)$.
Consequently $F\succeq BI$. Block multiplication verifies
\[
 (L+\sigma I+B\diag(\ind_S))^{-1}
 =J+WF^{-1}W^\top,\qquad J=\begin{pmatrix}C^{-1}&0\\0&0\end{pmatrix}.
\]
Thus $0\preceq A^{-1}-J\preceq\epsilon I$, where
\[
 \epsilon=\frac{1+\|C^{-1}D\|_2^2}{B}
 \leq\frac{1+9n^4}{B}.
\]
Here $\|C^{-1}\|_2\leq n^2$ by Lemma~\ref{lem:cover}, and
$\|D\|_2\leq3$ since it is a submatrix of the cubic adjacency matrix.
The ordered eigenvalues satisfy
$\lambda_j(J)\leq\lambda_j(A^{-1})\leq\lambda_j(J)+\epsilon$
by the min--max principle. For $x\geq0$ and $0<s<1$,
$0\leq(x+\epsilon)^s-x^s\leq\epsilon^s$:
the difference is decreasing in $x>0$, since $s-1<0$, and its
maximum is attained at $x=0$. Apply this scalar inequality to the
matched eigenvalues and sum all $n$ terms, including the $|S|$ zeros
of $J$. Since $\tr J^s=\tr C^{-s}$, this proves \eqref{eq:softbound}.
\end{proof}
\noindent Numerical check: \texttt{scripts/verify\_t2\_soft\_limit.py}.

\begin{theorem}[Hardness for fixed fractional exponent]\label{thm:hardness}
For each fixed $s\in(0,1)$, finding an exact minimizing set of total
cardinality $k$ is NP-hard in each of the following regimes:
\begin{enumerate}
\item hard grounding on connected unit-weight cubic graphs, with any
fixed $\sigma\geq0$;
\item soft grounding on the same graphs, with any fixed $\sigma\geq0$
and an input integer penalty $\beta$ of polynomial magnitude;
\item soft grounding with $\sigma=0$ and any fixed positive rational
penalty $\beta_0$, on connected rationally weighted cubic graphs.
\end{enumerate}
For $\sigma=0$, the specified-singleton-seed formulation is also
NP-hard under polynomial-time Turing reductions.
\end{theorem}
\begin{proof}
Use vertex cover on connected $3$-regular graphs, the NP-complete
source problem used in Section 4 of Li et al.~\cite{li2019}.
The nontrivial instances have $1\leq k<n$. A cover of size at most
$k$ exists exactly when one of size $k$ exists, since a cover can
be enlarged. Lemma~\ref{lem:cover} proves that whenever a size-$k$
cover exists, every hard-objective minimizer is a cover. An optimal-set
oracle therefore decides vertex cover by returning a set whose edge
coverage can be checked in polynomial time. This proves item 1
without needing an exact representation of a real objective threshold.

For item 2, write $\delta=\delta_{s,\sigma}$. Fix integers
$p>1/s$ and $C>0$ such that $C^{-s}\leq\delta/8$, and choose
\[
 B=C(1+9n^4)n^p.
\]
These integers depend only on the fixed parameters $s,\sigma$.
The chosen $B$ has polynomial magnitude and $O(\log n)$ bits.
By Lemma~\ref{lem:softlimit} the error for every nonempty proper
set is at most
$nC^{-s}n^{-ps}\leq\delta/8$. Put $b=(n-k)(3+\sigma)^{-s}$.
A cover has soft value in $[b,b+\delta/8]$, whereas a noncover
has soft value at least $b+\delta$. Every soft minimizer is therefore
a cover whenever one exists. Again check the returned set's edge
coverage rather than comparing real-valued oracle outputs.

For item 3 use the same construction at $\sigma=0$, and replace
every unit edge weight by $c=\beta_0/B$. The new Laplacian is $cL$,
and exactly
\[
 cL+\beta_0\diag(\ind_S)=c(L+B\diag(\ind_S)),\qquad
 \tr(cL+\beta_0\diag(\ind_S))^{-s}
 =c^{-s}\tr(L+B\diag(\ind_S))^{-s}.
\]
The positive common factor preserves all minimizers. The rational
weights have $O(\log n)$ bits, and the support remains connected and
cubic. This is a polynomial reduction with fixed finite $\beta_0$.

Finally suppose the oracle requires a specified nonempty seed at
$\sigma=0$. Make one call for each singleton seed $\{v\}$ with the
same total cardinality $k$ (or $k-1$ additional choices).
If a size-$k$ cover exists, one such seed belongs to it, and the
separation arguments above ensure that the corresponding oracle
returns a cover. If no cover exists, no returned set passes the edge
coverage test. These $n$ calls are a polynomial-time Turing reduction.
\end{proof}
\noindent Numerical check: \texttt{scripts/verify\_t2\_hardness.py}.

\begin{corollary}[Strictly convex hard-grounding objectives]\label{cor:hardextended}
For a fixed $C^2$ function $f$ with $f''(x)\geq c_f>0$ on
$(0,6+\sigma]$, the cubic-graph objective obeys
\[
 \tr f(L_{-S}+\sigma I)\geq(n-|S|)f(3+\sigma)+c_f e_U.
\]
Equality with the baseline characterizes vertex covers. Consequently,
hard-grounding optimal-set selection is NP-hard for every fixed $s>0$
with $f(x)=x^{-s}$, and for every fixed $\tau>0$ with
$f(x)=e^{-\tau x}$, at any fixed $\sigma\geq0$.
The singleton-seed qualification in Theorem~\ref{thm:hardness} applies.
\end{corollary}
\begin{proof}
Repeat the Taylor inequality in Lemma~\ref{lem:cover} with $c_f$:
the linear terms sum to zero and the quadratic terms sum to
$c_f\tr B_U^2/2=c_f e_U$. A cover has $B_U=0$ and achieves the
baseline, whereas a noncover exceeds it by at least $c_f$.
Thus an optimal-set oracle decides the connected cubic vertex-cover
problem exactly as in Theorem~\ref{thm:hardness}. Powers use
$c_f=s(s+1)(6+\sigma)^{-s-2}$ for every $s>0$; heat uses
$c_f=\tau^2e^{-\tau(6+\sigma)}>0$. Checking returned sets avoids
real threshold encoding. The same singleton-seed calls prove the
seeded assertion. At $\tau=0$ the fixed-cardinality heat objective
is constant, so the heat statement requires $\tau>0$.
\end{proof}
\noindent Numerical check: \texttt{scripts/verify\_t2\_extended.py}.

\section{T3: scalar marginal oracles}
For this section let $A\succ0$ be the current matrix, let $b=e_v$ be
an available candidate in its coordinates, and put
\[
 A_u=A+ubb^\top,\quad C=A[-v,-v],\quad
 R_t=(A+tI)^{-1},\quad r_j(t)=b^\top R_t^j b\quad(j\geq1).
\]
For a finite penalty $\beta$, let $B=A_\beta$; for hard grounding let
$B=C$. Write $\Delta_f=\tr f(A)-\tr f(B)$ and
$\Delta_s=\Delta_{x^{-s}}$. These results apply to an already
soft-grounded matrix or to the ungrounded coordinates of an already
hard-grounded matrix. They hold for general positive definite $A$.

\begin{lemma}[Scalar resolvent marginal]\label{lem:resolvent}
For $t\geq0$, $r_2(t)=-r'_1(t)$ and
\begin{equation}\label{eq:qone}
 q_1(t):=\tr(A+tI)^{-1}-\tr(B+tI)^{-1}
 =\begin{cases}
 \displaystyle\frac{\beta r_2(t)}{1+\beta r_1(t)},&\beta<\infty,\\[6pt]
 \displaystyle\frac{r_2(t)}{r_1(t)},&\beta=\infty.
 \end{cases}
\end{equation}
Consequently, for $0<s<1$,
\begin{equation}\label{eq:fractionaloracle}
 \Delta_s=\frac{\sin(\pi s)}{\pi}
              \int_0^\infty t^{-s}q_1(t)dt,
 \qquad \Delta_1=q_1(0).
\end{equation}
\end{lemma}
\begin{proof}
Differentiating $(A+tI)R_t=I$ gives $R'_t=-R_t^2$, hence
$r'_1=-r_2$. For the finite update, direct multiplication verifies
the Sherman--Morrison identity
\[
 (A+tI+\beta bb^\top)^{-1}
 =R_t-\frac{\beta R_tbb^\top R_t}{1+\beta b^\top R_tb}.
\]
Taking traces gives the first formula. For hard grounding reorder
the candidate first and write $R_t$ in blocks on $\{v\},U$. The
block inverse identity gives
\[
 (C+tI)^{-1}=(R_t)_{UU}
       -\frac{(R_t)_{Uv}(R_t)_{vU}}{(R_t)_{vv}}.
\]
Its trace difference from $R_t$ is
$r_1+\|(R_t)_{Uv}\|^2/r_1=r_2/r_1$; the formula also holds when
$U$ is empty. The scalar identity
\[
 x^{-s}=\frac{\sin(\pi s)}{\pi}
           \int_0^\infty\frac{t^{-s}}{x+t}dt
\]
follows from $t=xu$ and the beta integral
$\int_0^\infty u^{-s}/(1+u)du=\Gamma(s)\Gamma(1-s)=\pi/\sin(\pi s)$.
Diagonalization gives the corresponding trace identity. The integrals
for each trace are finite, so subtracting them and using
\eqref{eq:qone} proves \eqref{eq:fractionaloracle}. At $s=1$ use
the trace difference at $t=0$ directly.
\end{proof}
\noindent Numerical check: \texttt{scripts/verify\_t3\_resolvent.py}.

\begin{proposition}[Positive scalar representation for all powers]\label{prop:allpowers}
Define, for integers $m\geq1$,
\[
 q_m(t)=\frac{(-1)^{m-1}}{(m-1)!}\frac{d^{m-1}}{dt^{m-1}}q_1(t).
\]
Then $q_m(t)=\tr(A+tI)^{-m}-\tr(B+tI)^{-m}\geq0$.
For any $s>0$ and any integer $m>s$,
\begin{equation}\label{eq:allpowers}
 \Delta_s=\frac{\Gamma(m)}{\Gamma(s)\Gamma(m-s)}
       \int_0^\infty t^{m-s-1}q_m(t)dt.
\end{equation}
For integer $s=m$ one can instead use $\Delta_m=q_m(0)$.
Each $q_m(t)$ is determined by the scalar moments
$r_1(t),\ldots,r_{m+1}(t)$ from the original matrix $A$.
\end{proposition}
\begin{proof}
Repeated differentiation of the shifted inverse gives
$d^{m-1}(A+tI)^{-1}/dt^{m-1}=(-1)^{m-1}(m-1)!(A+tI)^{-m}$,
which proves the trace identity. For a soft update the ordered
eigenvalues of $A_\beta$ are at least those of $A$, by the min--max
principle. For a hard update, writing $d=\dim A$, interlacing gives
$\lambda_i(A)\leq\lambda_i(C)$ for $1\leq i\leq d-1$.
Because $x\mapsto(x+t)^{-m}$ decreases, in the soft case every matched
eigenvalue contribution is nonnegative; in the hard case this also
leaves the positive contribution of $\lambda_d(A)$. This proves
$q_m\geq0$.

The substitution $t=xu$ in the beta integral gives
\[
 \int_0^\infty t^{m-s-1}(x+t)^{-m}dt
 =x^{-s}\frac{\Gamma(m-s)\Gamma(s)}{\Gamma(m)}.
\]
Apply it to the finite eigenvalue sums and subtract. The integral
of the difference is finite: near zero its integrand is bounded
by a constant times $t^{m-s-1}$, and near infinity by a constant
times $t^{-s-1}$. This proves \eqref{eq:allpowers} as a positive
representation, including $s\geq1$.

Here is an explicit scalar evaluation rule. Put
\[
 F(t)=\begin{cases}1+\beta r_1(t),&\beta<\infty,\\r_1(t),&\beta=\infty.
 \end{cases}
 \qquad q_1(t)=-\frac{d}{dt}\log F(t).
\]
The latter identity follows immediately from \eqref{eq:qone} and
$r'_1=-r_2$. In a neighborhood of $z=0$,
\[
 r_1(t+z)=\sum_{j\geq0}(-1)^j r_{j+1}(t)z^j,
 \quad F(t+z)=\sum_{j\geq0}f_jz^j,\quad
 \log F(t+z)=g_0+\sum_{j\geq1}g_jz^j.
\]
The resolvent series converges for $|z|<\lambda_{\min}(A)+t$;
$F(t)>0$ also ensures a local logarithm. For soft grounding
$f_0=1+\beta r_1$ and $f_j=\beta(-1)^jr_{j+1}$ for $j\geq1$;
for hard grounding $f_j=(-1)^jr_{j+1}$ for all $j\geq0$.
Equating coefficients in $F'=F(\log F)'$ gives, successively,
\begin{equation}\label{eq:logrecurrence}
 g_k=\frac{kf_k-\sum_{j=1}^{k-1}jg_jf_{k-j}}{kf_0},\qquad
 q_m(t)=(-1)^m m g_m.
\end{equation}
Thus moments through order $m+1$ suffice. For example,
\[
 q_2(t)=\begin{cases}
 \displaystyle\frac{2\beta r_3}{1+\beta r_1}
       -\frac{\beta^2r_2^2}{(1+\beta r_1)^2},&\beta<\infty,\\[6pt]
 \displaystyle\frac{2r_3}{r_1}-\frac{r_2^2}{r_1^2},&\beta=\infty.
 \end{cases}
\]
Its nonnegativity, and that of every higher-order expression, follows
from the trace/interlacing argument above.
\end{proof}
\noindent Numerical check: \texttt{scripts/verify\_t3\_powers.py}.

\begin{proposition}[Duhamel and positive penalty integrals]\label{prop:duhamel}
For $\tau>0$ and $0<\beta\leq\infty$,
\begin{equation}\label{eq:heatpenalty}
 \Delta_{e^{-\tau x}}=\tau\int_0^\beta
                b^\top e^{-\tau A_u}b\,du.
\end{equation}
For every $s>0$ there is also the representation
\begin{align}
 \Delta_s
 &=s\int_0^\beta b^\top A_u^{-(s+1)}b\,du\label{eq:powerpenalty}\\
 &=\frac1{\Gamma(s)}\int_0^\infty t^s
                \int_0^\beta b^\top e^{-tA_u}b\,du\,dt.\notag
\end{align}
All integrands are nonnegative. At $\tau=0$ the soft heat drop is zero
and the hard heat drop is one.
\end{proposition}
\begin{proof}
Differentiate the finite-dimensional evolution equation for
$e^{-\tau A_u}$ with respect to $u$. Variation of constants yields
the Duhamel identity
\[
 \frac{\partial}{\partial u}e^{-\tau A_u}
 =-\int_0^\tau e^{-(\tau-r)A_u}bb^\top e^{-rA_u}\,dr.
\]
Taking traces and using cyclicity combines the two exponentials,
so the derivative of the trace is $-\tau b^\top e^{-\tau A_u}b$.
Integrate from $0$ to $\beta<\infty$ to obtain \eqref{eq:heatpenalty}.
The Laplace representation of powers, followed by Tonelli applied
to this nonnegative heat drop, gives the double integral in
\eqref{eq:powerpenalty}. Integrating in $t$ first gives
\[
 \frac1{\Gamma(s)}\int_0^\infty t^s e^{-tA_u}dt
 =\frac{\Gamma(s+1)}{\Gamma(s)}A_u^{-(s+1)}
 =sA_u^{-(s+1)}.
\]
Thus \eqref{eq:powerpenalty} holds. The values are finite because
the two original traces are finite.

To justify the hard endpoint even for general positive definite $A$,
interlacing for a rank-one update bounds the first $d-1$ increasing
eigenvalues of $A_u$ between those of $A$; the last eigenvalue tends
to infinity, with $\lambda_d(A_u)/u\to1$, by the trace identity.
The determinant is affine in the candidate diagonal entry, giving
\[
 \frac{\det(A_u-zI)}{u}
 =\frac{\det(A-zI)}{u}+\det(C-zI)\longrightarrow\det(C-zI).
\]
Every convergent subsequence of the bounded eigenvalues must
therefore have exactly the eigenvalue multiset of $C$: the diverging
factor divided by $u$ tends to one, so the products of the remaining
factors converge to the characteristic polynomial of $C$.
It follows that the first $d-1$ eigenvalues converge to those of
$C$, in order. Both $x^{-s}$ and $e^{-\tau x}$ vanish at infinity,
so their trace limits are the corresponding trace of $C$.
Monotone convergence in the nonnegative penalty integrals proves
the formulas with $\beta=\infty$. At $\tau=0$ the traces count
dimensions, giving the two asserted endpoint values directly.
\end{proof}
\noindent Numerical check: \texttt{scripts/verify\_t3\_duhamel.py}.

\begin{theorem}[One Lanczos run for all shifts]\label{thm:lanczos}
Run $h$ steps of Lanczos on $A$ starting at the unit vector $b=e_v$.
In exact arithmetic write its orthonormal basis and tridiagonal as
\[
 AV_h=V_hT_h+\gamma_h v_{h+1}e_h^\top,\qquad V_he_1=b.
\]
At termination use $\gamma_h=0$. For every $t\geq0$, the shifted
Galerkin vector and its scalar quantities are
\begin{align}
 y_h(t)&=V_h(T_h+tI)^{-1}e_1,\notag\\
 \widehat r_1(t)&=e_1^\top(T_h+tI)^{-1}e_1,\label{eq:projectedmoments}\\
 \widehat r_2(t)&=e_1^\top(T_h+tI)^{-2}e_1=\|y_h(t)\|^2,
 \qquad -\widehat r'_1(t)=\widehat r_2(t),\notag\\
 \widehat r_j(t)&=e_1^\top(T_h+tI)^{-j}e_1\quad(j\geq1).\notag
\end{align}
If $\underline a\leq\lambda_{\min}(A)$ is positive, define
\[
 \rho_h(t)=\gamma_h|e_h^\top(T_h+tI)^{-1}e_1|,\qquad
 \epsilon_h(t)=\frac{\rho_h(t)}{\underline a+t}.
\]
Then
\begin{align}
 0\leq r_1-\widehat r_1&\leq\frac{\rho_h^2}{\underline a+t},\label{eq:r1bound}\\
 \bigl(\max\{0,\sqrt{\widehat r_2}-\epsilon_h\}\bigr)^2
 \leq r_2&\leq(\sqrt{\widehat r_2}+\epsilon_h)^2.\label{eq:r2bound}
\end{align}
At invariant-subspace termination all the projected moments are exact.
For any of the power or heat functions above, the projected marginal is
\begin{equation}\label{eq:reducedmarginal}
 \widehat\Delta_f=
 \begin{cases}
 \tr f(T_h)-\tr f(T_h+\beta e_1e_1^\top),&\beta<\infty,\\
 \tr f(T_h)-\tr f(T_h[2:h,2:h]),&\beta=\infty.
 \end{cases}
\end{equation}
It is computable entirely from the same $T_h$ and is exact at
invariant-subspace termination.
\end{theorem}
\begin{proof}
Binomial expansion, applied in both directions, gives
$\mathcal K_h(A+tI,b)=\mathcal K_h(A,b)$.
Projecting the shifted equation onto $V_h$ yields the stated Galerkin
vector. Orthonormality and $V_he_1=b$ give the first two scalar
identities; differentiating the small inverse gives their derivative
identity. Higher projected moments are scalar matrix functions of the
same tridiagonal; if $T_h=Q\diag(\theta_i)Q^\top$, then
\[
 \widehat r_j(t)=\sum_{i=1}^h Q_{1i}^2(\theta_i+t)^{-j}.
\]
Thus all shifts and all needed derivative moments share one run.

The Lanczos relation gives the residual
$z_h=b-(A+tI)y_h=-\gamma_hv_{h+1}e_h^\top(T_h+tI)^{-1}e_1$,
whose norm is $\rho_h$. If $x=(A+tI)^{-1}b$, then
$\|x-y_h\|\leq\rho_h/(\underline a+t)=\epsilon_h$.
Galerkin orthogonality gives $y_h^\top z_h=0$ and hence
\[
 r_1-\widehat r_1
 =(x-y_h)^\top(A+tI)(x-y_h)
 =z_h^\top(A+tI)^{-1}z_h.
\]
This proves \eqref{eq:r1bound}. The triangle and reverse-triangle
inequalities applied to $x=y_h+(x-y_h)$ give \eqref{eq:r2bound}.

For penalty parameter $u$ the projected matrix is exactly
$V_h^\top A_uV_h=T_h+ue_1e_1^\top$. Applying
Lemmas/Propositions~\ref{lem:resolvent}--\ref{prop:duhamel} to this
positive definite reduced matrix gives its scalar integral formulas;
they equal \eqref{eq:reducedmarginal}. In particular, the hard
projection deletes the first reduced coordinate, not a coordinate of
an independently generated basis.

When $\gamma_h=0$, the Krylov subspace is invariant. Since $A$ is
symmetric, its orthogonal complement is also invariant. In a completed
orthonormal basis,
\[
 A=T_h\oplus A_\perp,\qquad
 A_u=(T_h+ue_1e_1^\top)\oplus A_\perp.
\]
The vector $b$ lies entirely in the first block, so all its scalar
moments equal the reduced moments. For hard grounding, restriction
to $b^\perp$ is orthogonally similar to
$T_h[2:h,2:h]\oplus A_\perp$. In either update the complement's
trace cancels. This proves exactness of \eqref{eq:reducedmarginal},
including early breakdown and $h=1$.

In exact arithmetic the three-term run uses $h$ large sparse
matrix--vector products and $O(dh)$ vector work. For $q$ shifts,
tridiagonal solves and vector norms compute
$\widehat r_1,\widehat r_2,\rho_h$ in $O(qh)$ additional work.
Moments through $m+1$ can be obtained by repeated tridiagonal solves
in $O(qmh)$, followed by $O(qm^2)$ work in
\eqref{eq:logrecurrence}. Power/heat evaluations in
\eqref{eq:reducedmarginal} or the penalty integrals involve only small
matrices and require no further large-matrix products.
\end{proof}
\noindent Numerical check: \texttt{scripts/verify\_t3\_lanczos.py}.

\section{T4: certified marginal bounds and fixed positive quadrature}
Throughout this section $\|b\|=1$, $0<a<\lambda_{\min}(A)$,
$\lambda_{\max}(A)\leq B$, $f$ is CM, $f_\infty=f(\infty)$,
and $g=-f'$. The function $g$ is CM: its Laplace measure is
$t\mu(dt)$, and it is finite at every positive argument.

\begin{theorem}[Shared-penalty Gauss--Radau certificate]\label{thm:radau}
Let $T_h,\gamma_h$ be obtained by one Lanczos run from $b$.
Set $J_h(u)=T_h+ue_1e_1^\top$. If $\gamma_h>0$, define
\begin{equation}\label{eq:radau}
 H_{h+1}(u)=
 \begin{pmatrix}
 J_h(u)&\gamma_h e_h\\
 \gamma_h e_h^\top&a+\gamma_h^2 e_h^\top(J_h(u)-aI)^{-1}e_h
 \end{pmatrix}.
\end{equation}
At invariant termination set $H=J$ and use the exact value.
For every $u\geq0$,
\[
 G_h(u):=e_1^\top g(J_h(u))e_1
 \leq F(u):=b^\top g(A_u)b
 \leq U_h(u):=e_1^\top g(H_{h+1}(u))e_1.
\]
For finite $\beta$, integrating from $0$ to $\beta$ gives a marginal
interval. For hard grounding the interval is
\[
 f_\infty+\int_0^\infty G_h(u)du
 \ \leq\Delta_f\leq\
 f_\infty+\int_0^\infty U_h(u)du.
\]
All endpoints are finite. The lower endpoint equals the reduced trace
drop in \eqref{eq:reducedmarginal}. Put
$d=b^\top Ab$ and $v=\|Ab-db\|^2$. For $K+d\geq2a$, a common bound
on the exact and Radau upper tails is
\begin{equation}\label{eq:penaltytail}
 \int_K^\infty F(u)du,\quad\int_K^\infty U_h(u)du
 \ \leq\frac{4v g(a)}{d+K}
                 +2\bigl(f((d+K)/2)-f_\infty\bigr).
\end{equation}
\end{theorem}
\begin{proof}
The spaces $\mathcal K_h(A_u,b)$ and $\mathcal K_h(A,b)$ agree:
successively expanding $(A+ubb^\top)^j b$ gives $A^jb$ plus a
combination of earlier Krylov vectors, and the reverse inclusion
follows by replacing $u$ by $-u$. More precisely the same Lanczos
recurrence works with only the first diagonal changed by $u$, since
$b^\top v_j=0$ for $j>1$. Its residual coefficient is still $\gamma_h$.

Let $\nu_u$ be the spectral probability measure of $(A_u,b)$.
The $h$-node rule defined by $J_h(u)$ has positive weights and is exact
through degree $2h-1$. One direct verification extends the Lanczos
tridiagonal to its invariant dimension: a walk from coordinate 1
that visits coordinate $h+1$ and returns needs at least $2h$ steps.
The augmented matrix $H$ shares the first $h$ diagonals and the first
$h$ off-diagonals with this full tridiagonal. Its changed last diagonal
can first affect a return walk at length $2h+1$, so its spectral rule
is exact through degree $2h$.
The Schur complement of $H-aI$ is zero, so $a$ is an eigenvalue and
$H\succeq aI$. Its other nodes are the zeros of the degree-$h$
orthogonal polynomial for $(x-a)\nu_u$: exactness makes
$\det(xI-H)/(x-a)$ orthogonal to degrees $0,\ldots,h-1$ for that
measure. Positivity of the measure places these roots inside the
convex hull of its support. This last assertion follows from the
usual sign-change argument: fewer than $h$ interior sign changes
would give a lower-degree polynomial whose product with it has a
fixed nonzero sign, contradicting orthogonality.
The spectral weights are squares of first eigenvector coordinates,
are positive before breakdown, and sum to one.

Hermite interpolation at the Gauss nodes (double interpolation)
has remainder
$g^{(2h)}(\xi)\prod_j(x-\theta_j)^2/(2h)!\geq0$.
For Radau, interpolate once at $a$ and twice at the other $h$ nodes.
The remainder is
$g^{(2h+1)}(\xi)(x-a)\prod_j(x-\theta_j)^2/(2h+1)!\leq0$.
These remainder formulas follow by repeated Rolle's theorem, and
hold also at the interpolation nodes by continuity. Integrating
and using polynomial exactness proves both quadrature inequalities.

The Duhamel proof extends to CM $f$ by Tonelli applied to its
nonnegative Laplace mixture. It gives $\Delta_f=\int_0^\beta F(u)du$
for a finite penalty, and $\Delta_f=f_\infty+\int_0^\infty F(u)du$
for hard grounding: the diverging eigenvalue contributes
$f_\infty$ in the soft limit. Apply the same identity to $T_h$ for
the lower endpoint.

For finiteness and an explicit tail, both the exact and Radau spectral
measures of the first vector have mean $d+u$, variance $v$, and
support in $[a,\infty)$. The mass below $(d+u)/2$ is at most
$4v/(d+u)^2$, by integrating $(x-d-u)^2$ on that event.
Since $g$ is decreasing, either bilinear form is at most
$4v g(a)/(d+u)^2+g((d+u)/2)$. Integration proves
\eqref{eq:penaltytail}.

There is also a directly evaluable finite-sum certificate for the
penalty integration. The exact $F$ is decreasing: write $g$ as a
Laplace mixture and apply Duhamel to its bilinear form; the negative
derivative is a positive mixture of
$\int_0^t (b^\top e^{-(t-r)A_u}b)(b^\top e^{-rA_u}b)dr$.
Differentiation on bounded $u$-intervals is justified by an integrable
exponential majorant at a positive spectral lower bound.
For any partition of $[0,K]$, lower sums using $G_h$ at right endpoints
and upper sums using $U_h$ at left endpoints bracket the exact integral.
For hard grounding add $f_\infty$ and the tail bound to the upper sum.
The lower reduced trace drop is an alternative, tighter certified
lower endpoint requiring no penalty integration.
\end{proof}
\noindent Numerical check: \texttt{scripts/verify\_t4\_radau.py}.

\begin{theorem}[Width and the dependence on the CM function]\label{thm:width}
Define
\[
 \rho(a,C)=\frac{\sqrt{2C-a}+\sqrt a}{\sqrt{2C-a}-\sqrt a},\qquad
 E_D(a,C)=\frac{2g(a/2)}{\rho(a,C)-1}\rho(a,C)^{-D}.
\]
The reduced marginal $L_h=\widehat\Delta_f$ gives the following
single-run intervals, with no numerical integration:
\begin{align}
 L_h\leq\Delta_f&\leq L_h+2\beta E_{2h-1}(a,B+\beta)
                 &&\text{(soft)},\label{eq:widthsoft}\\
 L_h\leq\Delta_f&\leq L_h+(B-a)E_{2h-2}(a,B)
                 &&\text{(hard)}.\label{eq:widthhard}
\end{align}
For relative width $\eta$, it suffices to make the added term at most
$\eta\beta g(B+\beta)$ in the soft case or $\eta f(B)$ in the hard
case. For each fixed power $s>0$ this gives
$h=O(\sqrt\kappa((s+1)\log\kappa+\log(1/\eta)))$, with the
corresponding interval condition number. For heat the bound retains
the parameter $\tau B$ (or $\tau(B+\beta)$).
There is no relative convergence rate uniform over the entire heat
family that depends only on $\kappa,h$.
\end{theorem}
\begin{proof}
For a rank-one update define the spectral count difference
$\xi(x)=N_A(x)-N_B(x)$. Interlacing gives $0\leq\xi\leq1$.
In the soft case it vanishes outside $[a,B+\beta]$,
$\int\xi=\beta$, and $\Delta_f=\int g(x)\xi(x)dx$.
In the hard case it equals one above $B$, and integration by parts
gives $\Delta_f=f(B)+\int_a^B g(x)\xi(x)dx$.
These identities follow first for a finite step function $\xi$ by
telescoping over its jumps, then for the constant tail using
$\int_B^\infty g=f(B)-f_\infty$. The same statements apply to
the reduced matrices, with spectral count $\widehat\xi$.

The trace drops of polynomials of degree at most $2h$ (soft) or
$2h-1$ (hard) agree exactly. To verify this, for large $z$ put
$m_A(z)=b^\top(zI-A)^{-1}b=z^{-1}\sum_{j\geq0}(b^\top A^jb)z^{-j}$.
The first $2h-1$ moments agree with those of $T_h$.
The hard determinant identity is
$m_A(z)=\det(zI-C)/\det(zI-A)$; the coefficient of $z^{-p}$
in $\log(zm_A(z))$ is $(\tr A^p-\tr C^p)/p$ and depends only
on moments through $p$. The soft identity is
$\det(zI-A_\beta)/\det(zI-A)=1-\beta m_A(z)$;
its logarithmic coefficient is the soft trace drop divided by $p$
and depends only on moments through $p-1$. The asserted degrees
follow. All series converge for sufficiently large $z$.
Therefore $\int p(x)(\xi-\widehat\xi)dx=0$ for every polynomial
of degree $2h-1$ (soft) or $2h-2$ (hard).

On the Bernstein ellipse for $[a,C]$ with parameter $\rho(a,C)$,
the leftmost point is $a/2$. The Laplace representation implies
$|g(z)|\leq g(a/2)$ there. Expanding in Chebyshev polynomials,
the contour coefficient bound is $2g(a/2)\rho^{-j}$; it follows
by applying Cauchy's coefficient formula after the map
$z=(a+C)/2+(C-a)(w+w^{-1})/4$.
Summing the tail yields a degree-$D$ polynomial $p$ with
$\|g-p\|_{[a,C]}\leq E_D(a,C)$.
For soft grounding, $\int|\xi-\widehat\xi|\leq2\beta$.
For hard grounding, $|\xi-\widehat\xi|\leq1$ on an interval
of length $B-a$. Subtract the matching polynomial in the spectral
count formulas to obtain the absolute errors in
\eqref{eq:widthsoft}--\eqref{eq:widthhard}; their one-sided sign
comes from Theorem~\ref{thm:radau}.

The spectral formulas also give
$\Delta_f\geq\beta g(B+\beta)$ (soft) and
$\Delta_f\geq f(B)$ (hard). With $K_s=2\beta$, $C=B+\beta$,
$d_*=\beta g(C)$, or $K_s=B-a$, $C=B$, $d_*=f(B)$,
an explicit sufficient degree is
\[
 D\geq\max\left\{0,
 \frac{\log(2K_sg(a/2)/(\eta d_*(\rho-1)))}{\log\rho}\right\}.
\]
Use $D=2h-1$ or $2h-2$, respectively. Constant functions are handled
exactly. For $g(x)=sx^{-s-1}$ the logarithm of the prefactor is
$O((s+1)\log\kappa+\log(1/\eta))$; also
$\log\rho=\Theta(\kappa^{-1/2})$. For heat,
$g(x)=\tau e^{-\tau x}$, giving the stated $\tau$ dependence.

For the last assertion take the $n\times n$ tridiagonal with diagonal
2 and off-diagonal $-1/2$, starting at its endpoint. Its eigenvalues
are $2-\cos(j\pi/(n+1))$ (substitute the sine eigenvectors), so its
condition number is less than 3. The $h<n$ step projection is the
same tridiagonal of size $h$. The exact hard heat drop is asymptotic
to $e^{-\tau\lambda_{\min}(A_n)}$, since the deleted principal
matrix has a strictly larger smallest eigenvalue. The projected drop
is asymptotic to $e^{-\tau\lambda_{\min}(A_h)}$.
Their ratio tends to zero as $\tau\to\infty$. Arbitrarily large $n$
and bounded condition number therefore preclude a function-independent
relative iteration bound.
\end{proof}
\noindent Numerical check: \texttt{scripts/verify\_t4\_convergence.py}.

\begin{theorem}[Fixed positive surrogate with logarithmic shift count]\label{thm:fixedquad}
Suppose every nonempty admissible grounded matrix has spectrum in
$[a,B]$, $a>0$, and put $\kappa=B/a$. For each $s>0$ and
$0<\eta<1$, there is a fixed positive rule
\[
 \widetilde f_s(x)=\sum_{j=1}^q w_j(x+t_j)^{-m},\quad
 t_j,w_j>0,\quad m=\lfloor s\rfloor+1,
\]
such that its trace objective is nonincreasing and supermodular at
every rule order and, for every soft or hard marginal,
\begin{equation}\label{eq:fixedrelative}
 (1-\eta)\Delta_{x^{-s}}\leq\Delta_{\widetilde f_s}\leq\Delta_{x^{-s}}.
\end{equation}
One explicit choice is
\[
 p=m+3+\left\lceil\log_4(2/\eta)\right\rceil,\qquad
 q\leq(\lceil\log_2\kappa\rceil+2)p.
\]
Thus $q=O_s((1+\log\kappa)\log(1/\eta))$ for a fixed exponent.
For $0<s<1$ the rule uses only first inverse powers ($m=1$);
for $s=1$ one can use the exact single shift $t=0$ instead.
\end{theorem}
\begin{proof}
Use the positive integral in Proposition~\ref{prop:allpowers}.
Partition its shift domain into $[0,a]$, geometric intervals
$[a,2a],[2a,4a],\ldots$ truncated at $B$, and $[B,\infty)$.
On $[0,a]$ use $t=az$ and $p$-node Gauss--Jacobi for the weight
$z^{m-s-1}$ on $[0,1]$. On every interior interval use $p$-node
Gauss--Legendre, including the positive factor $t^{m-s-1}$ in
the sampled integrand. On the tail use $z=B/t$ and Gauss--Jacobi
with weight $z^{s-1}$. The tail integrand becomes
$B^{-s}(1+xz/B)^{-m}$, a positive multiple of
$(x+B/z)^{-m}$. All nodes and coefficients are positive and
independent of $S$ or the candidate. Write the resulting scalar
approximation as $F_p(x)$.

Here is a uniform error bound for both $f_s$ and $g_s=-f_s'$.
Put $\alpha=s+1-m\in[0,1)$ and $D=2p-1$.
For an interior interval with center $c$ write $y=(t-c)/c$;
its doubling ratio ensures $|y|\leq1/3$. After extracting a
positive factor, the integrands for the function and its negative
derivative have the form
$(1+y)^{-\alpha}(1+\theta y)^{-r}$ with $0\leq\theta\leq1$
and $r=m,m+1$. On $|y|=2/3$ their modulus is at most
$3^{\alpha+r}$, while their minimum on the real interval is at
least $(4/3)^{-\alpha-r}$.
Cauchy's coefficient bound shows that the Taylor tail beyond degree
$D$ is at most $2\,3^{\alpha+r}2^{-(D+1)}$.
The quadrature integrates the Taylor polynomial exactly; positivity
and exactness on constants bound its error plus the integral's error
by twice this remainder times the interval mass. Dividing by the
true integral gives relative error at most
$4\,4^{\alpha+r}4^{-p}\leq4^{m+3-p}$.

For the lower endpoint the weight absorbs the singularity, leaving
$(x+az)^{-r}$. Expanding about $z=1/2$ gives a normalized variable
of radius $a/(2x+a)\leq1/3$ and the same, or smaller, bound.
For the upper endpoint the derivative integrand, under its fixed
weight, is proportional to $z(1+xz/B)^{-m-1}$.
Approximate its denominator factor to degree $D-1$ and multiply
by $z$. The error argument now uses exactness on $z$, with radius
$x/(2B+x)\leq1/3$. The bound is
$8\,4^{m+1}4^{-p}\leq4^{m+3-p}$.
The function tail needs degree $D$ and has a smaller bound.
These estimates hold for every $x\in[a,B]$. Summing positive panel
integrals gives, with $\epsilon=4^{m+3-p}\leq\eta/2$,
\[
 |F_p/f_s-1|\leq\epsilon,\qquad
 |(-F_p')/g_s-1|\leq\epsilon.
\]
Define $\widetilde f_s=F_p/(1+\eta/2)$. Both it and its negative
derivative lie between $1-\eta$ and 1 times their true counterparts.

For soft grounding integrate the derivative bounds against the
nonnegative spectral count $\xi$. For hard grounding use
$\Delta_f=f(B)+\int_a^B(-f')\xi$, so the endpoint value is controlled
as well. This proves \eqref{eq:fixedrelative} without taking a
relative error of two nearly equal traces. In particular it treats
the hard ratio $r_2/r_1$ rather than a log-resolvent substitute.
Every shifted inverse power is CM, so Theorem~\ref{thm:cm} proves
supermodularity of every finite positive rule, including orders
below the accuracy threshold. There are at most
$\lceil\log_2\kappa\rceil+2$ panels, proving the count.
\end{proof}
\noindent Numerical check: \texttt{scripts/verify\_t4\_quadrature.py}.

\paragraph{Spectral bounds used by the rule.}
One may take $B=2d_{\max}+\sigma+\beta$ for soft grounding and
$B=2d_{\max}+\sigma$ for hard grounding. If $\sigma>0$, take
$a=\sigma$. For a connected weighted graph at $\sigma=0$, put
$w_*=\min_e w_e$. A nonempty hard seed permits $a=w_*/n^2$.
For a soft seed a valid bound is
$a=(2n^2/w_*+2n/\beta)^{-1}$: along a path from $u$ to a seeded
vertex $v$, $x_u^2\leq2x_v^2+2n\sum_{e\text{ on path}}(\nabla_ex)^2$;
sum over $u$ and compare with the grounded quadratic form.
The hard bound follows from the same path inequality with $x_v=0$,
and the upper bounds follow from the edge-energy inequality
$(x_i-x_j)^2\leq2x_i^2+2x_j^2$.

\section{T5: incremental scalar caches}
\begin{proposition}[Exact two-moment updates]\label{prop:updates}
At one shift let $R=(A+tI)^{-1}$, $x=Re_v$, $y=R^2e_v$, and
$z=r_{2,v}=x^\top x$. Put $c=(\beta^{-1}+r_{1,v})^{-1}$, using
$\beta^{-1}=0$ for hard grounding. Every remaining candidate $u$ obeys
\begin{align}
 r_{1,u}^+&=r_{1,u}-c x_u^2,\label{eq:updateone}\\
 r_{2,u}^+&=r_{2,u}-2c x_u y_u+c^2 z x_u^2.\label{eq:updatetwo}
\end{align}
Once $x,y$ are available, updating all candidates at $q$ shifts costs
$O(nq)$ arithmetic and $O(nq)$ cache storage. Selected-column generation
requires two solves per shift, or one retained Krylov basis from $e_v$
with both small functions $(T+tI)^{-1}e_1,(T+tI)^{-2}e_1$.
\end{proposition}
\begin{proof}
Sherman--Morrison gives $R^+=R-cxx^\top$ for soft grounding.
For hard grounding the same matrix has zero row and column $v$;
its complementary block is $(A[-v,-v]+tI)^{-1}$ by the block inverse
identity in Lemma~\ref{lem:resolvent}. Thus diagonal entries, and
squared-column norms at remaining coordinates, can be computed before
or after deleting the zero row and column. Squaring gives
\[
 (R-cxx^\top)^2=R^2-c(yx^\top+xy^\top)+c^2(x^\top x)xx^\top.
\]
Taking diagonals proves both updates. Two shifted solves produce
$x$ and $y$; a retained Lanczos basis gives their projected versions
from the two indicated small vectors, exactly at termination.
Reconstructing both vectors at all shifts costs $O(nqh)$ after an
$h$-step basis is stored, in addition to the large matrix--vector
products. Only the cache update itself costs $O(nq)$.
\end{proof}
\noindent Numerical check: \texttt{scripts/verify\_t5\_updates.py}.

\begin{theorem}[Cache-error propagation]\label{thm:cache}
Suppose the current scalar errors are bounded by $e_{1,u},e_{2,u}$.
For one update suppose the errors in $c,x_u,y_u,z$ are bounded by
$e_c,e_x,e_y,e_z$, and both exact and approximate quantities have
absolute bounds $C,X,Y,Z$, respectively. Then valid updated errors are
\begin{align}
 e_{1,u}^+&=e_{1,u}+e_cX^2+2CXe_x,\label{eq:cacheone}\\
 e_{2,u}^+&=e_{2,u}+2e_cXY+2C(e_xY+Xe_y)
       +2Ce_cZX^2+C^2e_zX^2+2C^2ZXe_x.\label{eq:cachetwo}
\end{align}
If $\ell\leq r_{1,v},\widehat r_{1,v}$ is positive, one may use
$e_c=e_{1,v}/(\beta^{-1}+\ell)^2$ and $C=(\beta^{-1}+\ell)^{-1}$.
If $a_t\leq\lambda_{\min}(A+tI)$ and approximate selected columns
have residuals $p=e_v-(A+tI)\widehat x$,
$q=\widehat x-(A+tI)\widehat y$, valid column-error bounds are
\[
 e_x=\|p\|/a_t,\qquad e_y=(\|q\|+e_x)/a_t.
\]
These deterministic bounds remain valid under adaptive choices of
the selected vertex.
\end{theorem}
\begin{proof}
For a product use telescoping with the stated envelopes. In particular,
$|\widehat x^2-x^2|\leq2Xe_x$ and
$|\widehat c^2-c^2|\leq2Ce_c$.
The error in $cx^2$ is at most $e_cX^2+2CXe_x$;
the error in $2cxy$ is at most $2e_cXY+2C(e_xY+Xe_y)$;
the error in $c^2zx^2$ is at most
$2Ce_cZX^2+C^2e_zX^2+2C^2ZXe_x$.
Adding the old cache error proves the formulas.
The reciprocal difference identity gives the bound on $e_c$.
The residual equations give
$\widehat x-x=-(A+tI)^{-1}p$ and
$\widehat y-y=(A+tI)^{-1}(\widehat x-x-q)$, proving the solve bounds.

Useful deterministic envelopes are $\ell=1/B_t$, $X=1/a_t$,
$Y=Z=1/a_t^2$, with estimates clipped to these known ranges.
Projection onto a known interval containing the truth cannot increase
absolute error. SDDM inverses and their squares are entrywise
nonnegative, so columns may be clipped to $[0,X]$ and $[0,Y]$.
To see inverse positivity directly, choose $d\geq\max_i A_{ii}+t$;
$A+tI=dI-N$ with $N\geq0$ and $\rho(N)<d$, and expand the inverse
as a convergent nonnegative Neumann series.

For an explicit multiround bound set $E_i=\max_u e_{i,u}$ and
$D=\beta^{-1}+\ell$. With uniform column bounds one step satisfies
\[
 \begin{pmatrix}E_1^+\\E_2^+\end{pmatrix}
 \leq
 \begin{pmatrix}
 1+X^2/D^2&0\\
 (2XY+2CZX^2)/D^2&1+C^2X^2
 \end{pmatrix}
 \begin{pmatrix}E_1\\E_2\end{pmatrix}
 +\begin{pmatrix}
 2CXe_x\\(2CY+2C^2ZX)e_x+2CXe_y
 \end{pmatrix}.
\]
The coefficient matrices are nonnegative. Their ordered products and
the corresponding propagated additive terms give a bound for any
number of rounds by induction. This induction is pathwise and
therefore applies even when each selected vertex is chosen from
the current approximate cache.
\end{proof}
\noindent Numerical check: \texttt{scripts/verify\_t5\_cache\_errors.py}.

\section{T6: all-node initialization}
Let $A$ be the fixed initial positive definite SDDM matrix, with
$aI\preceq A\preceq BI$, $N$ active vertices, and $m_E$ active
off-diagonal edges. All shifts in the initialization are fixed before
the random sketch is drawn. In this section $d$ denotes sketch
dimension and $q$ denotes the number of requested shifts.

\begin{lemma}[Uniform JL on the entire resolvent curve]\label{lem:continuum}
Let $Q\in\R^{d\times N}$ have independent $N(0,1/d)$ entries,
$0<\epsilon,\delta<1$, and $\kappa=B/a$. Put
\[
 M_*=\sqrt{2N/\delta},\quad
 K=2+\left\lceil\frac{4M_*\log\kappa}{\epsilon}\right\rceil,
 \quad d\geq\frac{32}{\epsilon^2}\log\frac{4NK}{\delta}.
\]
With probability at least $1-\delta$, simultaneously for every
vertex $u$ and every $t\in[0,\infty)$,
\begin{equation}\label{eq:continuumjl}
 (1-\epsilon)r_{2,u}(t)
 \leq\|Q(A+tI)^{-1}e_u\|^2
 \leq(1+\epsilon)r_{2,u}(t).
\end{equation}
Consequently integrating the middle expression over $[t,\infty)$
approximates $r_{1,u}(t)$ with the same factor, simultaneously for
every $u,t$.
\end{lemma}
\begin{proof}
For a fixed unit vector $w$, $d\|Qw\|^2$ is chi-squared with $d$
degrees of freedom. Its moment generating function
$(1-2\theta)^{-d/2}$, Chernoff's method, and
$x-\log(1+x)\geq x^2/4$ for $0<x<1$ give
\begin{equation}\label{eq:chisquare}
 \Prob\{|\|Qw\|^2-1|>x\}\leq2e^{-dx^2/8}.
\end{equation}
For the lower tail use $-x-\log(1-x)\geq x^2/2$; these scalar
inequalities follow by differentiating both sides from zero.

Normalize $w_u(t)=(A+tI)^{-1}e_u/\|(A+tI)^{-1}e_u\|$.
Writing $R=(A+tI)^{-1}$, differentiation gives
$w'_u=-(R-(w_u^\top Rw_u)I)w_u$. Its squared norm is the variance
of the eigenvalues of $R$ under the spectral weights of $w_u$.
For a variable in $[l,r]$, its variance is at most $(r-l)^2/4$:
average $(X-l)(X-r)\leq0$ and then maximize
$(r-\E X)(\E X-l)$. Hence
\[
 \|w'_u(t)\|\leq\frac{B-a}{2(a+t)(B+t)}.
\]
The total length is at most $\tfrac12\log\kappa$.
The parameter
$\phi(t)=\tfrac12\log(B(a+t)/(a(B+t)))$ bounds the arc distance
between any two parameters by their $\phi$ difference, and
$w_u(\infty)=e_u$. Choose a deterministic net with mesh
$r_* =\epsilon/(8M_*)$ in $\phi$, including the limit at infinity.
At most $K$ net points per vertex suffice; when $\kappa=1$ each
curve is constant.

Apply \eqref{eq:chisquare} with $x=\epsilon/2$ and a union bound
over these $NK$ unit vectors. The failure probability is at most
$\delta/2$. Also $\E\|Q\|_F^2=N$, so Markov's inequality gives
$\|Q\|_2\leq\|Q\|_F\leq M_*$ except with probability $\delta/2$.
On their intersection, for any point and its nearest net point,
the projected norms differ by at most $M_*r_* =\epsilon/8$.
Squaring
$\sqrt{1-\epsilon/2}-\epsilon/8\leq\|Qw_u(t)\|
\leq\sqrt{1+\epsilon/2}+\epsilon/8$
gives \eqref{eq:continuumjl} for $0<\epsilon<1$.
This proves the uncountable event by a finite deterministic net.

Finally spectral decomposition gives
$\int_t^\infty (A+zI)^{-2}dz=(A+tI)^{-1}$.
All integrands in \eqref{eq:continuumjl} are nonnegative, so its
simultaneous event may be integrated without a second union bound.
For a finite implementation, the $s=1,m=2$ rule of
Theorem~\ref{thm:fixedquad}, applied on $[a+t,B+t]$, replaces this
integral by positive evaluations at extra shifts. Its relative
factor $[1-\eta,1]$ combines with JL to give
$[(1-\epsilon)(1-\eta),1+\epsilon]$. Thus no finite-grid assumption
about the unevaluated continuum is needed.
\end{proof}
\noindent Numerical check: \texttt{scripts/verify\_t6\_continuum\_jl.py}.

\begin{theorem}[Two RHS families: certified all-node initialization]\label{thm:initialization}
Factor $A=C^\top C$, where $C$ has $O(m_E+N)$ nonzeros.
Let $Q\in\R^{d\times N}$ and $Z\in\R^{d\times r}$ be independent
Gaussian matrices with entries $N(0,1/d)$, where $r$ is the row count
of $C$, and put $P=ZC$. For every requested shift define
\[
 Y_j=Q(A+t_jI)^{-1},\qquad X_j=P(A+t_jI)^{-1}.
\]
Thus only the $2d$ rows of $P,Q$ are right-hand sides, independently
of the number of vertices and shifts. Define the approximate outputs
\begin{equation}\label{eq:allnodediags}
 \widehat r_{2,u}(t_j)=\|\widehat Y_j[:,u]\|^2,\qquad
 \widehat r_{1,u}(t_j)
   =\|\widehat X_j[:,u]+\sqrt{t_j}\widehat Y_j[:,u]\|^2.
\end{equation}
For $0<\epsilon,\delta<1$, take
\[
 d\geq128\epsilon^{-2}\log(12Nq/\delta),\qquad
 a_j=a+t_j,\quad B_j=B+t_j.
\]
Require the two Frobenius residual tests
\begin{align}
 \|Q-\widehat Y_j(A+t_jI)\|_F
       &\leq\frac{\epsilon a_j}{16B_j},\label{eq:resY}\\
 \|P-\widehat X_j(A+t_jI)\|_F
       &\leq\frac{\epsilon a_j}{16\sqrt{B_j}}.\label{eq:resX}
\end{align}
If the solves satisfy these tests with total failure probability at
most $\delta/3$, then with probability at least $1-\delta$ all
$2Nq$ estimates have relative error at most $\epsilon$.
Using nearly-linear SDDM solvers, the initialization takes
$\widetilde O((m_E+N)q\epsilon^{-2})$ time; the suppressed factors
include logarithms of $N,\kappa,q,1/\epsilon,1/\delta$.
The outputs and streaming accumulation need $O(Nq)$ storage in
addition to graph and solver workspace.
\end{theorem}
\begin{proof}
For an SDDM matrix write the off-diagonals as $-w_{uv}$ and its row
sums as nonnegative $k_u$. Give $C$ one incidence row
$\sqrt{w_{uv}}(e_u-e_v)^\top$ per edge and a row $\sqrt{k_u}e_u^\top$
per nonzero row sum. Direct multiplication gives $C^\top C=A$
and the nonzero bound. Grounded principal submatrices retain this
form, including boundary-edge contributions in the row sums.

For a fixed $j,u$, $Y_j[:,u]$ has independent centered Gaussian
coordinates with variance $r_{2,u}(t_j)/d$.
The rows of $P+\sqrt{t_j}Q$ have covariance
$(A+t_jI)/d$, since $Z$ and $Q$ are independent. Therefore
$X_j[:,u]+\sqrt{t_j}Y_j[:,u]$ has coordinate variance
$r_{1,u}(t_j)/d$. Independence between different shifts or between
these two estimators is unnecessary. Equation~\eqref{eq:chisquare}
with $x=\epsilon/4$ and a union bound gives simultaneous squared-norm
error $\epsilon/4$ with failure at most
$4Nq\exp(-d\epsilon^2/128)\leq\delta/3$.

The solution error's Frobenius norm is at most its residual norm
divided by $a_j$. Thus the error in each $Y$ column is at most
$\epsilon/(16B_j)$, and the error in each combined $X+\sqrt{t_j}Y$
column is at most
\[
 \frac{\epsilon}{16\sqrt{B_j}}
       +\frac{\epsilon\sqrt{t_j}}{16B_j}
 \leq\frac{\epsilon}{8\sqrt{B_j}}.
\]
The exact target squared norms satisfy
$r_{2,u}\geq B_j^{-2}$ and $r_{1,u}\geq B_j^{-1}$.
The triangle inequality therefore bounds normalized approximate
norms by $\sqrt{1\pm\epsilon/4}\pm\epsilon/8$ (with a still
smaller perturbation for $r_2$). Squaring gives factors between
$1-\epsilon$ and $1+\epsilon$. Union with the solver event proves
the probability assertion.

For completeness, standard SDDM solves with relative energy error
can enforce the residual tests. If a batch has right-hand side
matrix $D$ and every solution has relative $A_j$-energy error $\xi$,
then, summing the squared energy errors,
\[
 \|\text{residual}\|_F
 \leq\sqrt{B_j}\,\xi
      \left(\sum_\ell D_\ell A_j^{-1}D_\ell^\top\right)^{1/2}
 \leq\xi\sqrt{B_j/a_j}\,\|D\|_F.
\]
Taking
\[
 \xi\leq\frac{\epsilon}{16\kappa_j^{3/2}
              \max\{1,\|Q\|_F,\|P\|_F/\sqrt{B_j}\}},
 \qquad\kappa_j=B_j/a_j,
\]
suffices for both tests. Equivalently, individual residual budgets
can be allocated with squared budgets summing to the displayed
Frobenius budget. These are deterministic, checkable conditions.
The standard nearly-linear SDDM solver is used as the linear-system
primitive~\cite{kyng2016}; the sketch construction follows the
factor-and-solve principle of Spielman--Srivastava~\cite{ss2011}.
Their guarantees give relative energy accuracy in time nearly linear
in nonzeros, with logarithmic dependence on $1/\xi$ and the failure
budget. For example, the high-probability factorization variant in
Appendix B of~\cite{kyng2016}, followed by its iterative refinement,
provides this primitive; SDDM systems reduce to a grounded Laplacian
by adding one cemetery vertex for the row sums. Here $\E\|Q\|_F^2=N$ and
$\E\|P\|_F^2=\tr A\leq NB$. Markov bounds make both norms at most
$\sqrt{6N/\delta}$ and $\sqrt{6NB/\delta}$, respectively, except
with combined probability $\delta/3$. Hence
$\log(1/\xi)=O(\log(N\kappa/(\epsilon\delta)))$ on that event.

There are $2dq$ solves, plus $O(d(m_E+Nq))$ sketch generation and
accumulation work. In primitive notation the time is
$O(d\sum_j 2T_{\rm SDDM}(A_j,\xi_j,\delta/(6dq))
   +d(m_E+Nq))$, yielding the stated bound.
Process one Gaussian row pair at a time and accumulate squared
coordinates into the $2Nq$ outputs; only one pair of length-$N$
right-hand sides and the solver workspace are needed at a time.
The inputs $P,Q$ are shared across all shifts.

The same shared right-hand sides also admit a multi-shift Krylov
implementation. With depth $h$ it costs
$O(dh(m_E+Nq))$ including vector recovery/updates at every shift.
The nearly-linear bound above is supplied by the SDDM primitive;
the shared unpreconditioned Krylov count retains $h$ explicitly.
General shift-dependent preconditioners do not satisfy the simple
shared-space identity used by the unpreconditioned construction.

Finally the uniform initialization event gives intervals
$[\widehat r/(1+\epsilon),\widehat r/(1-\epsilon)]$ for every entry.
Applying Theorem~\ref{thm:cache} maintains their containment along
any adaptive sequence of selected nodes. Random solver calls in
later rounds can be assigned conditional failure budgets with a
bounded sum; a union bound conditional on each preceding history
then suffices. No union over exponentially many possible selected
sets, or fresh independence of the original sketch after selection,
is required for this cache-update use.
For $0<s<1$, let $r_{i,u,j}\in[l_{i,u,j},u_{i,u,j}]$ be the
maintained intervals and put $b_0=\beta^{-1}$, or zero for hard
grounding. Intersect these intervals with the known positive
spectral ranges $[B_j^{-1},a_j^{-1}]$ and $[B_j^{-2},a_j^{-2}]$,
respectively. Each first-resolvent marginal lies in
\[
 \left[\frac{l_{2,u,j}}{b_0+u_{1,u,j}},
       \frac{u_{2,u,j}}{b_0+l_{1,u,j}}\right].
\]
Positive summation with the fixed surrogate weights gives bounds
$[L,U]$ for its marginal. By \eqref{eq:fixedrelative},
$[L,U/(1-\eta)]$ then contains the original fractional marginal.
These are valid candidate upper bounds throughout the cache event.
At $s=1$ use the same ratio at $t=0$, with $\eta=0$.
\end{proof}
\noindent Numerical check: \texttt{scripts/verify\_t6\_initialization.py}.

\begin{corollary}[Higher moments from the same Gaussian families]\label{cor:higherjl}
For every integer $\ell\geq1$, the squared columns of
$Q(A+tI)^{-\ell}$ estimate $r_{2\ell,u}(t)$, and those of
$(P+\sqrt tQ)(A+tI)^{-\ell}$ estimate $r_{2\ell-1,u}(t)$.
For any fixed maximum order $J$, all $NJq$ moments have simultaneous
relative accuracy with $d=O(\epsilon^{-2}\log(NJq/\delta))$ and
$O(dq\lceil J/2\rceil)$ solves, with logarithmic precision factors.
\end{corollary}
\begin{proof}
After multiplication the two row covariance matrices are
\[
 (A+tI)^{-2\ell}/d\quad\text{and}\quad(A+tI)^{-(2\ell-1)}/d.
\]
Apply the chi-squared tail and union bound exactly as above, over
at most $NJq$ moments. Repeated solves obtain these powers without
forming a matrix power. For an approximate chain
$V_\ell A_t\simeq V_{\ell-1}$, let $\rho_\ell$ be the Frobenius
norm of this residual, and $E_\ell$ the Frobenius error from the
exact chain with the same starting right-hand sides. Then
\[
 E_\ell\leq(E_{\ell-1}+\rho_\ell)/(a+t),\qquad E_0=0.
\]
This follows by multiplying the difference equation by $A_t^{-1}$.
For even moments impose $E^Q_\ell\leq(\epsilon/8)(B+t)^{-\ell}$;
for odd moments impose
$E^P_\ell+\sqrt t E^Q_\ell
 \leq(\epsilon/8)(B+t)^{-\ell+1/2}$.
These are sufficient because the corresponding exact target norms
are at least $(B+t)^{-\ell}$ and $(B+t)^{-\ell+1/2}$.
The same norm perturbation argument proves relative accuracy.
Choose residual budgets successively to meet these inequalities.
For fixed $J$ their relative precision is polynomial in
$\epsilon,1/\kappa,1/N,\delta$, so the nearly-linear primitive
adds only logarithmic precision factors. Each right-hand side is
solved at most $\lceil J/2\rceil$ times per shift.
\end{proof}
\noindent Numerical check: \texttt{scripts/verify\_t6\_higher\_moments.py}.

\begin{proposition}[Direct forest-count estimators and their variance]\label{prop:forest}
For $A_t=L+tI$, $t>0$, let $\rho(u)$ be the root of $u$ in a
rooted forest sampled with weight proportional to
$t^{\#\text{roots}}\prod_{e\text{ in forest}}w_e$.
Then $\Prob(\rho(u)=v)=t(A_t^{-1})_{uv}$.
In particular the self-root indicator divided by $t$ estimates
$r_{1,u}$, and the collision indicator for $u$'s roots in two
independent forests, divided by $t^2$, estimates $r_{2,u}$.
For unit-weight complete graphs and $0<t\leq1$, both Bernoulli
probabilities lie in $[1/N,2/N]$. Each direct sample-mean estimator
therefore needs $\Omega(N/\epsilon^2)$ independent samples even
to make its relative root-mean-square error at most $\epsilon$.
\end{proposition}
\begin{proof}
Use Wilson's exact rooted-forest sampler, equivalently add a cemetery
vertex with edge weight $t$ to every vertex and delete it after
sampling its rooted tree~\cite{wilson1996}. In this sampler $u$
may be processed first, by the order invariance of cycle popping.
Its eventual root is the last vertex of its walk before killing;
loop erasure preserves that endpoint. The continuous-time walk
has generator $-L$ and independent killing rate $t$.
Its absorption distribution is
$\int_0^\infty t e^{-tz}(e^{-zL})_{uv}dz=t(A_t^{-1})_{uv}$.
This also follows from the first-step equations
$(L+tI)H=tI$. Independence of two forests gives collision probability
$\sum_v(tR_{uv})^2=t^2r_{2,u}$, using symmetry of $R$.

For the complete graph the Laplacian is $N(I-\ind\ind^\top/N)$.
The self-root and collision probabilities are therefore
\[
 p_1=\frac1N+\left(1-\frac1N\right)\frac{t}{N+t},\qquad
 p_2=\frac1N+\left(1-\frac1N\right)\left(\frac{t}{N+t}\right)^2.
\]
For $0<t\leq1$ both are in $[1/N,2/N]$.
The mean of $R_s$ independent Bernoulli samples with success
probability $p$ has relative mean-square error $(1-p)/(R_sp)$.
Consequently $R_s\geq(N/2-1)/\epsilon^2$ is necessary.
Scaling by $1/t$ or $1/t^2$ does not change relative error.
For homogeneous killing, a simple upper bound on Wilson walk work
per forest is $N(1+d_{\max}/t)$: each walk step is absorbed with
probability at least $t/(d_{\max}+t)$, and at most $N$ walks start.
Hitting an existing tree only shortens these walks.
\end{proof}
\noindent Numerical check: \texttt{scripts/verify\_t6\_forest\_comparison.py}.

\paragraph{Initialization choice and comparison.}
The main initialization uses the two-factor Gaussian construction:
it estimates both diagonal families directly, shares $2d$ right-hand
sides across shifts, and has an SDDM-solver complexity independent of
individual root probabilities. The continuum theorem provides a
second derivation of the suggested $r_2$-integration route, with
additional positive-quadrature shifts when implemented finitely.
For $0<s<1$, the fixed surrogate uses first inverse powers, so these
two cached moments are exactly the quantities in its scalar ratios.
At $s=1$ only the zero shift is needed.

Rooted-forest sampling changes its killing distribution with the
shift. ForestCFCM and SchurCFCM~\cite{xia2025} use voltage/path
estimators, JL compression, adaptive error tests, and, in the latter
case, auxiliary roots and a Schur complement. Their Lemma 3.9 gives
a uniform sample bound containing
$\epsilon^{-2}\operatorname{diam}(G)^2
 d_{\max}(S)^{2\operatorname{diam}(G)+2}\log N$;
their walk cost is controlled by an absorbing-walk trace.
This identifies complementary graph-dependent costs to measure in
the later experimental comparison. The construction above supplies
uniform relative diagonal guarantees for arbitrary weighted SDDM
inputs, while Theorem~\ref{thm:radau} supplies deterministic
candidate-level CM marginal certificates.

\section{T7: integer marginals and certified lazy greedy}
This section analyzes stale-bound CELF; the fresh all-node variant
in T7b supplies the worst-case nearly-linear-per-round guarantee.
T12 analyzes the maintained-sketch hybrid used by the final engine.
The CELF algorithm retains upper bounds by supermodularity and
refreshes candidates from the current matrix. The incremental
two-moment updates of T5 are an optional acceleration, separate
from the guarantee below. We first resolve integer-power evaluation.
In this section $p\geq1$ denotes an integer exponent, and $m_E$
denotes the graph's edge count.

\begin{proposition}[Positive integer marginal representation]\label{prop:integerpositive}
At any nonnegative shift let $R=(A+tI)^{-1}$, $x=Rb$,
$c=(\beta^{-1}+r_1)^{-1}$, using $\beta^{-1}=0$ for hard grounding,
and $K=R-cxx^\top\succeq0$. Define
\[
 a_j=cr_{j+1},\qquad d_j=cx^\top K^{j-1}x\quad(j\geq1),
 \qquad H_0=1,\quad H_j=\sum_{i=1}^j d_i H_{j-i}.
\]
The integer marginal admits the noncancelling formula
\begin{equation}\label{eq:integerpositive}
 q_p(t)=\sum_{j=1}^p j d_j H_{p-j},\qquad
 a_p\leq q_p(t)\leq p a_p.
\end{equation}
All terms are nonnegative. In particular, with $\nu$ the probability
measure defined below,
\begin{equation}\label{eq:q2variance}
 q_2=2a_2-a_1^2=a_2+\operatorname{Var}_{\nu}(Z),
 \qquad a_1^2\leq a_2.
\end{equation}
\end{proposition}
\begin{proof}
For soft grounding $K=(A+tI+\beta bb^\top)^{-1}$.
For hard grounding it is the zero-padded inverse of the deleted
principal matrix. These facts follow from the T3 block identities,
so $q_p=\tr R^p-\tr K^p$ and $K\succeq0$ in both cases.
Write $D(z)=\sum_{j\geq1}d_jz^j=czx^\top(I-zK)^{-1}x$.
For sufficiently small $|z|$, the determinant lemma gives
\[
 \frac{\det(I-zK)}{\det(I-zR)}=\frac1{1-D(z)},\qquad
 \sum_{p\geq1}\frac{q_p}{p}z^p=-\log(1-D(z)).
\]
Every $d_j$ is nonnegative by positive semidefiniteness of $K$.
The reciprocal power series has coefficients $H_j$ satisfying the
stated positive recurrence. Differentiating the logarithmic series
proves \eqref{eq:integerpositive}'s equality.

Alternatively expand words in $R=K+cxx^\top$ inside $\tr R^p$.
Marking one of the $\ell$ occurrences of $cxx^\top$ and cyclically
rotating the trace gives
\[
 q_p=\sum_{\ell=1}^p\frac p\ell
       \sum_{\substack{j_1+\cdots+j_\ell=p\\j_i\geq1}}
       d_{j_1}\cdots d_{j_\ell}.
\]
This counting identity follows because there are $p$ possible
positions for a marked occurrence and each word is counted $\ell$
times. The same expansion in $cx^\top R^{p-1}x=cr_{p+1}=a_p$
gives the identical sum without the factor $p/\ell$.
Every factor is nonnegative and $1\leq p/\ell\leq p$, proving
both magnitude bounds.

For the variance form diagonalize $R$, with eigenvalues $z_i>0$
and spectral weights $w_i$ of $b$. Put mass $cw_i z_i$ at $z_i$
and mass $1-cr_1$ at zero. These masses sum to one; in the soft
case the zero mass equals $\beta^{-1}/(\beta^{-1}+r_1)$, and it
is zero in the hard case. The resulting distribution $\nu$ has
$\E_\nu Z^j=a_j$. Thus $a_1^2\leq a_2$, and the formula for
$q_2$ from T3 becomes \eqref{eq:q2variance}. The subtracted term
is at most half of $2a_2$, so second-order cancellation is uniformly
bounded, independently of dimension and condition number.
\end{proof}
\noindent Numerical check: \texttt{scripts/verify\_t7\_integer\_positive.py}.

\begin{theorem}[Relative moment-to-marginal certificate]\label{thm:momentstability}
Suppose $|\widehat r_j/r_j-1|\leq\zeta<1$ for
$1\leq j\leq p+1$. Evaluate $\widehat q_p$ by the scalar logarithm
recurrence of T3 from these moments and the exact penalty. Then
\begin{equation}\label{eq:momentcondition}
 \left|\frac{\widehat q_p}{q_p}-1\right|
 \leq C_p(\zeta):=2^p((1-\zeta)^{-p}-1).
\end{equation}
In particular a desired relative marginal error $\eta_0<1$ is
certified by
\[
 \zeta\leq1-(1+\eta_0/2^p)^{-1/p}.
\]
The absolute-term condition number of the exact logarithmic
polynomial is at most $2^p-1$. Independently, moment intervals
$r_1\in[l_1,u_1]$, $r_{p+1}\in[l_{p+1},u_{p+1}]$, $l_1>0$,
give the cancellation-free coarse marginal interval
\begin{equation}\label{eq:coarseinteger}
 \left[\frac{l_{p+1}}{\beta^{-1}+u_1},
       \frac{p u_{p+1}}{\beta^{-1}+l_1}\right].
\end{equation}
\end{theorem}
\begin{proof}
The Taylor series $F(t-z)/F(t)$ from T3 equals
$1+\sum_{j\geq1}a_jz^j$, where $a_j=cr_{j+1}$.
Hence
\[
 q_p=p\sum_{\ell=1}^p\frac{(-1)^{\ell+1}}\ell
       \sum_{\substack{j_1+\cdots+j_\ell=p\\j_i\geq1}}
       a_{j_1}\cdots a_{j_\ell}.
\]
For the probability measure in Proposition~\ref{prop:integerpositive},
Holder's inequality gives $a_j\leq a_p^{j/p}$ for $j\leq p$.
Thus every product in the inner sum is at most $a_p$.
The number of compositions is $\binom{p-1}{\ell-1}$, so the sum
of absolute terms is at most
$a_p\sum_{\ell=1}^p(p/\ell)\binom{p-1}{\ell-1}
 =(2^p-1)a_p\leq(2^p-1)q_p$.

The perturbed reciprocal coefficient satisfies
$(1+\zeta)^{-1}\leq\widehat c/c\leq(1-\zeta)^{-1}$.
Consequently $\theta^{-1}\leq\widehat a_j/a_j\leq\theta$,
where $\theta=(1+\zeta)/(1-\zeta)$.
A product of $\ell$ factors changes by at most
$(\theta^\ell-1)$ times its true value, even when the perturbed
moments are not a consistent moment sequence. Therefore
\[
 |\widehat q_p-q_p|
 \leq a_p\sum_{\ell=1}^p\binom p\ell(\theta^\ell-1)
 =a_p((1+\theta)^p-2^p)\leq q_p C_p(\zeta).
\]
Solving $C_p(\zeta)\leq\eta_0$ proves the precision choice.
Finally \eqref{eq:coarseinteger} follows directly from
$cr_{p+1}\leq q_p\leq pcr_{p+1}$ and positivity of the denominator.
For $p=2$, \eqref{eq:momentcondition} is
$4((1-\zeta)^{-2}-1)=8\zeta+O(\zeta^2)$.
These bounds propagate input-moment uncertainty; finite-precision
evaluation can additionally use the positive recurrence or validated
scalar arithmetic.
\end{proof}
\noindent Numerical check: \texttt{scripts/verify\_t7\_moment\_stability.py}.

\paragraph{Initialization and permanent keys.}
Fix the admissible seed $S_0$ and let $W=V\setminus S_0$ contain
$N$ eligible vertices. The budget $1\leq k\leq N$ counts additional vertices.
For integer exponent $p$, the only shift is zero and the initial
moments are $r_1,\ldots,r_{p+1}$. A relative moment estimate with
the precision in Theorem~\ref{thm:momentstability} gives a marginal
estimate $\widehat\Delta_v$ with error $\eta_0$; initialize the key
$U_v=\widehat\Delta_v/(1-\eta_0)$. The coarse upper endpoint in
\eqref{eq:coarseinteger} is also sufficient and needs only $r_1$
and $r_{p+1}$. This choice uses a constant moment tolerance even
when $p$ grows.

For noninteger $s$, let $\widetilde f_s$ be the fixed positive rule
with relative marginal error $\eta_q$. Apply the integer moment
certificate to every shifted kernel and sum with its positive
weights. If this gives an estimate $\widehat\Delta_v$ of the
surrogate marginal with error $\eta_0$, a valid original-objective
key is
\[
 U_v=\frac{\widehat\Delta_v}{(1-\eta_0)(1-\eta_q)}.
\]
All these intervals hold simultaneously on the single initial JL
and solver event. Their tolerances may be fixed constants; final
greedy accuracy is controlled by the deterministic refreshes.
For a general CM function, deterministic initial keys are available
directly: $f(a)$ for hard grounding and
$\min\{\beta g(a),f(a)-f_\infty\}$ for soft grounding. These follow
from the spectral count $0\leq\xi\leq1$ and $\int\xi=\beta$
in the soft case. They include the heat objective.

\begin{theorem}[Certified CELF with permanent upper bounds]\label{thm:celf}
Suppose the initial keys bound every marginal at $S_0$ with
probability at least $1-\delta$. Let $0<\theta<1$.
Maintain a max-heap of unselected vertices with their current keys.
At each round, use the following rule:
\begin{enumerate}
\item Pop a vertex $v$ with largest key. If it has not been refreshed
in this round, run Lanczos once from $e_v$ on the current grounded
matrix, extending that run until T4 gives
$L_v\leq\Delta_v(S)\leq H_v$ with $L_v\geq(1-\theta)H_v$.
Replace its key by $\min\{U_v,H_v\}$ and record the round.
\item Let $U_{\rm next}$ be the largest key left in the heap, or
zero if it is empty. If $L_v\geq(1-\theta)U_{\rm next}$, select
$v$. Otherwise return it to the heap and repeat.
\end{enumerate}
On the initial event every stored key remains valid throughout all
adaptive rounds. Every accepted vertex has marginal at least
$(1-\theta)$ times the current largest marginal. If $R$ counts
candidate/round refreshes, then
\begin{equation}\label{eq:Rbound}
 R\leq\sum_{j=0}^{k-1}(N-j)=kN-\frac{k(k-1)}2.
\end{equation}
For $F(U)=\Phi_f(S_0)-\Phi_f(S_0\cup U)$ and
$\mathrm{OPT}_k=\max_{|U|\leq k}F(U)$, the returned set satisfies
\begin{equation}\label{eq:lazyguarantee}
 F(U_k)\geq\left[1-\left(1-\frac{1-\theta}{k}\right)^k\right]
                  \mathrm{OPT}_k
 \geq(1-e^{-1}-\theta)\mathrm{OPT}_k
\end{equation}
with probability at least $1-\delta$. In particular set
$\theta=\epsilon$ for the requested guarantee.
\end{theorem}
\begin{proof}
Supermodularity gives
$\Delta_v(S')\leq\Delta_v(S)$ whenever $S\subseteq S'$ and
$v\notin S'$. Hence every valid old upper bound is a valid future
upper bound. A fresh T4 upper bound is valid on the current matrix,
so the minimum of old and fresh bounds is valid too. Induction
proves this invariant along the entire adaptive history on the
single initial event. No cache-update error or adaptive reuse of a
Gaussian sketch is involved.

At acceptance every other true marginal is at most
$U_{\rm next}$, while the chosen marginal is at least $L_v$.
This proves the $(1-\theta)$ maximum-marginal property.
If a freshly evaluated vertex fails the acceptance test, then
$H_v\leq L_v/(1-\theta)<U_{\rm next}$, so its new key is strictly
below another key. A vertex already refreshed in this round, when
it later becomes heap-top, necessarily passes: its stored lower
bound is at least $(1-\theta)$ times its key, and that key is at
least the next key. Thus no vertex is refreshed twice in a round,
and the procedure terminates after at most $N-j$ refreshes in
round $j$. Ties use the non-strict acceptance inequality.

T4 intervals converge and are exact at invariant-subspace
termination. For its integration-free interval $H=L+E$, it is
sufficient that $E\leq\theta d_*$, where $d_*$ is T4's positive
lower bound on the marginal: since $H\geq\Delta\geq d_*$,
$L=H-E\geq(1-\theta)H$. The positive-marginal cases follow from
T1; zero or constant objectives can be selected directly. Therefore
each refresh is one finite, possibly extended, Lanczos run.

The gain $F$ is normalized, nonnegative, monotone and submodular.
For an optimal additional set $O$ with $|O|\leq k$, telescoping
its additions and using diminishing returns gives
\[
 \mathrm{OPT}_k-F(U_j)
 \leq F(U_j\cup O)-F(U_j)
 \leq\sum_{v\in O\setminus U_j}\Delta_v(S_0\cup U_j)
 \leq k\max_v\Delta_v(S_0\cup U_j).
\]
Each accepted step closes at least $(1-\theta)/k$ of this gap.
Induction gives the first inequality in \eqref{eq:lazyguarantee}.
The bound $(1-(1-\theta)/k)^k\leq e^{-(1-\theta)}$ and
$e^{\theta-1}-e^{-1}=\int_0^\theta e^{x-1}dx\leq\theta$ give
the second one. The proof is deterministic conditional on the
initial event, so the final failure probability is at most $\delta$.
\end{proof}
\noindent Numerical check: \texttt{scripts/verify\_t7\_celf.py}.

\begin{theorem}[Refresh-sensitive complexity]\label{thm:lazycost}
Let $n$ be the largest active matrix dimension, $m_E$ its initial
edge count, and let $h_r$ be the Lanczos depth of refresh $r$.
Let $J$ be the largest initial moment order, $q$ the initial shift
count, $\zeta$ the required moment tolerance, and
$d=O(\zeta^{-2}\log(nqJ/\delta))$.
With $L=\lceil J/2\rceil$, initialization uses $2dqL$ SDDM
right-hand-side solves and polynomially small energy tolerances.
In terms of the SDDM solve primitive, a valid bound is
\begin{align}\label{eq:lazycost}
 T_{\rm init}
 &=O\left(dqL\,T_{\rm SDDM}(\xi,\delta/(6dqL))
       +d(m_E+nqL)+nqJ^2\right),\notag\\
 T_{\rm total}
 &=T_{\rm init}
   +O\left(\sum_{r=1}^R\bigl((m_E+n)h_r+C_f(h_r)\bigr)
       +(N+R+k)\log(N+1)\right).
\end{align}
For integer $s=p$, $q=1$, $J=p+1$, and scalar work can be bounded
by $C_f(h)=O(ph+p^2\log(h+1))$. For a general CM scalar evaluation
oracle, a conservative choice is $C_f(h)=O(h^3)$ plus scalar
function-evaluation costs. For fixed integer $p$ and constant
coarse-initialization tolerance, the SDDM part is nearly linear;
the refresh part retains $R$ and $h=\max_r h_r$ explicitly as
$O(R(m_E+n)h)$ plus its displayed scalar work.
For powers, T4 supplies the sufficient depth
$O(\sqrt{\kappa_*}((s+1)\log\kappa_*+\log(1/\theta)))$,
capped by the current dimension, where $\kappa_*$ bounds the
spectral intervals used by the refreshes. Doubling changes only
the constant in this depth bound.
\end{theorem}
\begin{proof}
T6's two Gaussian families are advanced through $L$ inverse powers
at each shift. That is $2dqL$ right-hand-side solves; simultaneous
JL concentration is over at most $nqJ$ scalar moments. Accumulating
moments costs $O(dnqL)$, generating the original RHS pairs costs
$O(d(m_E+n))$, and evaluating the logarithmic recurrence at every
node/shift costs $O(nqJ^2)$. Coarse integer keys need just the two
endpoint moments and $O(n)$ final scalar work instead.

An explicit precision dependence also follows from T6. Normalize
the largest eigenvalue to one, put $\kappa=B/a$, and suppose the
initial Gaussian RHS Frobenius norms are bounded by $M$ after
normalization. A sufficient common relative energy tolerance is
\[
 \xi=\frac{\zeta}{64L\kappa^{2L}(M+1)}.
\]
The exact inverse-power chain has norm at most $\kappa^{L-1}M$
before its last solve. Inductively its approximate norm is at most
$\kappa^{L-1}(M+1)$. The energy-to-residual inequality in T6
then bounds each residual by
$\zeta/(64L\kappa^{L+1/2})\leq\zeta/(64L\kappa^L)$.
The recurrence $E_j\leq\kappa(E_{j-1}+\rho_j)$ bounds every
chain error by $\zeta/64$. This meets T6's norm-error budgets;
after shifting, normalized $\sqrt t$ coefficients are at most one.
Thus $\log(1/\xi)=O(L\log\kappa+log(L(M+1)/\zeta))$.
The Gaussian norm event gives $M=O(\sqrt{n/\delta})$ with the
allocated failure probability. Fixed $J$ yields logarithmic
precision factors; growing $J$ retains its $L\log\kappa$ term.

Each refresh uses $h_r$ masked sparse products and vector
recurrences, costing $O((m_E+n)h_r)$. Evaluate stopping conditions
at doubling depths and at breakdown. For integer powers, repeated
tridiagonal solves compute $r_1(T),\ldots,r_{p+1}(T)$ in $O(ph)$,
and their logarithmic recurrence costs $O(p^2)$. Summing over the
doubling checks gives $C_f(h)$ as stated. The positive small-matrix
representation is also available. General CM functions may instead
use the small trace differences and T4's analytic width, giving
the conservative cubic bound. At full termination the complement's
trace cancels exactly as in T3.

Every unsuccessful fresh evaluation causes one heap reinsertion;
every successful selection removes one vertex. A previously fresh
vertex is accepted when it returns to the top. Therefore heap
operations number $O(N+R+k)$ and each costs $O(\log(N+1))$.
Grounding is implemented by diagonal masks or restricted coordinates,
so no dense inverse maintenance is required. The three-term Lanczos
count is in exact arithmetic; explicit full reorthogonalization
adds $O(n\sum_r h_r^2)$ if chosen. This proves \eqref{eq:lazycost}.
\end{proof}
\noindent Numerical check: \texttt{scripts/verify\_t7\_complexity.py}.

\paragraph{The $s=1$ comparison.}
At $s=1$, this algorithm has one zero shift and two initial moments.
For constant initialization tolerance its cost is
$\widetilde O(m_E+n)+O((m_E+n)\sum_r h_r)$, plus heap and scalar
work. ApproxGreedy's original analysis~\cite{li2019}, including
its explicitly stated error dependence, gives
$\widetilde O(km_E\epsilon^{-2})$ using nearly-linear SDDM solves
to re-estimate all marginals each round. These sparse-work bounds
match in order when $\sum_r h_r=\widetilde O(k\epsilon^{-2})$;
the refresh-sensitive bound keeps the realized sum explicit.
The unconditional CELF refresh bound is \eqref{eq:Rbound}.
T7b supplies an explicit fresh all-node variant with an unconditional
$\widetilde O(km_E\epsilon^{-2})$ bound at $p=1$.

ForestCFCM and SchurCFCM~\cite{xia2025}, Theorems 3.11 and 4.7,
display the graph-dependent bound
$\widetilde O(k\epsilon^{-2}n D^2d_{\max}^{2D+2})$ under their
stated realistic-graph walk-cost assumptions, where $D$ is graph
diameter. Their adaptive sampling and Schur auxiliary roots target
these graph-dependent costs. The quantities to compare with our
refresh work are their forest sample counts and walk costs versus
$\sum_r h_r$ and sparse products. The original CFCC papers also
select the first leader; their free-first-leader guarantees use a
different normalization. Equation~\eqref{eq:lazyguarantee} compares
sets of $k$ additions to the same fixed admissible seed $S_0$.

\section{T7b: fresh all-node sketches and a polynomial exponent factor}
The positive reciprocal series in T7 sharpens the earlier
absolute-monomial perturbation bound. This sharpening uses only
moment accuracy; perturbed moments may be mutually inconsistent.

\begin{theorem}[Polynomial-precision integer marginal estimation]
\label{thm:polymoment}
Use the notation of Proposition~\ref{prop:integerpositive}, at any
fixed shift. If $|\widehat r_j/r_j-1|\leq\zeta<1$ for
$j=1,\ldots,p+1$, the T3 logarithmic recurrence satisfies
\begin{equation}\label{eq:polymoment}
 \left|\frac{\widehat q_p}{q_p}-1\right|
 \leq \left(1+\frac{4\zeta}{1-\zeta}\right)^p-1.
\end{equation}
For $0<\eta\leq1$, choosing $\zeta\leq\eta/(16p)$ makes this
error at most $\eta$. The assertion holds for soft and hard grounding,
independently of dimension and condition number.
\end{theorem}
\begin{proof}
Work with formal power series through degree $p$. Set
\[
 F(z)=1+\sum_{j\geq1}a_jz^j,\qquad
 D(z)=\sum_{j\geq1}d_jz^j.
\]
The determinant identities in Proposition~\ref{prop:integerpositive}
give $F^{-1}=1-D$, with $a_j,d_j\geq0$, and hence
\begin{equation}\label{eq:renewalmoment}
 a_j=d_j+\sum_{i=1}^{j-1}a_i d_{j-i}.
\end{equation}
Also $q_p=p[z^p]\log F$, $q_p\geq a_p$, and Holder's inequality
for the probability measure constructed there gives
$a_{j_1}\cdots a_{j_\ell}\leq a_p$ whenever
$j_1+\cdots+j_\ell=p$ and $j_i\geq1$.

Put $\widehat a_j=\widehat r_{j+1}/(\beta^{-1}+\widehat r_1)$.
Positivity of the exact penalty yields
\[
 \frac{1-\zeta}{1+\zeta}\leq\frac{\widehat a_j}{a_j}
 \leq\frac{1+\zeta}{1-\zeta},\qquad
 |\widehat a_j-a_j|\leq\alpha a_j,\quad
 \alpha=\frac{2\zeta}{1-\zeta}.
\]
For $E(z)=\sum_{j=1}^p(\widehat a_j-a_j)z^j$, define
$H=E/F=E(1-D)$. For $1\leq j\leq p$, identity
\eqref{eq:renewalmoment} implies
\[
 |[z^j]H|\leq\alpha\left(a_j+
          \sum_{i=1}^{j-1}a_i d_{j-i}\right)
 =\alpha(2a_j-d_j)\leq2\alpha a_j.
\]
The formal identity $\log(F+E)-\log F=\log(1+H)$ then gives
\begin{align*}
 |\widehat q_p-q_p|
 &\leq p\sum_{\ell=1}^p\frac{(2\alpha)^\ell}{\ell}
   \sum_{\substack{j_1+\cdots+j_\ell=p\\j_i\geq1}}
            a_{j_1}\cdots a_{j_\ell}\\
 &\leq a_p\sum_{\ell=1}^p\binom p\ell(2\alpha)^\ell
 =a_p((1+2\alpha)^p-1).
\end{align*}
Division by $q_p\geq a_p$ proves \eqref{eq:polymoment}.
All identities concern finitely many coefficients, so convergence
of a series or realizability of the perturbed moments is unnecessary.
Finally $\zeta\leq\eta/(16p)$ implies
$4p\zeta/(1-\zeta)\leq4\eta/15$.
Thus the bound is at most $e^{4\eta/15}-1\leq\eta$:
convexity bounds it by $\eta(e^{4/15}-1)$ on $[0,1]$, and
$e^{4/15}\leq(1-4/15)^{-1}=15/11<2$.
\end{proof}
\noindent Numerical check: \texttt{scripts/verify\_t7b\_stability.py}.

\paragraph{FreshRoundGreedy (all-node worst-case variant).}
Fix a seed $S_0$, $k$ additional vertices, integer $p\geq1$,
$0<\epsilon\leq1/2$, and $0<\delta<1$. Write $n$ for an upper
bound on the active dimension and set
\begin{equation}\label{eq:freshparams}
 \eta=\epsilon/2,\quad\zeta=\epsilon/(32p),\quad
 J=p+1,\quad L=\lceil J/2\rceil,\quad
 d=\left\lceil128\zeta^{-2}\log\frac{12nJk}{\delta}\right\rceil.
\end{equation}
At each of the $k$ rounds:
\begin{enumerate}
\item Form the current SDDM matrix $A$ and its sparse factor
$A=C^\top C$ as in T6. Draw entirely new independent Gaussian
families $Q,Z$, with variance $1/d$ entries, and put $P=ZC$.
These families and the solver randomness are independent of all
previous rounds.
\item Apply T6's repeated solves to these families through depth
$L$, obtaining every $r_{j,v}=e_v^\top A^{-j}e_v$, $1\leq j\leq J$,
with relative tolerance $\zeta$. Specifically, squared columns of
$QA^{-\ell}$ estimate order $2\ell$, and those of $PA^{-\ell}$
estimate order $2\ell-1$. Discard orders above $J$.
\item Evaluate $\widehat\Delta_v=\widehat q_{p,v}(0)$ for every
eligible vertex by the scalar logarithmic recurrence, select a
maximizer, and ground it. Discard this round's sketches.
\end{enumerate}
Each round's conditional failure budget is $\delta/k$, including
all its candidate estimates and all its solves.

\begin{theorem}[Fresh-round all-node greedy guarantee]\label{thm:freshgreedy}
FreshRoundGreedy returns $k$ additional vertices $U_k$ with
\begin{equation}\label{eq:freshgain}
 \Phi_p(S_0)-\Phi_p(S_0\cup U_k)
 \geq(1-e^{-1}-\epsilon)
   \max_{|U|\leq k}\bigl(\Phi_p(S_0)-\Phi_p(S_0\cup U)\bigr)
\end{equation}
with probability at least $1-\delta$, for hard and soft grounding.
\end{theorem}
\begin{proof}
Condition on the complete history before round $i$. Its matrix,
factor, and eligible vertices are now fixed. The new Gaussian
families therefore have exactly the covariance and chi-squared
laws used in T6. There are at most $nJ$ moment estimates;
\eqref{eq:chisquare}, with deviation $\zeta/4$, and
\eqref{eq:freshparams} give conditional JL failure at most
\[
 2nJ\exp(-d\zeta^2/128)\leq\delta/(6k).
\]
The residual/norm budgets in Theorem~\ref{thm:freshcost} give
conditional failure at most $\delta/(3k)$ for each of the Gaussian
norm and solver events. Their intersection implies relative moment
error at most $\zeta$, by T6. In particular the total conditional
failure is at most $\delta/k$, for every history.

Let $E_i$ be this round's simultaneous event. Taking conditional
expectations and then a union bound gives
\[
 \Prob\left(\bigcup_{i=1}^k E_i^c\right)
 \leq\sum_{i=1}^k\E\bigl[\Prob(E_i^c\mid\mathcal H_{i-1})\bigr]
 \leq\delta.
\]
Independence between the events $E_i$ themselves is unnecessary;
it is the fresh random draws conditional on the history that
justify each probability bound.

On their intersection, Theorem~\ref{thm:polymoment} gives
$(1-\eta)\Delta_v\leq\widehat\Delta_v\leq(1+\eta)\Delta_v$
for all eligible vertices in every round. If $v$ maximizes the
estimate and $u$ maximizes the true marginal, then
\[
 \Delta_v\geq\frac{\widehat\Delta_v}{1+\eta}
 \geq\frac{1-\eta}{1+\eta}\Delta_u
 \geq(1-\epsilon)\Delta_u.
\]
The normalized seeded gain is monotone and submodular by T1.
The same gap-closing induction as in Theorem~\ref{thm:celf}
therefore gives the stronger factor
$1-(1-(1-\epsilon)/k)^k\geq1-e^{-1}-\epsilon$.
For hard grounding the last allowed deletion, if any, uses the
one-dimensional marginal and leaves zero trace; no inverse of
an empty matrix is formed.
\end{proof}
\noindent Numerical check: \texttt{scripts/verify\_t7b\_greedy.py}.

\begin{theorem}[Worst-case cost with explicit polynomial dependence]
\label{thm:freshcost}
Suppose certified bounds $0<a\leq\lambda_{\min}(A)$ and
$\lambda_{\max}(A)\leq B$ hold for every round, and put $\kappa=B/a$.
Using the nearly-linear SDDM primitive of T6, FreshRoundGreedy has
\begin{equation}\label{eq:freshcost}
 T=\widetilde O\bigl(p^4 k(m_E+n)\epsilon^{-2}\bigr).
\end{equation}
Thus one explicit choice is $f(p)=p^4$; on connected graphs
$m_E+n=O(m_E)$ gives the requested
$\widetilde O(f(p)km_E\epsilon^{-2})$.
The suppressed factors are logarithmic in
$n,k,p,\kappa,1/\epsilon,1/\delta$, with no suppressed polynomial
factor in $p$. The count of SDDM right-hand-side solves alone is
$O(kdL)=\widetilde O(p^3k\epsilon^{-2})$.
These are input-wise worst-case bounds, independent of the observed
CELF refresh count, in the standard capped high-probability solver
model; solver failure is included in $\delta$.
\end{theorem}
\begin{proof}
Normalize every matrix by the common $B$, and also replace the
soft penalty by $\beta/B$. All original marginal values acquire
the same positive factor $B^{-p}$ on converting back, so every
selection is unchanged. The normalized spectrum is in
$[1/\kappa,1]$. Its factor satisfies
$\E\|P\|_F^2=\tr A\leq n$ and $\E\|Q\|_F^2\leq n$.
Set
\begin{equation}\label{eq:freshprecision}
 M=\sqrt{6nk/\delta},\qquad
 \xi=\frac{\zeta}{64L\kappa^{2L}(M+1)}.
\end{equation}
Markov's inequality gives simultaneous starting norms at most $M$
except with conditional probability $\delta/(3k)$. Use relative
energy accuracy $\xi$ at every inverse-power solve, with per-solve
conditional failure at most $\delta/(6kdL)$. There are at most
$2dL$ such solves per round. The norm test and a capped solve budget
may terminate with an arbitrary feasible answer on failure, so the
work bound also holds on unsuccessful executions.

For either family, let $V_0$ be its starting batch and $\widehat V_j$
its approximate chain. Write $\rho_j=\|\widehat V_j A-
\widehat V_{j-1}\|_F$ and
$E_j=\|\widehat V_j-V_0A^{-j}\|_F$. Relative energy accuracy
implies $\rho_j\leq\xi\sqrt\kappa\|\widehat V_{j-1}\|_F$.
Inductively $\|\widehat V_{j-1}\|_F\leq
\kappa^{j-1}(M+1)$, whence for $j\leq L$
\[
 \rho_j\leq\frac{\zeta}{64L\kappa^{L+1/2}}
 \leq\frac{\zeta}{64L\kappa^L},\qquad
 E_j\leq\kappa(E_{j-1}+\rho_j)\leq\zeta/64.
\]
The last inequality follows by summing at most $L$ terms, each
amplified by at most $\kappa^L$, starting with $E_0=0$.
It also closes the stated induction on approximate batch norms.
Every normalized target moment is at least one. Thus these
absolute column errors, combined with the relative $\zeta/4$
JL event, meet T6's sufficient $\zeta/8$ norm budgets and give
relative moment error at most $\zeta$.

There are at most $2dL$ solves per round, each on a matrix with
$O(m_E+n)$ nonzeros. Sparse Gaussian-factor multiplication and
moment accumulation cost $O(d(m_E+nL))$ per round. The scalar
recurrences cost $O(np^2)$ and the maximization costs $O(n)$.
The precision satisfies
\[
 \log(1/\xi)=2L\log\kappa+
       \log\bigl(64L(M+1)/\zeta\bigr).
\]
A nearly-linear SDDM solve costs
$\widetilde O((m_E+n)\log(1/\xi))$, including logarithmic
failure-probability factors. Since $d=\widetilde O(p^2\epsilon^{-2})$
and $L\leq p$, explicitly counting this precision cost gives
\[
 \widetilde O\bigl(k(m_E+n)p^2L^2\epsilon^{-2}\bigr)
       +O(knp^2)
 \subseteq\widetilde O\bigl(p^4k(m_E+n)\epsilon^{-2}\bigr).
\]
The $p^3$ solve count and the extra factor from solve precision are
therefore distinguished. The argument applies to every input and
every adaptive matrix sequence on the simultaneous success event.
\end{proof}
\noindent Numerical check: \texttt{scripts/verify\_t7b\_complexity.py}.

\paragraph{Algorithmic role and the $p=1$ comparison.}
CELF supplies practical acceleration by avoiding most candidate
refreshes and using permanent upper bounds. FreshRoundGreedy
supplies the worst-case bound \eqref{eq:freshcost}. Each has its
own proved guarantee. At $p=1$, $J=2$, $L=1$, and
$\Delta_v=r_{2,v}/(\beta^{-1}+r_{1,v})$, so fresh all-node
estimation uses two Gaussian solve families per round and costs
$\widetilde O(km_E\epsilon^{-2})$ on connected graphs. This is
the same asymptotic order as Li et al.'s ApproxGreedy~\cite{li2019}.
The comparison is between the marginal-selection phases;
\eqref{eq:freshgain} uses the common prescribed seed $S_0$,
whereas their complete CFCC procedure also chooses its first leader.
For growing $p$, the positive reciprocal representation replaces
the earlier $\zeta=O(\eta/(p2^p))$ sufficient tolerance by
$\zeta=O(\eta/p)$, eliminating the $4^p$ sketch-size factor.

\section{T8: positioning}
Li et al.~\cite{li2019} establish supermodularity and scalable greedy
selection for group current-flow closeness, the hard-grounding
$s=1$ member of our family. Xia and Zhang~\cite{xia2025} accelerate
that objective through rooted forests and Schur complements, while
Wang et al.~\cite{wang2026} study group Markov centrality based on
hitting times. We study node grounding for the full completely
monotone trace family. Its $s=1$ member enables direct comparison
on group current-flow closeness, while higher inverse powers
increase sensitivity to poorly grounded spectral modes.
For hard grounding, the computational core maintains current-matrix
moment sketches under node removal and retains refreshed moments
through exact algebraic updates. A uniform sketch event supports
adaptive reuse and budgeted candidate refinement. Integer marginals
admit positive representations and explicit moment-error bounds;
fresh-round sketches and Gauss--Radau lazy certification provide
complementary guaranteed algorithms. FastFrac~\cite{fastfrac}
instead adds edges
under positive fractional traces, and POLAR~\cite{polar} treats
negative fractional edge objectives through submodularity-ratio
analysis. Grounding changes killing rather than transition rates,
yielding the stronger diminishing-drop structure used here.
The inverse-power spectral limit connects this family to the
grounded-eigenvalue objective of Zhou et al.~\cite{zhou2023}, with
explicit fixed-budget baseline-aware and augmented-budget spectral
guarantees.

\noindent Numerical checks: \texttt{scripts/verify\_t8\_positioning.py}
and the four T12 checks below.

\section{T9: inverse powers and the smallest eigenvalue}
All matrices in this section have positive dimension. For hard
grounding, budgeted sets leave at least one ungrounded vertex.

\begin{theorem}[Uniform inverse-power spectral rate]\label{thm:spectralrate}
Let $X\succ0$ have dimension $d$ and ordered eigenvalues
$\lambda_1\leq\cdots\leq\lambda_d$. For every $s>0$,
\begin{equation}\label{eq:spectralrate}
 \lambda_1^{-1}\leq(\tr X^{-s})^{1/s}
          \leq d^{1/s}\lambda_1^{-1}.
\end{equation}
The relative excess is at most $e^{(\log d)/s}-1$; it is at most
$\eta$ whenever $s\geq\log d/\log(1+\eta)$.
If $\lambda_1$ has multiplicity $r<d$ and
$z=\lambda_1/\lambda_{r+1}<1$, then
\[
 0\leq\log\bigl(\lambda_1(\tr X^{-s})^{1/s}\bigr)
           -\frac{\log r}{s}
 \leq\frac{d-r}{rs}z^s.
\]
At a fixed cardinality and matrix dimension $d$, any set whose
power-trace objective is within a multiplicative factor $\alpha\geq1$
of its minimum satisfies
\begin{equation}\label{eq:alphaspectral}
 \lambda_{\min}(X_{\rm out})
     \geq(\alpha d)^{-1/s}\lambda_k^*,
\end{equation}
where $\lambda_k^*$ is the maximum smallest eigenvalue at that
cardinality and the same seed. In particular an exact power-trace
minimizer has factor $d^{-1/s}$.
\end{theorem}
\begin{proof}
Factor the smallest eigenvalue out of the trace:
\[
 (\tr X^{-s})^{1/s}
 =\lambda_1^{-1}
    \left(\sum_{i=1}^d(\lambda_1/\lambda_i)^s\right)^{1/s}.
\]
The sum lies between 1 and $d$, giving \eqref{eq:spectralrate}
and its stated rate. With multiplicity $r$ it lies between $r$
and $r+(d-r)z^s$. Taking logarithms and using
$\log(1+x)\leq x$ proves the refined inequality. If $r=d$, the
sum equals $d$ exactly.

Let $O_\lambda$ maximize the smallest eigenvalue and $O_s$ minimize
the trace, at the same budget. Then
\[
 \lambda_{\min}(X_{\rm out})^{-s}
 \leq\Phi_s({\rm out})\leq\alpha\Phi_s(O_s)
 \leq\alpha\Phi_s(O_\lambda)\leq\alpha d(\lambda_k^*)^{-s}.
\]
Taking the inverse $s$th root proves \eqref{eq:alphaspectral}.
The scalar root is used to compare minimizers and spectral values;
the greedy gain remains defined using $\Phi_s=\tr X^{-s}$.
\end{proof}
\noindent Numerical check: \texttt{scripts/verify\_t9\_spectral\_limit.py}.

\begin{theorem}[Baseline-aware translation of the gain guarantee]\label{thm:gaintranslation}
Fix $s>0$, a seed, and $k$ additional anchors. Put
$B_s=\Phi_s(S_0)$, $\Psi_s^*=\min_{|U|=k}\Phi_s(S_0\cup U)$,
and $C_s=B_s/\Psi_s^*$. Let $d$ be the common final dimension
and $\lambda_k^*$ the optimal smallest eigenvalue.
If a returned set satisfies a gain factor $\gamma\in(0,1]$, put
$\rho=1-\gamma$. Then
\begin{align}
 \frac{\lambda_{\min}(X_{\rm out})}{\lambda_k^*}
 &\geq\left[\rho B_s(\lambda_k^*)^s+\gamma d\right]^{-1/s}
 \label{eq:gaintranslation}\\
 &\geq\left[d(\rho C_s+\gamma)\right]^{-1/s}.\notag
\end{align}
Thus a contrast bound $C_s\leq C$ at the chosen exponent and
\[
 s\geq\frac{\log(d(\rho C+\gamma))}{\log(1+\eta)}
\]
certify a factor $1/(1+\eta)$ on the smallest eigenvalue.
For CELF take $\gamma$ from \eqref{eq:lazyguarantee}; the statement
holds with probability at least $1-\delta$.
Without a baseline-contrast condition, any fixed gain factor
$\gamma<1$ alone admits arbitrarily small spectral ratios, even on
a connected weighted four-vertex star and at every integer exponent
$s\geq1$.
\end{theorem}
\begin{proof}
The gain guarantee is exactly the residual-objective inequality
\[
 \Phi_s({\rm out})\leq\rho B_s+\gamma\Psi_s^*.
\]
Since $\Psi_s^*\leq d(\lambda_k^*)^{-s}$ and
$\lambda_{\min}(X_{\rm out})^{-s}\leq\Phi_s({\rm out})$,
the first spectral bound follows. Also
$B_s(\lambda_k^*)^s=C_s\Psi_s^*(\lambda_k^*)^s\leq C_s d$,
giving the second bound and the stated sufficient exponent.
Only the CELF initial event is random.

For the final assertion, take a star with center 0, leaves 1,2,3,
and respective edge weights $b/2,b,1$, where $0<b<1$. The fixed
seed is $S_0=\{0\}$, so the current matrix is
$\diag(b/2,b,1)$. Two additional anchors are allowed.
The optimal choice $\{1,2\}$ leaves the eigenvalue 1; the choice
$\{1,3\}$ leaves the eigenvalue $b$.
Their gain ratio is exactly
\[
 \frac{(2/b)^s+1}{(2/b)^s+b^{-s}}
 =\frac{2^s+b^s}{2^s+1}\geq\frac23>1-e^{-1}
 \qquad(s\geq1).
\]
Nevertheless the spectral ratio is $b$, which tends to zero.
For a larger prescribed $\gamma<1$, replace the weight $b/2$ by
$b/M$ with integer $M\geq\max\{2,\gamma/(1-\gamma)\}$.
The gain ratio is then $(M^s+b^s)/(M^s+1)\geq\gamma$ for
every integer $s\geq1$, and the spectral ratio is still $b$.
This is a counterexample to an implication from the gain factor
alone: it distinguishes that factor from the multiplicative
trace condition in \eqref{eq:alphaspectral}.
For fixed $\rho>0$ and $\lambda_k^*>\lambda_{\min}(X_{S_0})$,
the first bound in \eqref{eq:gaintranslation} tends to
$\lambda_{\min}(X_{S_0})/\lambda_k^*$ as $s\to\infty$, by
Theorem~\ref{thm:spectralrate}. Increasing the exponent makes the
surrogate sharper while the baseline term remains explicit.
\end{proof}
\noindent Numerical check: \texttt{scripts/verify\_t9\_gain\_translation.py}.

\begin{corollary}[A spectral bicriteria guarantee]\label{cor:spectralbudget}
Compare with the optimal set of $k$ additional anchors, with
remaining dimension $d_k$. Let the initial matrix have dimension
$d_0$, let $a>0$ lower-bound its smallest eigenvalue, and let
$B$ upper-bound $\lambda_k^*$. Run the same certified lazy rule
with step factor $1-\theta$ for $\ell\geq k$ rounds.
Whenever the budget is feasible and hard grounding leaves a vertex,
\begin{equation}\label{eq:budgetresidual}
 \Phi_s(S_\ell)
 \leq e^{-(1-\theta)\ell/k}B_s+\Psi_{s,k}^*.
\end{equation}
For a target $\eta>0$, choose an integer
\begin{align}
 s&\geq\left\lceil\frac{\log(2d_k)}{\log(1+\eta)}\right\rceil,
 \label{eq:budgetchoices}\\
 \ell&\geq\max\left\{k,
 \left\lceil\frac{k}{1-\theta}
       \left(\log\frac{d_0}{d_k}+s\log\frac Ba\right)\right\rceil
              \right\}.\notag
\end{align}
Then, on the same probability-$1-\delta$ event,
$\lambda_{\min}(X_{S_\ell})\geq\lambda_k^*/(1+\eta)$.
\end{corollary}
\begin{proof}
The maximum-marginal argument in Theorem~\ref{thm:celf} may
compare every round with a fixed size-$k$ optimum, even after
more than $k$ rounds. Its residual recurrence is
\[
 \Phi_s(S_{j+1})-\Psi_{s,k}^*
 \leq\left(1-\frac{1-\theta}{k}\right)
               (\Phi_s(S_j)-\Psi_{s,k}^*).
\]
If the parenthesized gap is negative, monotonicity also implies
this inequality. Iteration and the exponential bound prove
\eqref{eq:budgetresidual}. Now $B_s\leq d_0a^{-s}$ and
$\Psi_{s,k}^*\leq d_k(\lambda_k^*)^{-s}$.
The second choice in \eqref{eq:budgetchoices} bounds the first
term in \eqref{eq:budgetresidual} by
$d_k B^{-s}\leq d_k(\lambda_k^*)^{-s}$. Therefore
\[
 \lambda_{\min}(X_{S_\ell})^{-s}
 \leq\Phi_s(S_\ell)\leq2d_k(\lambda_k^*)^{-s}.
\]
The first choice of exponent proves the spectral factor.
The permanent-key argument remains valid for all $\ell$ rounds
without additional adaptive-sketch events. For growing integer
$s$, the coarse endpoint-moment initialization avoids the
$2^s$ precision factor of the fine logarithmic-polynomial estimate;
the power-chain precision term in Theorem~\ref{thm:lazycost}
continues to be counted explicitly.
\end{proof}
\noindent Numerical check: \texttt{scripts/verify\_t9\_bicriteria.py}.

\paragraph{Relation to the grounded-eigenvalue objective.}
Zhou et al.~\cite{zhou2023} establish that the smallest grounded
eigenvalue is monotone but not submodular. The inverse-power family
provides supermodular trace objectives at every finite exponent,
together with a uniform approximation to the reciprocal smallest
eigenvalue after the monotone scalar root. Equations
\eqref{eq:gaintranslation} and \eqref{eq:budgetchoices} quantify,
respectively, the fixed-budget baseline-aware guarantee and the
resource-augmented spectral guarantee furnished by this route.

\section{T10: a fixed-budget spectral approximation barrier}
We quantify the cubic-cover reduction of Zhou et al.~\cite{zhou2023}
with a $3$ versus $2$ gap and extend it to a prescribed seed.
For a maximization problem, a $\gamma$-approximation means a
feasible set of exactly the given budget with value at least
$\gamma$ times optimum, where $0<\gamma\leq1$.

\begin{theorem}[Cubic spectral gap and its seeded extension]
\label{thm:spectralhardness}
Unless $\mathrm P=\mathrm{NP}$, for every fixed $\gamma>2/3$
there is no polynomial-time $\gamma$-approximation to
\[
 \max_{|S|=k}\lambda_{\min}(L_{-S})
\]
on connected unweighted cubic graphs, with $1\leq k<n$ supplied
as input. The same barrier holds on connected positively weighted
graphs with a prescribed singleton seed, $k$ additional anchors,
and any fixed rational $\sigma\geq0$ for
$\lambda_{\min}((L+\sigma I)_{-(S_0\cup U)})$.
For randomized algorithms succeeding with probability at least
$2/3$ on every input, the corresponding consequence is
$\mathrm{NP}=\mathrm{RP}$.
\end{theorem}
\begin{proof}
Use the NP-complete vertex-cover problem on connected cubic
graphs; it remains NP-complete even with planarity and
3-connectivity~\cite{uehara1996}. A cover of size at most $k$ can
be enlarged to size exactly $k<n$.
For every nonempty $S\subsetneq V$, the grounded matrix is
\[
 L_{-S}=3I-B_U,\qquad U=V\setminus S,
\]
where $B_U$ is the induced adjacency matrix. If $S$ is a vertex
cover then $B_U=0$ and its smallest eigenvalue is exactly 3.
If $S$ is not a cover, take an uncovered edge $uv$ and the unit
vector $x=(e_u+e_v)/\sqrt2$ on $U$. Direct evaluation gives
$x^\top L_{-S}x=2$, so its smallest eigenvalue is at most 2.
Every grounded matrix also has smallest eigenvalue at most 3,
by a coordinate Rayleigh quotient. Thus the optimum is 3 on
YES instances and at most 2 on NO instances.

A $\gamma>2/3$ approximation on a YES instance returns value
strictly larger than 2, so it must return a vertex cover. Check
this combinatorial property of the returned set. On a NO instance
no returned size-$k$ set can pass the check. This decides the
NP-complete problem without computing or comparing eigenvalues.
It is a polynomial many-one gap reduction with an exact integral
YES/NO separation. The budget is part of the input; the argument
applies to algorithms obeying that budget exactly.

For the prescribed-seed version, take the same cubic graph $H$,
scale each original edge weight to a positive integer $W$, and add
a new seed vertex $z$ joined to every original vertex with weight
one. Ground $z$ in advance and choose $k$ additional original
vertices. For $U$ the remaining original vertices the matrix is
\begin{equation}\label{eq:seedhardgap}
 (L+\sigma I)[U,U]=(3W+1+\sigma)I-W B_U.
\end{equation}
Let $D=3W+1+\sigma$. A cover gives eigenvalue $D$; an uncovered
edge gives Rayleigh quotient $D-W$. Choose the constant integer
$W$ large enough that
\[
 (3\gamma-2)W>(1-\gamma)(1+\sigma).
\]
This is exactly $\gamma D>D-W$ and is possible for every fixed
$\gamma>2/3$. For fixed $\gamma,\sigma$, the integer $W$ is a
constant independent of the source instance, so construction and
encoding are polynomial. The graph is connected and all its
weights are positive; the mandatory seed makes this a direct
many-one reduction for the same seeded formulation used in T9.
Again the output is checked only for vertex coverage in $H$.

For a randomized approximation, the cover check accepts every
YES instance with probability at least $2/3$ and never accepts a
NO instance. It therefore places this NP-complete decision problem
in $\mathrm{RP}$. Closure under polynomial reductions and
$\mathrm{RP}\subseteq\mathrm{NP}$ yield
$\mathrm{NP}=\mathrm{RP}$.
\end{proof}
\noindent Numerical check: \texttt{scripts/verify\_t10\_hardness.py}.

\subsection{Budget-preserving amplification by parallel paths}
The following construction leaves the anchor budget unchanged and
allows every new internal vertex to be selected. For an integer
$\ell\geq2$, define
\begin{equation}\label{eq:pathconstants}
 a_\ell=2-2\cos\frac{\pi}{2\ell-1},\qquad
 b_\ell=2-2\cos\theta_\ell,
\end{equation}
where $\theta_\ell$ is the unique root in
$(0,\pi/(2\ell-1))$ of
\begin{equation}\label{eq:pathroot}
 3\cos((\ell-\tfrac12)\theta)=\cos(\theta/2).
\end{equation}

\begin{lemma}[Parallel-path YES/NO bounds]\label{lem:pathamplification}
Let $H$ be a connected simple cubic graph. Form $G_{M,\ell}$ by
replacing every edge of $H$ by $M$ internally vertex-disjoint
paths of length $\ell$, with unit edge weights. The original
vertices remain as the paths' endpoints. Then:
\begin{enumerate}
\item If $S$ is an original vertex cover with $|S|<|V(H)|$, then
\begin{equation}\label{eq:pathyes}
 \lambda_{\min}(L(G_{M,\ell})_{-S})
 \geq\frac{a_\ell}{1+2(1+a_\ell)/(3M)}.
\end{equation}
\item Let $T$ be any set of $k$ vertices of $G_{M,\ell}$, including
arbitrary internal vertices. If $T\cap V(H)$ is not a cover and
$M(\ell-1)\geq8k$, then
\begin{equation}\label{eq:pathno}
 \lambda_{\min}(L(G_{M,\ell})_{-T})
 \leq b_\ell+\frac{24k}{M(\ell-1)}.
\end{equation}
\item The constants satisfy $0<b_\ell<a_\ell\leq1$,
$b_\ell/a_\ell>c_*$ for every finite $\ell$, and
\begin{equation}\label{eq:pathlimit}
 \lim_{\ell\to\infty}\frac{b_\ell}{a_\ell}
 =c_*:=\left(\frac{2\arccos(1/3)}{\pi}\right)^2.
\end{equation}
\end{enumerate}
\end{lemma}
\begin{proof}
Put $h=\ell-1$. The quadratic form
\[
 q(y)=\sum_{j=1}^{h-1}(y_j-y_{j+1})^2+y_h^2
\]
has smallest eigenvalue $a_\ell$. Indeed its matrix is the path
tridiagonal with first diagonal entry 1, all other diagonal entries
2, and off-diagonals $-1$ (for $h=1$ the matrix is $[1]$).
Its eigenvectors have entries $\cos((j-\tfrac12)\vartheta)$,
with $\vartheta=(2r-1)\pi/(2h+1)$, $r=1,\ldots,h$;
substitution gives the distinct eigenvalues $2-2\cos\vartheta$.
Thus $q(y)\geq a_\ell\|y\|^2$.

When $S$ is a cover, each component after grounding is either a
star with $3M$ length-$\ell$ rays ending at ground, or an internal
path with both endpoints grounded. For a star let $x$ be its center
value, $y^{(r)}$ the internal values of ray $r$, and write
$Y=\sum_r\|y^{(r)}\|^2$. Its energy is
\[
 E=\sum_{r=1}^{3M}\bigl((x-y^{(r)}_1)^2+q(y^{(r)})\bigr)
 \geq a_\ell Y.
\]
Averaging $x^2\leq2(x-y^{(r)}_1)^2+2(y^{(r)}_1)^2$ over the
rays yields $x^2\leq2(E+Y)/(3M)$. Hence
\[
 x^2+Y\leq\left(\frac{1+2/(3M)}{a_\ell}+
                    \frac2{3M}\right)E.
\]
The Rayleigh principle gives \eqref{eq:pathyes}. A path grounded
at both ends has energy $q(y)+y_1^2\geq a_\ell\|y\|^2$,
so the same lower bound holds for the whole matrix.

For the NO bound, choose an original uncovered edge $uv$.
Give $u,v$ value 1 and every other original vertex value 0.
On every length-$\ell$ path with endpoint values $x,y\in\{0,1\}$,
give its internal vertex $j$ the value
\begin{equation}\label{eq:pathtrial}
 z_j=\frac{x\sin((\ell-j)\theta_\ell)
               +y\sin(j\theta_\ell)}{\sin(\ell\theta_\ell)}.
\end{equation}
These values satisfy
$2z_j-z_{j-1}-z_{j+1}=b_\ell z_j$ at every internal vertex.
At each of $u,v$, there is one original path type with endpoint
values $(1,1)$ and two types with endpoint values $(1,0)$.
The total outward flux, divided by $M$, is
\[
 3-\frac{3\sin((\ell-1)\theta_\ell)+\sin\theta_\ell}
             {\sin(\ell\theta_\ell)}=0,
\]
by \eqref{eq:pathroot} and the sine subtraction formula.
All other original vertices have value zero. If $D$ is the sum of
squared internal values across one copy of each original path type,
then the full trial vector has squared norm $MD+2$ and energy
exactly $b_\ell MD$.

All trial values lie in $[0,2]$. To see this, note that
$\ell\theta_\ell/2<\ell\pi/(4\ell-2)\leq\pi/3$ and
$\ell\theta_\ell<\pi$. The $(1,1)$ path values equal
\[
 \frac{\cos((j-\ell/2)\theta_\ell)}
      {\cos(\ell\theta_\ell/2)}\in[1,2],
\]
and the nonnegative $(1,0)$ and $(0,1)$ values are at most their
sum. In particular the $uv$ path alone gives $D\geq h$.

Now set the value of every internal vertex in $T$ to zero. The
original values already vanish on $T\cap V(H)$. Each altered
internal vertex has degree 2, so at most $2k$ edges change energy.
A changed edge has new energy at most 4, and the squared norm
falls by at most $4k$. The modified vector is nonzero because
$M h\geq8k$ and it is an admissible Dirichlet test vector. Its
Rayleigh quotient is at most
\[
 \frac{b_\ell MD+8k}{MD+2-4k}
 \leq b_\ell+\frac{(8+4b_\ell)k}{MD-4k}
 \leq b_\ell+\frac{24k}{Mh},
\]
using $b_\ell<1$ and $MD\geq Mh\geq8k$.
This handles internal anchors without restricting the candidate set.

For completeness, the left side minus the right side of
\eqref{eq:pathroot} equals 2 at zero and is negative at
$\pi/(2\ell-1)$. Its derivative is negative in between, since
$\sin((\ell-\tfrac12)\theta)\geq\sin(\theta/2)>0$ there and
$3(\ell-\tfrac12)>1/2$. This proves existence, uniqueness, and
$0<b_\ell<a_\ell\leq1$. Put
$t_\ell=(\ell-\tfrac12)\theta_\ell\in(0,\pi/2)$.
Equation~\eqref{eq:pathroot} becomes
$3\cos t_\ell=\cos(t_\ell/(2\ell-1))$.
Its right side is strictly less than 1 and converges uniformly to 1,
so $t_\ell>\arccos(1/3)$ and
$t_\ell\to\arccos(1/3)$. Concavity of sine on $[0,\pi/2]$ gives
\[
 \frac{b_\ell}{a_\ell}
 =\left(\frac{\sin(t_\ell/(2\ell-1))}
 {\sin(\pi/(2(2\ell-1)))}\right)^2
 \geq\left(\frac{2t_\ell}{\pi}\right)^2>c_*.
\]
Using $2-2\cos x=4\sin^2(x/2)$
and $\sin x/x\to1$ proves \eqref{eq:pathlimit}.
\end{proof}
\noindent Numerical check: \texttt{scripts/verify\_t10\_paths.py}.

\begin{theorem}[Amplified exact-budget spectral hardness]
\label{thm:amplifiedhardness}
For every fixed $\gamma>c_*$, where $c_*$ is defined in
\eqref{eq:pathlimit}, a polynomial-time $\gamma$-approximation
for exact-budget grounded-$\lambda_{\min}$ maximization on connected
simple unweighted graphs would imply $\mathrm P=\mathrm{NP}$.
The reduction preserves the source budget $k$ exactly and permits
all graph vertices as candidates. In particular this excludes a
same-budget spectral factor $1-1/e$. It also extends to a prescribed
singleton seed and every fixed rational $\sigma\geq0$ on connected
positively weighted graphs. A randomized approximation with success
probability at least $2/3$ would imply $\mathrm{NP}=\mathrm{RP}$.
\end{theorem}
\begin{proof}
It suffices to consider $\gamma\leq1$. Choose a constant integer
$\ell\geq2$ with $b_\ell/a_\ell<\gamma$, possible by
\eqref{eq:pathlimit}, and put $\Delta=\gamma a_\ell-b_\ell>0$.
Fix an integer constant $C>\max\{8,32/\Delta\}$.
Given a connected cubic vertex-cover instance $(H,k)$,
$1\leq k<|V(H)|$, construct $G_{Ck,\ell}$ and keep budget $k$.
It has $n+Ckm(\ell-1)$ vertices and $Ckm\ell$ edges, so its size
is polynomial. Every vertex, including every path-interior vertex,
is eligible for grounding. Define the constants
\[
 \alpha=\frac{a_\ell}{1+4/(3C)},\qquad
 \beta=b_\ell+\frac{24}{C(\ell-1)}.
\]
Because $a_\ell\leq1$,
\[
 \gamma\alpha-\beta
 \geq\Delta-\frac{4\gamma a_\ell/3+24/(\ell-1)}C
 >\Delta-\frac{32}C>0.
\]
If $H$ has a cover of size at most $k$, enlarge it to size $k$
and ground those original vertices. Equation~\eqref{eq:pathyes}
gives value at least $\alpha$. If an arbitrary returned size-$k$
set $T$ has original part $T\cap V(H)$ that is not a cover,
\eqref{eq:pathno} gives value at most $\beta$.
Thus a $\gamma$-approximation on a YES instance must return an
original part that is a cover, since its value is at least
$\gamma\alpha>\beta$. Check that original part combinatorially.
On a NO instance no original part of size at most $k$ can pass.
This decides cubic vertex cover and proves the claimed hardness.
The constants depend only on the fixed target factor; strict
inequalities allow fixed rational upper/lower enclosures when
choosing them, independently of the input instance.

The special factor $1-1/e$ has an exact rational certificate.
For $\ell=4$, the ND path matrix and the sine recurrence give
$a_4$ as the smallest positive zero of $F$, and $3F(b_4)=1$, where
$F(x)=1-6x+5x^2-x^3$ is strictly decreasing on $[0,1/5]$.
Direct rational evaluations give
$F(99/500)>0$, $F(1/5)<0$, and $3F(309/2500)<1$.
Consequently
\[
 \frac{b_4}{a_4}<\frac{309/2500}{99/500}
 =\frac{103}{165}<\frac58<1-\frac1e.
\]
The last inequality follows from
$e>1+1+1/2+1/6+1/24=65/24>8/3$.

For a fixed seed and shift, multiply every path-edge weight by
a constant positive integer $W$, add a new mandatory seed $z$,
and join $z$ to every other vertex by a unit-weight edge.
After grounding $z$, any additional set $T$ has matrix
$W L(G_{Ck,\ell})_{-T}+(1+\sigma)I$.
Choose $W$ so that
\[
 W(\gamma\alpha-\beta)>(1-\gamma)(1+\sigma).
\]
The YES lower bound $W\alpha+1+\sigma$ and the NO upper bound
$W\beta+1+\sigma$ are then separated by factor $\gamma$.
The graph and its rational encoding remain polynomial, and the
budget still counts exactly $k$ additions to the mandatory seed.
The same combinatorial output check proves the reduction.
For randomized algorithms this check has no false positives and
succeeds with probability at least $2/3$ on YES instances, giving
the $\mathrm{NP}=\mathrm{RP}$ consequence as in
Theorem~\ref{thm:spectralhardness}.
\end{proof}
\noindent Numerical check: \texttt{scripts/verify\_t10\_amplified\_hardness.py}.

\subsection{Structural constraints on further amplification}
\begin{proposition}[The cubic half-budget threshold is sharp]
\label{prop:cubicalgorithm}
On connected unweighted cubic graphs and budgets $n/2\leq k<n$,
a polynomial-time $2/3$-approximation exists. Together with
Theorem~\ref{thm:spectralhardness}, this is the sharp approximation
threshold in this graph/budget regime unless $\mathrm P=\mathrm{NP}$.
\end{proposition}
\begin{proof}
Start with any two-coloring. Whenever a vertex has at least two
neighbors of its own color, flip it. Each flip increases the number
of cut edges by at least one; at most $m$ flips occur. The final
color classes each induce maximum degree at most one. One class
has at least $n/2$ vertices. Keep any $n-k$ vertices of that class
and ground all others. The remaining adjacency matrix is a
matching plus isolated vertices, so its largest eigenvalue is at
most one and the grounded Laplacian has smallest eigenvalue at
least two. Every feasible smallest eigenvalue is at most three.
Scanning for a flippable vertex after each flip is already a
polynomial-time implementation.

A cubic vertex cover has size at least $m/3=n/2$. Hence the
hardness reduction may restrict source instances to $k\geq n/2$:
instances with $k<n/2$ are immediately NO and can be replaced by
the fixed NO instance $(K_4,2)$. Theorem~\ref{thm:spectralhardness}
therefore excludes every larger factor in the stated regime.
\end{proof}
\noindent Numerical check: \texttt{scripts/verify\_t10\_cubic\_threshold.py}.

\begin{proposition}[Restrictions on direct matrix transformations]
\label{prop:amplificationtransforms}
The following statements hold.
\begin{enumerate}
\item If SPD SDDM matrices $A,B$ each have a strictly negative
 off-diagonal entry, $A\otimes B$ has a positive off-diagonal entry.
 Every Schur complement of an SPD SDDM matrix is SDDM.
 The Kronecker sum $A\otimes I+I\otimes B$ is SDDM and has
 smallest eigenvalue $\lambda_{\min}(A)+\lambda_{\min}(B)$.
\item Let $f$ be continuous on $(0,\infty)$, with $f(A)$ SDDM
 for every SPD SDDM $A$. If
 $f(A)[U,U]=f(A[U,U])$ for every principal index set $U$, then
 $f(x)=cx+d$ with $c,d\geq0$. For $0<a<b$ and $f(b)>0$,
 $f(a)/f(b)\geq a/b$.
\item The map $\mathcal H(A)=((A^{-1})^{\circ2})^{-1}$ fails
 to commute with grounding even on the explicit SDDM pair
 \[
 A=\begin{pmatrix}2&-1&0\\-1&2&-1\\0&-1&2\end{pmatrix},\qquad
 \mathcal H(A)=\begin{pmatrix}2&-1/2&0\\-1/2&5/4&-1/2\\0&-1/2&2\end{pmatrix}.
 \]
 For $U=\{1,3\}$, $\mathcal H(A)[U,U]=2I$ whereas
 $\mathcal H(A[U,U])=4I$.
\end{enumerate}
\end{proposition}
\begin{proof}
If $A_{ij}<0$ and $B_{uv}<0$ with $i\ne j,u\ne v$, the
$((i,u),(j,v))$ entry of their tensor product is positive.
For the Schur assertion write an SPD SDDM matrix as
$\left(\begin{smallmatrix}P&Q\\Q^\top&D\end{smallmatrix}\right)$.
Its inverse block $D^{-1}$ is entrywise nonnegative: write
$D=sI-N$ with $N\geq0$ and $s>\rho(N)$ and use the convergent
Neumann series. Thus $P-QD^{-1}Q^\top$ has nonpositive
off-diagonals. If $r_1=P\ind+Q\ind\geq0$ and
$r_2=Q^\top\ind+D\ind\geq0$, its row sums are
$r_1-QD^{-1}r_2\geq0$. Positive definiteness follows by minimizing
the original quadratic form over the eliminated variables.
For example, $A^2$ for the displayed three-by-three $A$ also has
positive $(1,3)$ entry 1. These sign tests concern exact matrices
and exact static Schur complements. The Kronecker sum has
nonpositive off-diagonals and nonnegative row sums. Tensor products
of the two eigenvector bases diagonalize it with eigenvalues
$\lambda_i(A)+\lambda_j(B)$, proving its asserted smallest value.

For the second assertion take
$A=\left(\begin{smallmatrix}t&-u\\-u&t\end{smallmatrix}\right)$,
$t>u\geq0$, and restrict to a single index. Its eigenvalues are
$t-u,t+u$, so the commutation hypothesis gives
$f(t)=(f(t-u)+f(t+u))/2$. Every pair of positive numbers has this
form. Repeated midpoint subdivision gives affine interpolation
at dyadic rational points of every interval; continuity extends
it to all points, proving that $f(x)=cx+d$ with common constants.
The off-diagonal sign of $f(A)=cA+dI$ forces $c\geq0$; its
nonnegative row sums for arbitrarily small positive scalar $A$
force $d\geq0$. Finally
$b(ca+d)-a(cb+d)=d(b-a)\geq0$ gives the ratio inequality.

For the third assertion direct multiplication gives
\[
 A^{-1}=\frac14\begin{pmatrix}3&2&1\\2&4&2\\1&2&3\end{pmatrix}.
\]
Squaring entries and multiplying by the displayed $\mathcal H(A)$
gives the identity matrix. Both displayed matrices have
nonpositive off-diagonals and nonnegative row sums, and are
positive definite. The two principal-transform values now follow
by restricting to indices 1 and 3 and inverting scalar diagonals.
\end{proof}
\noindent Numerical check: \texttt{scripts/verify\_t10\_transforms.py}.

\begin{proposition}[Uniform pendant ratio bound]
\label{prop:pendantamplification}
Attach the same connected weighted rooted graph $H$ to every
vertex of a base graph, identifying its root with that vertex.
For any nonempty original grounding set $S$ leaving an original
vertex, keep every auxiliary vertex ungrounded. If the original
smallest grounded eigenvalue is $\mu$, the decorated graph's
smallest grounded eigenvalue is $g(\mu)$, where $g$ is increasing
in the non-strict sense, concave, positive on $(0,\infty)$, and
$g(0)=0$. In particular $g(a)/g(b)\geq a/b$ for $0<a<b$.
The selected set and its cardinality are unchanged in this identity.
\end{proposition}
\begin{proof}
Write the Laplacian of the attached graph, with root first, as
\[
 L_H=\begin{pmatrix}d&-b^\top\\-b&D\end{pmatrix},\quad
 b\geq0,\quad D\ind=b,\quad D\succ0.
\]
If there are no auxiliary vertices the assertion is $g(\mu)=\mu$.
Otherwise define
\[
 g(\mu)=\lambda_{\min}
 \begin{pmatrix}d+\mu&-b^\top\\-b&D\end{pmatrix}.
\]
The Rayleigh minimum of affine nondecreasing functions of $\mu$
is concave and nondecreasing. Connectivity gives $g(0)=0$ and
$g(\mu)>0$ for $\mu>0$ by the edge-energy form plus a positive
root penalty.

Diagonalize the base grounded Laplacian and apply the same
orthogonal change of coordinates to its attached copies. The
part belonging to ungrounded original roots splits into the
above blocks, one per base eigenvalue. The decorations of
selected roots give additional isolated $D$ blocks.
Let $\lambda_D=\lambda_{\min}(D)$. A unit minimizing eigenvector
$w$ can be chosen nonnegative, because taking entrywise absolute
values never increases the quadratic form of an SDDM matrix.
Then $b^\top w=(D\ind)^\top w=\lambda_D\ind^\top w>0$.
Testing the block defining $g(\mu)$ on $(t,w)$ for small $t>0$
gives a Rayleigh quotient strictly below $\lambda_D$, since
its numerator is $\lambda_D-2t b^\top w+(d+\mu)t^2$.
Thus the isolated $D$ blocks are above $g(\mu)$, and monotonicity
shows that the full smallest eigenvalue is $g$ of the smallest
base eigenvalue. Finally concavity at
$a=(a/b)b+(1-a/b)0$ gives $g(a)\geq(a/b)g(b)$.
\end{proof}
\noindent Numerical check: \texttt{scripts/verify\_t10\_pendant.py}.

\paragraph{Consequence for the T9 bicriteria guarantee.}
For $0<\eta<c_*^{-1}-1$, the spectral factor $1/(1+\eta)$ is
strictly larger than $c_*$. Theorem~\ref{thm:amplifiedhardness} rules out a
polynomial-time same-budget guarantee with that factor under the
stated complexity assumptions. T9's augmented budget provides a
resource-relaxed route to precisely such high-accuracy factors,
subject to its feasibility condition. The numerical pilot in
\texttt{notes/findings.md}, item 6, motivates comparing finite
powers on particular graphs; the hardness statement concerns
worst-case instances and the exact prescribed budget.

\section{Response to the Zhou comparison: finite removal and oracle precision}

\begin{proposition}[A finite-removal stress case]\label{prop:zhoustress}
Let $0<t<1/4$. On vertices $\{z,a,b,c\}$ put edges $za,ab$
of weight $t$ and $zb,bc,cz$ of weight one. Ground the prescribed
seed $z$ and choose one further vertex. Apply the weighted version
of the score in Zhou et al.\ \cite{zhou2023}, Eq.~(4),
$h_v=2u_v\sum_{w\notin S}w_{vw}u_w$, using the exact smallest
eigenvector. The score uniquely selects $b$, and its achieved
smallest eigenvalue divided by the optimum is at most $2t$.
For every real $p\geq1$, exact minimization of $\tr X^{-p}$
over the three one-vertex choices uniquely selects the spectral
optimum $a$.
\end{proposition}
\begin{proof}
In the order $a,b,c$, the current grounded matrix is
\[
 X=\begin{pmatrix}2t&-t&0\\-t&2+t&-1\\0&-1&2\end{pmatrix}.
\]
It is positive definite by grounded connectivity. Its off-diagonal
graph is connected, so its smallest eigenvector is strictly positive:
apply Perron--Frobenius to $dI-X$ with $d>\lambda_{\max}(X)$.
The scores obey
\[
 h_a=2tu_au_b>0,\quad h_c=2u_bu_c>0,\quad h_b=h_a+h_c.
\]
Consequently $b$ is the unique maximizer, independently of how
accurately the positive eigenvector is scaled. After selecting $b$
the matrix is $\diag(2t,2)$ and its smallest eigenvalue is $2t$.
After selecting $a$, the matrix is
\[
 \begin{pmatrix}2+t&-1\\-1&2\end{pmatrix}
 =\begin{pmatrix}2&-1\\-1&2\end{pmatrix}
   +\diag(t,0)\succeq I,
\]
so its smallest eigenvalue is at least one. After selecting $c$,
the Rayleigh quotient at the coordinate $a$ is $2t<1$.
This proves uniqueness of the spectral optimum and the ratio bound.
The inverse-power trace after selecting $a$ is at most two; after
either alternative it is at least $(2t)^{-p}>2$ for every $p\geq1$.
The inverse-power choice is therefore also uniquely $a$.
\end{proof}
\noindent Numerical check: \texttt{scripts/verify\_zhou\_response.py}.

\begin{proposition}[Tighter endpoint-moment upper bounds]
\label{prop:tightendpoint}
Use the notation of Proposition~\ref{prop:integerpositive} and let
$p\geq2$ be an integer. Put $v_2=a_2-a_1^2\geq0$. Then
\begin{align*}
 q_p&\leq U_2:=p a_p-(p-1)a_1^p-p(p-2)v_2a_1^{p-2}\\
    &\leq U_1:=p a_p-(p-1)a_1^p\leq p a_p.
\end{align*}
Both upper bounds are at least $a_p$. For $p=2$, $U_1=U_2=q_2$;
for $p=3$, $U_2=q_3$. If $R$ has rank one and hard grounding
removes its nonzero mode, both bounds equal $q_p=a_p$.

For certified intervals $a_j\in[\ell_j,u_j]$, with nonnegative
endpoints, set $v_- =\max\{0,\ell_2-u_1^2\}$. A certified
version is
\[
 q_p\leq p u_p-(p-1)\ell_1^p
                   -p(p-2)v_-\ell_1^{p-2}.
\]
At $p=2$ the last term is defined as zero. For $p=1$ use $q_1=a_1$.
\end{proposition}
\begin{proof}
Define the nonnegative word sums
\[
 W_\ell=\sum_{\substack{j_1+\cdots+j_\ell=p\\j_i\geq1}}
                d_{j_1}\cdots d_{j_\ell}.
\]
The proof of Proposition~\ref{prop:integerpositive} gives
$a_p=\sum_{\ell=1}^p W_\ell$ and
$q_p=\sum_{\ell=1}^p(p/\ell)W_\ell$. Therefore
\[
 p a_p-q_p=\sum_{\ell=2}^p p(1-1/\ell)W_\ell.
\]
Keep just the $\ell=p$ term for $U_1$. Since $d_1=a_1$,
$W_p=a_1^p$. For $p\geq3$ also keep the $\ell=p-1$ term.
Here $d_2=cx^\top(R-cxx^\top)x=a_2-a_1^2=v_2$ and
$W_{p-1}=(p-1)v_2a_1^{p-2}$, yielding $U_2$.
All discarded terms are nonnegative. For $p=2,3$ all positive
correction terms have been retained, giving the exact special cases.
The lower bound $q_p\geq a_p$ from the same proposition proves
$U_2,U_1\geq a_p$. In the rank-one hard case $K=0$, so only
$d_1$ is nonzero and only $W_p$ survives. Finally
$v_2\geq v_-$ and $a_1\geq\ell_1$, which proves the interval
bound by bounding each subtracted nonnegative product from below
and $a_p$ from above.
\end{proof}
\noindent Numerical check: \texttt{scripts/verify\_marginal\_upper.py}.

\begin{proposition}[Cache and deflation precision conditions]
\label{prop:responsescope}
If a cached moment $\widehat r^{(0)}$ is advanced by exact,
cache-independent decrements $D_i=r^{(i)}-r^{(i+1)}$, its signed
absolute error is constant. With decrements $\widehat D_i$, its
error after $h$ updates is
\[
 e^{(h)}=e^{(0)}+\sum_{i=0}^{h-1}(D_i-\widehat D_i).
\]
Even with exact decrements, its relative error can grow without bound
on connected grounded graphs.

For a positive definite current matrix $X$, let
$0<\lambda_1\leq\cdots\leq\lambda_d$ and $\phi_i$ be an orthonormal
eigenbasis. Fix $r<d$ and define
\[
 H_j(u)=\sum_{i\leq r}\lambda_i^{-j}\phi_i(u)^2,\quad
 T_j(u)=\sum_{i>r}\lambda_i^{-j}\phi_i(u)^2,\quad
 w_r(u)=\sum_{i\leq r}\phi_i(u)^2.
\]
If $H_j$ is exact and $|\widehat T_j-T_j|\leq\zeta T_j$, then
\[
 \frac{|H_j+\widehat T_j-r_j|}{r_j}
 \leq\zeta\frac{T_j}{r_j}.
\]
For $w_r(u)>0$, the tail fraction is at most
\[
 \min\left\{1,\frac{1-w_r(u)}{w_r(u)}
                  \left(\frac{\lambda_r}{\lambda_{r+1}}\right)^j\right\}.
\]
If $\phi_1(u)\ne0$, $w_r(u),\lambda_r$ in the denominator bound
can instead be replaced by $\phi_1(u)^2,\lambda_1$, retaining
$1-w_r(u)$ in the numerator. A bound using only the latter
spectral ratio and $j$ cannot force the tail fraction to zero
uniformly over all vertices and connected grounded matrices.
\end{proposition}
\begin{proof}
Subtract the exact update from the computed update and telescope
to obtain the error identity. Consider the graph with seed $z$,
edge $zu$ of weight $t$ and edge $uv$ of weight one. Before
grounding $u$, $r_1(v)=(1+t)/t$, and afterwards $r_1(v)=1$.
An initial signed relative error $-\zeta$ becomes final signed
relative error $-\zeta(1+t)/t$ under the exact decrement $1/t$.
It is unbounded as $t\downarrow0$.

The deflated error identity follows by subtracting $H_j+T_j$.
The inequalities $T_j\leq(1-w_r)\lambda_{r+1}^{-j}$ and
$r_j\geq H_j\geq w_r\lambda_r^{-j}$ give the first bound.
Using $r_j\geq\phi_1(u)^2\lambda_1^{-j}$ gives the second.
To check the uniform obstruction even with a connected follower
graph, fix any positive integer $j$, let $0<t<1$, set $b=t^{j+2}$,
and consider
\[
 X_t=\begin{pmatrix}t+b&-b\\-b&1+b\end{pmatrix}.
\]
This is the grounded matrix of seed edges of weights $t,1$ and
one follower edge of weight $b$. As $t\downarrow0$, its eigenvalues
are $t(1+o(1))$ and $1+o(1)$. The second coordinate of its normalized
smallest eigenvector is $O(b)$, by its second eigenvector equation.
At this coordinate, for $r=1$, $H_j=O(b^2t^{-j})=O(t^{j+4})$
while $T_j=1+o(1)$. Thus $T_j/r_j\to1$ even though
$(\lambda_1/\lambda_2)^j\to0$.
\end{proof}
\noindent Numerical check: \texttt{scripts/verify\_response\_precision.py}.

\section{T11: preconditioning, cache retention, and scores}

\subsection{A fixed seed preconditioner}
Write $W=V\setminus S_0$, $N=|W|$, $A_0=X_{S_0}\succ0$,
$T=S\setminus S_0$, $t=|T|$, and $U=W\setminus T$.
All inverses in this subsection act on their active coordinates.
An inverse preconditioner $B_0\succ0$ for $A_0$ is applied to a
vector on $U$ by zero extension to $W$, application of $B_0$, and
restriction back to $U$. Its active matrix is $B_U=B_0[U,U]$.

\begin{theorem}[Masked seed preconditioning]\label{thm:seedpcg}
Suppose $aA_0^{-1}\preceq B_0\preceq bA_0^{-1}$ for $0<a\leq b$.
Let $A=A_0[U,U]$ and $d=|U|>0$. Every eigenvalue of
$H=A^{1/2}B_UA^{1/2}$ is at least $a$, and at most $t$ of its
eigenvalues exceed $b$. Let $q\leq\min\{t,d\}$ count those
exceeding $b$, and set $\chi=b/a$. In exact arithmetic PCG obeys,
for $h\geq q$ and $\chi>1$,
\[
 \frac{\|e_h\|_A}{\|e_0\|_A}
 \leq 2\left(\frac{\sqrt\chi-1}{\sqrt\chi+1}\right)^{h-q}.
\]
It therefore reaches relative energy error $\xi$ in at most
\[
 \min\left\{d,\ t+
 \left\lceil\frac{\log(2/\xi)}{
 \log((\sqrt\chi+1)/(\sqrt\chi-1))}\right\rceil\right\}
\]
iterations. For $a=b$, termination occurs within $\min\{d,t+1\}$
iterations. These bounds hold simultaneously for every $T$ on the
single seed spectral-approximation event.

The identity-padded matrix $\overline A=A\oplus I_T$ differs from
$A_0$ by rank at most $2t$.
\end{theorem}
\begin{proof}
Order $U$ before $T$ and write
$A_0=\left(\begin{smallmatrix}A&C\\C^\top&D\end{smallmatrix}\right)$.
The block inverse formula gives
\[
 (A_0^{-1})[U,U]
 =A^{-1}+A^{-1}C(D-C^\top A^{-1}C)^{-1}C^\top A^{-1}.
\]
The Schur complement is positive definite. Congruence by $A^{1/2}$
thus yields $K=I+ZZ^\top$ with $\operatorname{rank}Z\leq t$.
Compression preserves the assumed Loewner inequalities, so
$aK\preceq H\preceq bK$. In particular $H\succeq aI$.
On the subspace $\ker Z^\top$, of codimension at most $t$,
the Rayleigh quotient of $H$ is at most $b$. The min--max
principle proves the outlier count. The preconditioned system
$B_U^{1/2}AB_U^{1/2}$ has the same spectrum as $H$.

For completeness, CG minimizes the energy norm of the error over
its affine Krylov space: its conjugate search directions are
pairwise orthogonal in the system inner product and exact line
minimization makes the final error orthogonal to their span.
Apply this property to $M=B_U^{1/2}AB_U^{1/2}$ and the transformed
error $\widetilde e=B_U^{-1/2}e$, for which
$\|\widetilde e\|_M=\|e\|_A$. The Krylov error candidates are
$P(M)\widetilde e_0$ for polynomials of degree at most $h$ with
$P(0)=1$. Diagonalizing $M$ bounds their $M$-norm by
$\max_{\lambda\in\spec(M)}|P(\lambda)|\|e_0\|_A$.
Since $M$ and $H$ have the same spectrum, this is the required
polynomial bound in terms of $H$.

Let $\mu_1,\ldots,\mu_q>b$ be the outliers, with multiplicity.
Multiply the degree-$(h-q)$ Chebyshev residual polynomial on
$[a,b]$, normalized to equal one at zero, by
$\prod_{i=1}^q(1-z/\mu_i)$. The product vanishes at every outlier
and has absolute value at most one on $[a,b]$. The Chebyshev
factor is bounded by
$2((\sqrt\chi-1)/(\sqrt\chi+1))^{h-q}$ there, proving the result.
For $a=b$ the polynomial $(1-z/a)\prod_i(1-z/\mu_i)$ annihilates
the spectrum. The usual degree-$d$ annihilating polynomial gives
the dimension cap. Finally
\[
 \overline A-A_0=
 \begin{pmatrix}0&-C\\-C^\top&I-D\end{pmatrix}
 =\begin{pmatrix}0&-C\\0&I-D\end{pmatrix}
  +\begin{pmatrix}0&0\\-C^\top&0\end{pmatrix},
\]
a sum of two matrices of rank at most $t$. The stronger one-sided
outlier conclusion follows from the block inverse calculation,
rather than from this rank count alone. Every argument is
deterministic for every principal subset once $B_0$ satisfies its
initial spectral inequalities.
\end{proof}
\noindent Numerical check: \texttt{scripts/verify\_t11\_pcg.py}.

\paragraph{Approximate Cholesky instantiation.}
Contract the seed to a vertex $g$, and represent any diagonal
killing by edges to $g$. Let $\mathcal L$ be the resulting connected
Laplacian, whose principal matrix off $g$ is $A_0$. The
SparseCholesky theorem of Kyng--Sachdeva~\cite{kyng2016},
Theorem 3.1, supplies $\widetilde{\mathcal L}$ with
$(1-\epsilon_c)\mathcal L\preceq\widetilde{\mathcal L}
\preceq(1+\epsilon_c)\mathcal L$ except with probability $\delta_c$.
Its expected construction work and factor storage are
$O(\epsilon_c^{-2}(m_E+N)\log^2((N+1)/\delta_c)\log(N+1))$;
Appendix B gives the high-probability work variant with the stated
sampling rule. The sparse inverse application has work $O(F+N)$,
where $F$ is the stored factor size.

Define $Q$ to have rows $e_u^\top-e_g^\top$, $u\in W$.
Then $B_0=Q\widetilde{\mathcal L}^{+}Q^\top$ equals the inverse
of $\widetilde{\mathcal L}[W,W]$: solving with balanced source
$(r,-\ind^\top r)$ and subtracting the ground potential solves
the grounded system. Restriction of the Laplacian spectral
inequalities and inversion give
$a=(1+\epsilon_c)^{-1}$, $b=(1-\epsilon_c)^{-1}$, hence
$\chi=(1+\epsilon_c)/(1-\epsilon_c)$. At $\epsilon_c=1/2$,
$\chi=3$ and the iteration bound is $t+O(\log(1/\xi))$.
This instantiation uses the imported factorization theorem;
Theorem~\ref{thm:seedpcg} proves its simultaneous reuse under
grounding. The numerical check also verifies the $Q$ identity.

\begin{theorem}[Fixed-factor incremental work]\label{thm:inccost}
Let $k$ denote additional anchors, $J=p+1$, $L=\lceil J/2\rceil$,
and let $d_{\rm JL}$ be the initial sketch dimension. Let $r_t$
be the number of candidate refreshes in round $t=0,\ldots,k-1$,
$R=\sum_t r_t$, and let $s_t\leq L$ selected-column solves be
reused for the update in rounds $t<k-1$. Suppose every solve uses
the fixed preconditioner and reaches relative energy tolerance
at least as accurate as $\xi_*$. Set
$\mathcal H_t=O(\min\{N-t,t+1+\sqrt\chi\log(2/\xi_*)\})$.
An implementation with independent batched PCG columns has work
\begin{align*}
 T_{\rm build}+O\bigg(& (m_E+N+F)\left[
 Jd_{\rm JL}\mathcal H_0+
 \sum_{t=0}^{k-1}Lr_t\mathcal H_t+
 \sum_{t=0}^{k-2}(J-s_t)\mathcal H_t\right]\\
 &+kNJ^2+(R+k)\log(N+1)\bigg).
\end{align*}
For constant $\chi$ and $F,T_{\rm build}=\widetilde O(m_E+N)$,
putting $\ell=1+\log(2/\xi_*)$ gives the simpler bound
\[
 \widetilde O\left((m_E+N)
 [1+p d_{\rm JL}\ell+p(R+k)(k+\ell)]
       +kNp^2+(R+k)\log(N+1)\right).
\]
If each candidate is refreshed at most once, $R\leq N$.
For a hybrid implementation using Jacobi on some solves, replace
each corresponding $\mathcal H_t$ by
$O(\min\{N-t,\sqrt{\kappa_{J,t}}\log(2/\xi_*)\})$, and its
per-iteration work by $O(m_E+N)$, where $\kappa_{J,t}$ is the
Jacobi-preconditioned condition number.
\end{theorem}
\begin{proof}
The odd and even T6 chains together use $Jd_{\rm JL}$ initial
right-hand-side solves. At a refreshed candidate, the vectors
$y_i=A^{-i}e_v$, $i\leq L$, give
$r_{2i}=\|y_i\|^2$ and $r_{2i-1}=y_{i-1}^\top y_i$,
so each refresh needs $L$ solves. The full selected chain through
$A^{-J}e_v$ requires $J$ solves, of which $s_t$ can be reused.
The final round requires no following cache update. Each PCG
iteration has one sparse product, one factor application, and
linear vector work. Apply Theorem~\ref{thm:seedpcg} to each
right-hand side. Finite batches execute independent CG recurrences;
bounding every column by $\mathcal H_t$ also bounds their padded
batch work. Generating and accumulating sketches, moments, and
chain inner products is absorbed in these solve terms.
The scalar moment recurrences and all-node cache passes cost
$O(kNJ^2)$. Rebuilding each heap costs $O(N)$ per round, and
refresh pops plus selections cost $O((R+k)\log(N+1))$.
Substitute the largest $t<k$ and $J,L=O(p)$ for the simplified bound.
The Jacobi expression is the same Chebyshev proof with all
eigenvalues in its full spectral interval and no deflated outliers.
The distinct-refresh cap follows by counting eligible vertices.
\end{proof}
\noindent Numerical check: \texttt{scripts/verify\_t11\_complexity.py}.

\paragraph{Tolerance and factor parameters.}
For $a_*I\preceq A\preceq B_*I$, a recomputed relative residual
$\tau$ implies relative energy error at most
$\sqrt{B_*/a_*}\,\tau$: apply the two extremal norm inequalities
to $A^{-1/2}r$ and $A^{-1/2}b$. Conversely the residual ratio is
at most $\sqrt{B_*/a_*}$ times the relative energy error for a zero
start. Thus a residual-based stopping implementation can use
$\tau=\xi_*/\sqrt{B_*/a_*}$ and the work bound with energy target
$\xi_* /(B_*/a_*)$, adding only a logarithmic condition factor.
The simultaneous seed spectral event covers all later right-hand
sides and sets. With the T6 initial moment budgets, fixed $p$
has logarithmic solve-precision cost. The mixed setting used in
the night update applies the factor for batch width at least four;
the hybrid formula accounts separately for its single-column solves.

\subsection{Exact updates and retained moments}

\begin{proposition}[All-order independent decrements]\label{prop:incall}
Let $R=X_S^{-1}$ be zero padded on already selected coordinates,
let $v$ be active, $x=Re_v$, $c=1/R_{vv}$, and
$K=R-cxx^\top$. For $i\geq0$ define $a_i=R^ix$,
$s_i=x^\top R^ix$, and, for each active $u$,
\[
 \gamma_0(u)=a_0(u),\qquad
 \gamma_t(u)=a_t(u)-c\sum_{i=0}^{t-1}s_i\gamma_{t-1-i}(u).
\]
Then $\gamma_t(u)=e_u^\top K^tx\geq0$ and the moment decrement is
\[
 D_j(u)=r_j(u,S)-r_j(u,S\cup\{v\})
      =c\sum_{i=0}^{j-1}a_i(u)\gamma_{j-1-i}(u)\geq0.
\]
These decrements depend on selected columns, not on the cached
candidate moments. After a last refresh at time $b$,
\[
 |\widehat r_j(u,S_t)-r_j(u,S_t)|
 \leq E_{j,u}^{(b)}+\sum_{i=b}^{t-1}\varepsilon_{j,u}^{(i)},
 \qquad
 |\widehat D_j^{(i)}-D_j^{(i)}|\leq\varepsilon_{j,u}^{(i)}.
\]
For exact columns all $\varepsilon$ are zero. In an implementation
whose refresh flag changes only from false to true, whose refresh
batch contains only false-flag candidates, and whose nonseed
selections always have true flags, a run starting with all $N$
candidate flags false has exactly $N-R$ never-refreshed candidates
after $R$ total refreshes. All of them remain unselected.
\end{proposition}
\begin{proof}
For small $|z|$, the rank-one inverse identity gives
\[
 (I-zK)^{-1}x=
 \frac{(I-zR)^{-1}x}{1+czx^\top(I-zR)^{-1}x}.
\]
Equating coefficients proves the recurrence for $K^tx$.
The telescoping identity
\[
 R^j-K^j=\sum_{i=0}^{j-1}R^i(R-K)K^{j-1-i}
\]
proves the displayed decrement. Grounded inverses and their zero-padded
versions are entrywise nonnegative, so every displayed product is
nonnegative. Subtracting the computed update from the true update
and telescoping proves the error envelope; refreshing restarts it
with that refresh's error budget. Finally each refresh changes one
distinct false flag to true, so exactly $R$ flags have changed.
Only true flags can be selected, proving the last assertion.
\end{proof}
\noindent Numerical check: \texttt{scripts/verify\_t11\_incremental.py}.

\paragraph{Computable decrement-error envelopes.}
Let a selected-column chain have true-residual bounds $\rho_i$ and
spectral lower bound $a_*>0$. With $a_{-1}=e_v$, its Euclidean
errors satisfy $E_{-1}=0$, $E_i\leq(E_{i-1}+\rho_i)/a_*$.
Use component intervals $\widehat a_i(u)\pm E_i$,
and intervals for $s_i$ centered at $\widehat a_0^\top\widehat a_i$
with radius
$\|\widehat a_0\|E_i+\|\widehat a_i\|E_0+E_0E_i$.
The reciprocal of a positive interval for $a_0(v)$ encloses $c$.
Propagating these intervals through the finite $\gamma_t,D_j$
recurrences encloses every true decrement by induction. Their
distance from the computed decrement gives the $\varepsilon$
above. This construction has the same $O(NJ^2)$ scalar order as
the update. The residual inequality follows by applying $R$ to
each chain equation; the inner-product radius follows by expansion
and Cauchy--Schwarz. The incremental check tests both constructions.

\begin{theorem}[Moment retention and killed returns]\label{thm:retention}
For an unselected $u$ define
$\rho_j(u,T)=r_j(u,S_0\cup T)/r_j(u,S_0)$, $j>0$.
There is a probability law on killed return paths from $u$ such that
\[
 \rho_j(u,T)=1-\Pr_{u,j}(H_T\leq t).
\]
Its unnormalized path/time measure is
$t^{j-1}\mathbf1\{X_t=u,t<\zeta_{S_0}\}\,dt$, with total
mass $\Gamma(j)r_j(u,S_0)$. In particular $0<\rho_j\leq1$.
For integer $j\geq1$, writing $\alpha_1=1-\rho_1$,
\[
 1-\rho_j\leq
 \min\left\{1,
 \frac{\sqrt{\alpha_1\Gamma(2j-1)r_{2j-1}(u,S_0)r_1(u,S_0)}}
      {\Gamma(j)r_j(u,S_0)}\right\}.
\]
The second expression is at most
$\sqrt{\binom{2j-2}{j-1}\alpha_1\kappa_0^{j-1}}$,
where $\kappa_0$ is the condition number of $X_{S_0}$.

Let $R_0=X_{S_0}^{-1}$. For nonempty $T$,
\[
 \alpha_1=
 \frac{R_0[u,T]R_0[T,T]^{-1}R_0[T,u]}{R_0[u,u]}
 \leq\Pr_u(H_T<\zeta_{S_0}).
\]
For $T=\{v\}$ this is
$\alpha_1=\Pr_u(H_v<\zeta_{S_0})\Pr_v(H_u<\zeta_{S_0})
=R_0[u,v]^2/(R_0[u,u]R_0[v,v])$.
In the contracted electrical network, put
$d_u=\mathcal R(u,g)$, $d_v=\mathcal R(v,g)$, and
$d_{uv}=\mathcal R(u,v)$. Then
$\alpha_1=(d_u+d_v-d_{uv})^2/(4d_ud_v)$.

If the initial relative error is at most $\zeta_0$, a
never-refreshed cache satisfies
\[
 \frac{|\widehat r_j(u,S)-r_j(u,S)|}{r_j(u,S)}
 \leq\frac{\zeta_0}{\rho_j(u,T)}+
       \frac{\sum_i\varepsilon_{j,u}^{(i)}}{r_j(u,S)}.
\]
\end{theorem}
\begin{proof}
Use T1's common killed walk, additionally killed at its first hit of
$T$. The spectral integral
$X^{-j}=\Gamma(j)^{-1}\int_0^\infty t^{j-1}e^{-tX}\,dt$
and the killed return identity give the stated measure and ratio.
All integrands are nonnegative, justifying Tonelli. Its mass is
finite since $X_{S_0}\succ0$, and the survivor diagonal is positive.
Under the $j=1$ law, the loss at order $j$ is the weighted
expectation of the hit indicator with weight $t^{j-1}$, divided
by its mean. Cauchy--Schwarz bounds its numerator by
$\sqrt{\alpha_1\E_{u,1}t^{2j-2}}$.
Since $\E_{u,1}t^h=\Gamma(h+1)r_{h+1}/r_1$, this proves the
bound. Spectral decomposition gives
$r_{2j-1}\leq\lambda_{\min}^{-(j-1)}r_j$ and
$r_1\leq\lambda_{\max}^{j-1}r_j$, proving the condition bound.

The zero-padded block inverse formula after grounding $T$ gives
\[
 R_T=R_0-R_0[:,T]R_0[T,T]^{-1}R_0[T,:].
\]
For the hitting interpretation, the strong Markov property for
the Green occupation kernel yields
\[
 R_0[u,u]-R_T[u,u]
 =\sum_{v\in T}\Pr_u(X_{H_T}=v,H_T<\zeta_{S_0})R_0[v,u].
\]
The identity $R_0[v,u]/R_0[u,u]=\Pr_v(H_u<\zeta_{S_0})$
follows by applying the same property at the first hit of $u$.
It is at most one, proving the probability bound and singleton
factorization. The voltage solution to a unit current from $u$
to $g$ has value $R_0[u,u]$ at $u$, so $d_u=R_0[u,u]$.
Similarly $d_{uv}=(e_u-e_v)^\top R_0(e_u-e_v)$, proving the
resistance expression. The last inequality divides
Proposition~\ref{prop:incall}'s absolute-error envelope by the
current moment.
\end{proof}
\noindent Numerical check: \texttt{scripts/verify\_t11\_retention.py}.

\paragraph{Higher-moment sensitivity.}
In the connected three-vertex graph with seed-edge weights $t,1$
and follower-edge weight $t$, let $T$ be the first follower and
$u$ the second. The initial matrix is
$\left(\begin{smallmatrix}2t&-t\\-t&1+t\end{smallmatrix}\right)$.
Its Green matrix is
\[
 R_0=\frac1{t(2+t)}
       \begin{pmatrix}1+t&t\\t&2t\end{pmatrix}.
\]
Thus $\rho_1=(2+t)/(2(1+t))\to1$ and
$\Pr_u(H_T<\zeta)=t/(1+t)\to0$.
Yet the term $R_0[u,T]^2R_0[T,T]$ in $(R_0^3)_{uu}$
is asymptotic to $1/(8t)$, while the current third moment is
$(1+t)^{-3}\leq1$. Hence $\rho_3\to0$. The retention check
verifies this example as well as the weighted-return integrals.

\begin{theorem}[Conditional guarantee for practical incremental keys]
\label{thm:practicalinc}
Fix $p\geq1$, $0<\epsilon\leq1/2$, and $k$ additional anchors.
Suppose the following conditions hold throughout a run:
\begin{enumerate}
\item Every never-refreshed active moment of order $1\leq j\leq p+1$
has retention $\rho_j\geq\gamma>0$. The initial simultaneous
relative error is at most $\zeta_0\leq\gamma\epsilon/(128p)$.
\item For never-refreshed entries, cumulative decrement error is
at most $\epsilon/(128p)$ times the current moment. For refreshed
entries, refresh error plus subsequent decrement errors is at
most $\epsilon/(64p)$ times the current moment.
\item Every candidate key is its current estimated marginal,
possibly multiplied by a factor at least one. The selected
candidate's evaluated value is at least $1-\theta$ times every
remaining key and at least every other evaluated value, with
$0\leq\theta\leq\epsilon/2$.
\end{enumerate}
Then the returned seeded gain is at least
$(1-e^{-1}-\epsilon)\mathrm{OPT}_k$.
If the retention and numerical conditions are enforced pathwise,
an initial T6 event of probability $1-\delta$ yields the same
success probability for the complete adaptive run. Sufficient
Gaussian dimension is
$d_{\rm JL}=O(p^2\gamma^{-2}\epsilon^{-2}
\log(N(p+1)/\delta))$, with T6's solve budgets.
\end{theorem}
\begin{proof}
Theorem~\ref{thm:retention} bounds every current moment's relative
error by $\zeta_*\leq\epsilon/(64p)$; this also holds for the
refreshed class by hypothesis. Theorem~\ref{thm:polymoment}
therefore bounds every marginal estimate relatively by
$\eta=\epsilon/4$. At acceptance, the evaluated value $b$ of the
selected vertex satisfies
$b\geq(1-\theta)(1-\eta)\max_u\Delta_u$:
every other vertex is either evaluated or has a remaining key,
and every such value/key is at least $(1-\eta)\Delta_u$.
The selected true marginal is at least $b/(1+\eta)$, hence
its fraction of the maximum is at least
\[
 (1-\theta)\frac{1-\eta}{1+\eta}
 \geq1-\theta-2\eta\geq1-\epsilon.
\]
Apply the T7 gap-closing proof. The initial JL event is uniform
over the initial candidates and moments; the exact identities and
pathwise numerical bounds apply to every subsequent adaptive
choice, so no sketch reuse on a new random matrix is assumed.
Substitution of $\zeta_0=\Theta(\gamma\epsilon/p)$ into T6
gives the dimension and the stated success probability.
\end{proof}
\noindent Numerical check: \texttt{scripts/verify\_t11\_practical.py}.

\paragraph{Checking the sufficient conditions.}
Valid intervals $[L_j(S),U_j(S)]$ and $[L_j(S_0),U_j(S_0)]$
give the computable retention lower bound $L_j(S)/U_j(S_0)$.
Alternatively an absolute error envelope $E_j(S)$ certifies the
current relative tolerance directly whenever
$E_j(S)\leq\zeta_*L_j(S)$. These follow by substituting the
true moments between their respective endpoints. Residual and
interval refinement can enforce the conditions used in
Theorem~\ref{thm:practicalinc}. The factor and initial-sketch
failure budgets add by a union bound when combining its gain
guarantee with Theorem~\ref{thm:inccost}'s work guarantee.

\paragraph{Refresh counts in the inspected night runs.}
The current \texttt{inc} implementation has the flag semantics of
Proposition~\ref{prop:incall}. Reading its completed run outputs,
with one seed and no exact initialization, gives the following
counts. Percentages use the initial $N=n-1$ candidates as denominator.
The reported refresh total is a distinct-candidate count in this mode.
\begin{center}
\begin{tabular}{lrrrr}
Graph / exponent / slack & $N$ & $R$ & Never refreshed & Percent\\\hline
Amazon / 1 / 0 &334862&736&334126&99.780\\
Amazon / 1 / 0.2 &334862&1232&333630&99.632\\
DBLP / 1 / 0 &317079&256&316823&99.919\\
DBLP / 1 / 0.2 &317079&448&316631&99.859\\
Amazon / 2 / 0.2 &334862&1504&333358&99.551
\end{tabular}
\end{center}
The incremental check verifies these arithmetic counts and the flag
invariant. The source paths and the unmeasured retention conditions
are recorded in \texttt{notes/gaps.md}.

\subsection{Low-order marginals and Zhou's spectral score}

\begin{proposition}[Finite-removal expansion and spectral weighting]
\label{prop:zhoutaylor}
Let $X$ be a grounded Laplacian with a simple smallest eigenpair
$(\lambda_1,\phi_1)$, $\|\phi_1\|=1$, and orthonormal eigenbasis
$\phi_i$. For an active candidate $v$, let $R=X^{-1}$,
$x=Re_v$, $c=1/r_1$, $D=cxx^\top$, and $a_j=cr_{j+1}$.
For $0\leq z\leq1$ put
$Q_p(z)=\tr R^p-\tr(R-zD)^p$. Its endpoint $Q_p(1)$ is the
hard-grounding marginal. For $p=2,3,4$,
\begin{align*}
 Q_2(z)&=2za_2-z^2a_1^2,\\
 Q_3(z)&=3za_3-3z^2a_1a_2+z^3a_1^3,\\
 Q_4(z)&=4za_4-z^2(4a_1a_3+2a_2^2)
                       +4z^3a_1^2a_2-z^4a_1^4.
\end{align*}
Their first-order coefficient satisfies
\[
 Q_p'(0)=\frac p{r_1(v)}\sum_i\lambda_i^{-(p+1)}\phi_i(v)^2.
\]
For the weighted Zhou score
$h_v=2\phi_1(v)\sum_{w\ne v}w_{vw}\phi_1(w)$,
\[
 h_v=2(X_{vv}-\lambda_1)\phi_1(v)^2.
\]
Whenever $X_{vv}>\lambda_1$,
\[
 Q_p'(0)=\frac{p\lambda_1^{-(p+1)}h_v}
                    {2r_1(v)(X_{vv}-\lambda_1)}
       +\frac p{r_1(v)}\sum_{i\geq2}
                          \lambda_i^{-(p+1)}\phi_i(v)^2.
\]
The tail sum is at most
$(1-\phi_1(v)^2)\lambda_2^{-(p+1)}$.
For every integer $p\geq2$,
\[
 0\leq zQ_p'(0)-Q_p(z)
 \leq\frac{p(p-1)}2 z^2 a_1^2\|R\|^{p-2}.
\]
An infinitesimal soft penalty instead has coefficient
$p\sum_i\lambda_i^{-(p+1)}\phi_i(v)^2$, without the division by
$r_1(v)$.
\end{proposition}
\begin{proof}
Expand the trace of $(R-zD)^p$ as noncommutative words and use
cyclicity of the trace. Since $D$ has rank one,
$\tr(DR^{j-1})=a_j$ and
$\tr(DR^iDR^h)=a_{i+1}a_{h+1}$.
For degree four the six words with two $D$'s have traces
$4a_1a_3+2a_2^2$; the four with three $D$'s give
$4a_1^2a_2$. The degree-two and degree-three counts give the
remaining formulas. Differentiating at zero gives $pa_p$ and
spectral decomposition gives its eigenmode expression. The
$v$-coordinate of $X\phi_1=\lambda_1\phi_1$ proves the identity
for $h_v$, and substitution gives the comparison. The tail bound
uses $\sum_{i\geq2}\phi_i(v)^2=1-\phi_1(v)^2$.

Put $K_z=R-zD\succeq0$. Differentiation gives
$Q_p'(z)=p\tr(DK_z^{p-1})$ and
\[
 -Q_p''(z)=p\sum_{i=0}^{p-2}\tr(DK_z^iDK_z^{p-2-i}).
\]
Write $D=a_1ww^\top$ with $\|w\|=1$. Each summand equals
$a_1^2(w^\top K_z^iw)(w^\top K_z^{p-2-i}w)$, between zero
and $a_1^2\|R\|^{p-2}$ since $0\preceq K_z\preceq R$.
Integrating twice gives the remainder bound. Finally
$R-zD=(X+\beta e_ve_v^\top)^{-1}$ for
$\beta=z/((1-z)r_1)$ when $z<1$, by Sherman--Morrison.
Thus $dz/d\beta=r_1$ at zero and the soft coefficient follows.
\end{proof}
\noindent Numerical check: \texttt{scripts/verify\_t11\_zhou.py}.

\paragraph{Mechanism and scope of simultaneous improvement.}
At $p=2,3,4$ the first-order modal weights are respectively
$\lambda^{-3},\lambda^{-4},\lambda^{-5}$. The finite-removal
polynomials additionally account for the normalization $r_1(v)$
and higher-order terms. They offer a precise mechanism for choices
that differ from a single-eigenvector score while retaining
weight on several slow modes. The four-node stress case in
Proposition~\ref{prop:zhoustress} gives strict improvement in both
metrics: its inverse-trace after selecting $a$ is at most two,
whereas after Zhou's choice it exceeds two, and its smallest
eigenvalue is larger.

The opposite spectral ordering also occurs within the same family.
Consider the grounded matrix
\[
 X=\begin{pmatrix}1&-9/10&0\\-9/10&3&-1/2\\0&-1/2&1\end{pmatrix}.
\]
Its three row sums are $1/10,8/5,1/2$, so it is realized by a
connected weighted graph with one seed attached by those weights.
The follower path and positivity imply that Zhou uniquely chooses
the middle vertex. Its resulting matrix is $I_2$, giving spectral
value one and inverse trace two. Deleting the first vertex gives
eigenvalues $2\pm\sqrt5/2$, inverse trace $16/11$, and inverse
power traces $168/121,1984/1331,24352/14641$ for $p=2,3,4$,
each below two. Deleting the third leaves a smallest eigenvalue
$2-\sqrt{181}/10<7/10$, so every $p\geq2$ inverse-power trace
then exceeds two. Thus each of these three exact power-greedy
choices improves CFCC but has spectral value $2-\sqrt5/2<1$.
These explicit calculations, checked by
\texttt{scripts/verify\_t11\_zhou.py}, distinguish a possible
joint improvement from a universal ordering.

\section{Road-network diagnosis: maintaining the sketches themselves}

\begin{proposition}[Current-matrix sketch maintenance for hard CFCC]
\label{prop:maintainedsketch}
Let $A_0=X_{S_0}=C^\top C\succ0$ have dimension $N$.
For $T\subset W$, write $R_T$ for the zero-padded inverse of
$A_0[W\setminus T,W\setminus T]$. Fix matrices $G,H$ and maintain
\[
 Y_1=GCR_T,\qquad Y_2=HR_T.
\]
When grounding an active vertex $v$, one solve for $x=R_Te_v$
updates both sketches exactly by
\[
 Y_i' =Y_i-\frac{Y_i[:,v]x^\top}{x_v},\qquad i=1,2.
\]
For every fixed $T$, if $G,H$ are independent Gaussian matrices
with $d$ rows and entry variance $1/d$, then
$\widehat r_1(u)=\|Y_1[:,u]\|^2$ and
$\widehat r_2(u)=\|Y_2[:,u]\|^2$ are unbiased estimates of the
\emph{current} moments $r_1(u,T),r_2(u,T)$.

For $0<\zeta<1$, $0<\delta<1$, and a maximum of $k<N$ new anchors,
put $M=\sum_{t=0}^k\binom Nt(N-t)$. If
\[
 d\geq8\zeta^{-2}\log(4M/\delta),
\]
both estimates have simultaneous relative error at most $\zeta$
at every active vertex of every such set, with probability at
least $1-\delta$. This event covers adaptive choices made from
the maintained sketches. Choosing the maximum estimated ratio
$\widehat r_2/\widehat r_1$, or exactly rescoring a candidate
pool containing that maximizer and choosing its best member,
achieves a maximum-marginal fraction at least
$\alpha=((1-\zeta)/(1+\zeta))^2$. For $k$ selections its
normalized gain is at least $(1-e^{-\alpha})\mathrm{OPT}_k$.
In particular $\zeta\leq\epsilon/4$ gives the T7
$(1-e^{-1}-\epsilon)$ factor. The sufficient row count has order
$O(\zeta^{-2}((k+1)\log(N+1)+\log(1/\delta)))$.

Each update costs $O(Nd)$ scalar work and uses $O(Nd)$ sketch
storage. The selected-column solve can be reused from rescoring;
no $R_T^2e_v$ solve is needed for this sketch update.

Moreover, for any row-vector error $e$ supported on the active
coordinates, its exact-column update $e'=e-e_vx^\top/x_v$ obeys
\[
 \|e'\|_{A_0}^2=\|e\|_{A_0}^2-e_v^2/x_v.
\]
Thus accumulated initial sketch-solve errors are contractive in
the fixed $A_0$ energy norm under exact subsequent columns.
\end{proposition}
\begin{proof}
Sherman--Morrison in its hard zero-padded form gives
$R_{T+v}=R_T-xx^\top/x_v$. Multiplying on the left by $GC$ or
$H$ proves the update because $Y_i[:,v]$ is that left factor
times $x$. On active coordinates $R_TA_0R_T=R_T$, and all
remaining rows and columns of both sides are zero. Consequently
\[
 \|CR_Te_u\|^2=r_1(u,T),\qquad
 \|R_Te_u\|^2=r_2(u,T).
\]
For a fixed deterministic set and candidate, Gaussian norm
concentration gives failure at most $2e^{-d\zeta^2/8}$ for each
of the two vectors. There are $2M$ vectors across all sets and
candidates, so the stated row count gives the joint event.
Since it holds simultaneously for this deterministic family,
it also holds along every adaptive selection path. The bound
$M\leq N(k+1)N^k$ gives the stated asymptotic dimension.

On the event each ratio is between $(1-\zeta)/(1+\zeta)$ and
$(1+\zeta)/(1-\zeta)$ times its true marginal. Comparing an
estimated maximizer to a true maximizer therefore gives $\alpha$.
Exact rescoring a pool containing that candidate can only improve
the true marginal at the same current set. The T7 induction gives
$1-e^{-\alpha}$; since $1-\alpha=4\zeta/(1+\zeta)^2\leq4\zeta$,
the final approximation follows. The two rank-one outer products
and recomputation of column norms use $O(Nd)$ operations and
storage, in addition to one selected-column solve.

For a supported row $e$, the active inverse equations imply
$eA_0x=e_v$ and $x^\top A_0x=x_v$. Expanding the squared energy
norm gives the last identity. The new row has zero at $v$ and
at every previous anchor, so the contraction can be iterated.
\end{proof}
\noindent Numerical check:
\texttt{scripts/verify\_road\_sketch\_update.py}.

\paragraph{Initial solve errors in the maintained sketches.}
For a matrix error $E$ with supported rows, define
$\|E\|_{F,A_0}^2=\tr(EA_0E^\top)$.
Cauchy--Schwarz with the representer $R_Te_u$ gives
$\|E[:,u]\|\leq\sqrt{r_1(u,T)}\|E\|_{F,A_0}$.
Let $B\geq\lambda_{\max}(A_0)$ and let $e_i$ bound the initial
energy errors for the two sketch matrices. The contraction above
keeps those bounds valid under exact-column updates. On the
simultaneous Gaussian event, the relative squared-norm errors
are at most
\[
 \zeta+2\sqrt{1+\zeta}\,e_1+e_1^2,
 \qquad
 \zeta+2\sqrt{1+\zeta}\sqrt B\,e_2+B e_2^2,
\]
respectively. To obtain the second bound, use the spectral
inequality $r_1\leq B r_2$. Both follow by expanding the
perturbed norm square and bounding the cross term. They are
also checked by the script above.

\section{T12: adaptive hybrid sketches and higher powers}

This section uses hard grounding and the fixed admissible seed
$S_0$. Write $W=V\setminus S_0$, $N=|W|$,
$A_0=C^\top C\succ0$, and let $R_T$ be the zero-padded inverse
on $W\setminus T$. The budget $k<N$ counts \emph{additional}
anchors. In particular, a run with one seed and twenty total
anchors has $k=19$ here.

\begin{proposition}[Higher-power sketch maintenance]
\label{prop:highermaintained}
Fix a matrix $F$ and an integer $L\geq1$, and store
$Y_\ell=FR_T^\ell$ for $1\leq\ell\leq L$.
Before grounding $v$, put $R=R_T$, $x=Re_v$, $c=1/x_v$,
$a_i=R^ix$, and $s_i=x^\top a_i$. Define
\[
 b_0=x,\qquad
 b_t=a_t-c\sum_{i=0}^{t-1}s_i b_{t-1-i}
 \quad(1\leq t<L).
\]
All levels update simultaneously by
\begin{equation}\label{eq:highermaintained}
 Y_\ell'=Y_\ell-c\sum_{i=0}^{\ell-1}
             Y_{i+1}[:,v]b_{\ell-1-i}^\top.
\end{equation}
Every vector on the right is its pre-update value.
Apply this construction to $F=H$ and $F=GC$, with $d$ rows each.
The squared columns of $HR_T^\ell$ and $GCR_T^\ell$ estimate,
respectively, $r_{2\ell}(u,T)$ and $r_{2\ell-1}(u,T)$.
Thus moments through $p+1$ need $L=\lceil(p+1)/2\rceil$.

Initialization uses at most $2dL$ right-hand-side solves.
An update uses at most $L$ selected-column solves,
$O(NdL^2)$ arithmetic, and $O(NdL)$ storage; existing selected
columns can be reused. The formula specializes to the two
rank-one updates of Proposition~\ref{prop:maintainedsketch}
when $p=1$. For $p=2$, it additionally gives
\[
 Y_2'=Y_2-cY_1[:,v](Rx-c\|x\|^2x)^\top
                 -cY_2[:,v]x^\top.
\]
\end{proposition}
\begin{proof}
Let $K=R-cxx^\top=R_{T+v}$. The generating-function calculation
in T11 gives $b_t=K^tx$; alternatively it follows by comparing
coefficients in
\[
 (I-zK)^{-1}x=
 \frac{(I-zR)^{-1}x}{1+czx^\top(I-zR)^{-1}x}.
\]
The noncommutative telescoping identity
\[
 R^\ell-K^\ell=c\sum_{i=0}^{\ell-1}R^ixx^\top K^{\ell-1-i}
\]
follows by expanding the consecutive differences
$R^{i+1}K^{\ell-1-i}-R^iK^{\ell-i}$ and summing.
Multiplication by $F$ proves \eqref{eq:highermaintained}, since
$FR^ix=Y_{i+1}[:,v]$ and $K$ is symmetric.
Snapshotting all columns $Y_\ell[:,v]$ before any update gives
an in-place implementation of this identity.

For every $T$, $R_TA_0R_T=R_T$. Therefore
\[
 \|R_T^\ell e_u\|^2=r_{2\ell}(u,T),\qquad
 \|CR_T^\ell e_u\|^2=r_{2\ell-1}(u,T).
\]
These identities prove the moment targets for isotropic sketches.
The chain $a_0,\ldots,a_{L-1}$ uses $L$ solves; its inner
products and the $b_t$ recurrence cost $O(NL^2)$.
There are $\sum_{\ell=1}^L\ell=O(L^2)$ outer products per
sketch family, costing $O(NdL^2)$. The stored levels and
selected chains take $O(NdL)$ space. Repeatedly solving each
of the two families of $d$ starting rows gives initialization.
For $L=2$, $b_1=Rx-c\|x\|^2x$, giving the final formula.
\end{proof}
\noindent Numerical check:
\texttt{scripts/verify\_t12\_higher\_sketch.py}.

\begin{theorem}[Uniform adaptive moment accuracy]
\label{thm:adaptivehybridjl}
Let all entries of $G,H$ be independent centered Gaussians of
variance $1/d$, or independent signs $\pm1/\sqrt d$.
For $J=p+1$, $0<\zeta,\delta<1$, put
\[
 M_{N,k}=\sum_{t=0}^k\binom Nt(N-t),\qquad
 d_{\rm reuse}=\left\lceil
    8\zeta^{-2}\log\frac{2JM_{N,k}}{\delta}\right\rceil.
\]
With $d\geq d_{\rm reuse}$ and exact sketch maintenance,
all $J$ moment estimates have relative error at most $\zeta$
at every active vertex of every $|T|\leq k$, with probability
at least $1-\delta$. The same event covers every adaptive
selection rule, refresh pattern, and stopping decision.
With $\zeta\leq\eta/(16p)$, the computed integer marginals
have relative error at most $\eta$ by
Theorem~\ref{thm:polymoment}.

For comparison, new independent sketches in each of $k$ rounds
need only
\[
 d_{\rm fresh}=\left\lceil
       8\zeta^{-2}\log\frac{2JNk}{\delta}\right\rceil
\]
for the same ideal-solve moment event along the chosen path.
In particular, the sufficient dimensions have orders
\begin{align*}
 d_{\rm reuse}
 &=O\!\left(p^2\eta^{-2}
       \left[(k+1)\log(N+1)+\log(p/\delta)\right]\right),\\
 d_{\rm fresh}
 &=O\!\left(p^2\eta^{-2}\log(Npk/\delta)\right).
\end{align*}
\end{theorem}
\begin{proof}
For a deterministic nonzero vector $w$, let $Z$ be one
unscaled sketch row dotted with $w/\|w\|$. In either distribution,
$\E Z^2=1$, $\E Z^4\leq3$, and
$\E e^{tZ}\leq e^{t^2/2}$. For signs the last inequality follows
from $\cosh x\leq e^{x^2/2}$ and independence. Integrating
against an auxiliary standard Gaussian yields
$\E e^{tZ^2}\leq(1-2t)^{-1/2}$ for $0<t<1/2$.
Chernoff's bound with $t=\zeta/[2(1+\zeta)]$ bounds the upper
tail of the average of $d$ independent $Z^2$ by
\[
 \exp\{-d[\zeta-\log(1+\zeta)]/2\}
 \leq e^{-d\zeta^2/8}.
\]
Here $\zeta-\log(1+\zeta)\geq\zeta^2/4$ on $[0,1]$, as
follows by differentiating the difference.
For the lower tail, $e^{-y}\leq1-y+y^2/2$ for $y\geq0$
gives $\E e^{-tZ^2}\leq\exp(-t+3t^2/2)$.
Taking $t=\zeta/3$ gives $e^{-d\zeta^2/6}$.
The combined two-sided failure is at most $2e^{-d\zeta^2/8}$.

The targets $R_T^\ell e_u$ and $CR_T^\ell e_u$ form a
\emph{deterministic} family with $JM_{N,k}$ members when only
the first $J$ moments are retained. They are nonzero at active
vertices. The preceding tail bound and a union bound prove the
claim. Their correlations, including shared rows across powers,
do not affect this union bound. Once the event holds, adaptively
choosing a member of this family cannot invalidate it.
The higher-power identity computes precisely these fixed-randomness
sketches. Unbiasedness for each deterministic $T$ supplies the
moment identity, while this simultaneous event supplies the
adaptive guarantee.

With fresh randomness, condition on the complete history at each
round; its current matrix is then fixed independently of the new
rows. A union bound over at most $NJ$ current moments and a
conditional failure budget $\delta/k$ proves the fresh formula.
Finally $M_{N,k}\leq N(k+1)N^k$ and
$\zeta=\eta/(16p)$ prove the stated orders.
\end{proof}
\noindent Numerical check:
\texttt{scripts/verify\_t12\_adaptive.py}.

\begin{theorem}[Hybrid retention and capped refreshes]
\label{thm:hybridcap}
At each state, suppose every not-yet-refreshed marginal estimate
satisfies
$(1-\eta)\Delta_u\leq\widehat\Delta_u\leq(1+\eta)\Delta_u$,
and refreshed candidates are maintained exactly by T11.
Use the estimate itself as the heap key for an unrefreshed
candidate and the exact marginal for a refreshed one.
Let $0\leq\theta<1$ be an acceptance tolerance. Process keys
in descending order, exactly rescoring batches of unrefreshed
candidates and retaining the best exact value $b$. Stop either
when $b\geq(1-\theta)K$, where $K$ is the largest remaining
key, or after a positive refresh budget has been reached, and
return the best processed candidate. Every processed refreshed
candidate is included in this comparison. If a budget stop occurs,
at least the initial maximum-key candidate has been processed.
Then the selected marginal is at least
\[
 \alpha\max_u\Delta_u,\qquad
 \alpha=\min\left\{\frac{1-\eta}{1+\eta},
                           (1-\theta)(1-\eta)\right\}.
\]
The $k$-round normalized gain is at least
$(1-e^{-\alpha})\mathrm{OPT}_k$.
In particular $\eta\leq\epsilon/4$ and $\theta\leq\epsilon/2$
give $(1-1/e-\epsilon)\mathrm{OPT}_k$ on the uniform event of
Theorem~\ref{thm:adaptivehybridjl}.

For $p=1$, put $\chi=(1+\zeta)/(1-\zeta)$ on its moment event.
The sharper factor is
$\alpha=\min\{\chi^{-2},(1-\theta)\chi^{-1}\}$; thus
$\zeta\leq\epsilon/8$ and $\theta\leq\epsilon/2$ suffice.
Arbitrarily many candidates can be kept exact after their first
refresh. Their total number of refreshes is $R\leq N$.
\end{theorem}
\begin{proof}
Both types of key lie between $(1-\eta)\Delta_u$ and
$(1+\eta)\Delta_u$. If $u_*$ maximizes the true marginal and
$w$ has maximum initial key, then
\[
 \Delta_w\geq\frac{\mathrm{key}(w)}{1+\eta}
 \geq\frac{1-\eta}{1+\eta}\Delta_{u_*}.
\]
A budget stop has processed $w$ and returns a candidate no worse
than $w$. At an acceptance stop, every processed true value is
at most $b$, and every unprocessed one is at most $K/(1-\eta)$.
Thus $b\geq(1-\theta)(1-\eta)\Delta_{u_*}$.
Taking the minimum covers both stopping rules, including an
empty remaining heap. Positive-budget batching may overshoot
its nominal size without changing this argument.

By T1 the gain is normalized, monotone, and submodular relative
to $S_0$. At each round its remaining gap to the best $k$-element
extension is at most $k$ times the largest current marginal.
The gap therefore contracts by at least $1-\alpha/k$ at every
step, proving $1-(1-\alpha/k)^k\geq1-e^{-\alpha}$.
Moreover $1-\alpha\leq\max\{2\eta,\eta+\theta\}\leq\epsilon$,
and $e^{-\alpha}-e^{-1}\leq1-\alpha$ proves the claimed factor.
For $p=1$ each ratio key is between $\Delta_u/\chi$ and
$\chi\Delta_u$. Repeat the same proof with these endpoints.
Here $1-\chi^{-2}\leq4\zeta$ and
$1-(1-\theta)\chi^{-1}\leq\theta+2\zeta$.
Exact retention follows from the independent-column decrement
identity in T11; each false-to-true refresh flag changes at most
once. Replacing any sketched value by its exact value preserves
all the preceding key inequalities, even for an adaptively
chosen subset of candidates.
\end{proof}
\noindent Numerical check:
\texttt{scripts/verify\_t12\_hybrid.py}.

\paragraph{Work of the maintained and hybrid variants.}
Let $\mathcal T_t$ bound one required-precision solve on the
current active matrix, and let $T_{\rm refresh}$ include all
candidate rescoring. Pure higher-power sketch maintenance costs
\[
 O\!\left(dm+2dL\mathcal T_0+
       L\sum_{t=0}^{k-1}\mathcal T_t+
       kNdL^2+kNp^2+T_{\rm refresh}\right),
\]
using $O(NdL)$ sketch memory in addition to solver storage.
The $kNp^2$ term evaluates integer marginals from moments.
For hybrid retention, the selected chain for scalar exact updates
extends to $p+1$ columns, replacing $L$ by $p+1$ in the selected
solve term; the T11 scalar work is $O(Nkp^2)$ and its storage
$O(Np)$. Heap rebuilding costs $O(kN)$ and removals
$O((R+kR+k)\log N)$, allowing every previously refreshed
candidate to be processed in every later round.
At $p=1$ this gives two initial sketch families, one selected
solve for sketch-only maintenance or two for the hybrid, and
$O(Nd)$ work per update. Reuse of refreshed chains reduces the
actual solve count. The adaptive union bound increases $d$;
it trades repeated all-node solves for maintained arrays, rather
than identifying the two algorithms' worst-case work.
The higher-sketch and hybrid scripts check these operation counts.

\begin{proposition}[Residual budgets for maintained sketches]
\label{prop:sketchresidual}
Let $a\leq\lambda_{\min}(A_0)$ and
$B\geq\lambda_{\max}(A_0)$. At a $p=1$ update, let
$\widetilde x$ be supported on the current active set $U$, with
$\widetilde x_v>0$ and
$\|A_0[U,U]\widetilde x_U-e_v\|_2\leq\rho$.
For either sketch, perform
\[
 \widetilde Y'=
 \widetilde Y-\widetilde Y[:,v]\widetilde x^\top/
                 \widetilde x_v+D,
\]
where $D$ has zero columns on all new and old anchors and
$\|D\|_{F,A_0}\leq\omega$.
If $e\geq\|\widetilde Y-Y\|_{F,A_0}$ before the update, then
\begin{equation}\label{eq:sketchresidual}
 e'=e+\frac{\|\widetilde Y[:,v]\|_2\rho}
                   {\sqrt a\,\widetilde x_v}+\omega
\end{equation}
is a valid bound afterwards. This recurrence is deterministic
along adaptive selections. At every state the two moment errors
are bounded by
\[
 \zeta+2\sqrt{1+\zeta}\,e_1+e_1^2,\qquad
 \zeta+2\sqrt{1+\zeta}\sqrt B\,e_2+B e_2^2
\]
on the ideal-sketch event, where $e_1,e_2$ are their current
accumulated energy bounds.

For higher powers, a sufficient alternative certificate uses
active-coordinate residuals
$\rho_\ell=\|\widetilde Y_\ell A_U-\widetilde Y_{\ell-1}\|_F$,
with $\widetilde Y_0=F[:,U]$ exact. Recursively let
$E_0=0$ and $E_\ell=(E_{\ell-1}+\rho_\ell)/a$.
Then $\|\widetilde Y_\ell-F[:,U]A_U^{-\ell}\|_F\leq E_\ell$.
For even moments take $e=B^\ell E_\ell$, and for odd moments
take $e=B^{\ell-1/2}E_\ell$. Their relative squared-norm error
is at most $\zeta+2\sqrt{1+\zeta}\,e+e^2$.
\end{proposition}
\begin{proof}
Let $x=R_Te_v$, $h=\widetilde x-x$, and $h_v=\widetilde x_v-x_v$.
Active inverse equations give
$\|h\|_{A_0}\leq\rho/\sqrt a$,
$\langle h,x\rangle_{A_0}=h_v$, and
$\|x\|_{A_0}^2=x_v$. Therefore
\[
 \left\|\frac{\widetilde x}{\widetilde x_v}
              -\frac{x}{x_v}\right\|_{A_0}^2
 =\frac{\|h\|_{A_0}^2-h_v^2/x_v}{\widetilde x_v^2}
 \leq\frac{\rho^2}{a\widetilde x_v^2}.
\]
Write $P=I-e_vx^\top/x_v$ and $E=\widetilde Y-Y$.
Subtracting the exact update yields
\[
 \widetilde Y'-Y'=EP+
 \widetilde Y[:,v]\left(x/x_v-\widetilde x/\widetilde x_v\right)^\top+D.
\]
Proposition~\ref{prop:maintainedsketch} contracts $EP$ in energy.
The energy norm of an outer product is the product of the
Euclidean norm of its left vector and the energy norm of its
right vector. The triangle inequality proves
\eqref{eq:sketchresidual}. All new errors remain supported on
unselected coordinates, so the argument iterates. Applying the
coordinate-representer bounds after
Proposition~\ref{prop:maintainedsketch} proves the moment bounds.

For higher powers, subtract the exact chain equation at each
level and multiply by $A_U^{-1}$, whose norm is at most $1/a$
by principal-submatrix interlacing. This proves the $E_\ell$
recurrence. The true target squared norms are at least
$B^{-2\ell}$ and $B^{-(2\ell-1)}$, respectively. On the ideal
JL event, their sketched norm is at most $\sqrt{1+\zeta}$
times the true target norm. Expanding a perturbed squared norm
and bounding its cross term proves the final error expression.
\end{proof}
\noindent Numerical check:
\texttt{scripts/verify\_t12\_residual.py}.

\paragraph{Instantiating numerical budgets.}
For example, a $p=1$ target moment tolerance $\bar\zeta<1$
is met by an ideal-sketch tolerance $\zeta=\bar\zeta/2$ and
$e_1,\sqrt B e_2\leq\bar\zeta/8$; substitution in the preceding
bound is at most $\bar\zeta$. Initial solve errors, subsequent
column residuals, and rounding terms share these budgets.
The same substitution applies to the higher-power normalized
errors $e$. Refreshed scalar entries use the T11 decrement
intervals and refresh-error bounds; if their resulting relative
marginal error is at most $\nu$, replacing exact values throughout
the hybrid proof loses at most the additional factor
$(1-\nu)/(1+\nu)$, with key tolerance enlarged to
$\max\{\eta,\nu\}$. This follows by comparing the largest
computed processed value with the largest true processed value
before applying the two stopping arguments. Residual evaluation
for every higher-power level costs $O(dL(m+N))$ per checked
state and is charged separately when used. These statements
specify numerical conditions, in addition to the sketch event,
for an inexact implementation of the preceding algorithms.

\begin{thebibliography}{10}
\bibitem{uehara1996} R. Uehara,
``NP-complete problems on a 3-connected cubic planar graph and
 their applications,'' Technical Report 4, 1996.
\url{https://www.jaist.ac.jp/~uehara/pdf/triangle.pdf}.
\bibitem{li2019} H. Li, R. Peng, L. Shan, Y. Yi, and Z. Zhang,
``Current Flow Group Closeness Centrality for Complex Networks,''
WWW, 2019; arXiv:1802.02556, Sections 3--5.
\bibitem{ss2011} D. A. Spielman and N. Srivastava,
``Graph Sparsification by Effective Resistances,'' SIAM J. Comput.,
2011; arXiv:0803.0929, Section 4.
\bibitem{kyng2016} R. Kyng and S. Sachdeva,
``Approximate Gaussian Elimination for Laplacians: Fast, Sparse, and
Simple,'' FOCS, 2016; arXiv:1605.02353.
\bibitem{wilson1996} D. B. Wilson,
``Generating Random Spanning Trees More Quickly than the Cover Time,''
STOC, 1996.
\bibitem{xia2025} H. Xia and Z. Zhang,
``Fast Maximization of Current Flow Group Closeness Centrality,''
arXiv:2504.04472v1, 2025, Sections III--IV.
\bibitem{zhou2023} X. Zhou, R. Wang, W. Li, and Z. Zhang,
``Maximizing the Smallest Eigenvalue of Grounded Laplacian Matrix,''
arXiv:2110.12576v4, 2023.
\bibitem{wang2026} G. Wang, H. Xia, R. Zhang, and Z. Zhang,
``Fast Algorithms for Group Markov Centrality Optimization,''
KDD, 2026, pp. 1412--1423, doi:10.1145/3770854.3780201.
\bibitem{fastfrac} ``FastFrac: Fast Edge Recommendation for Fractional
Spectral Traces on Graphs,'' provided manuscript, Section IV.
\bibitem{polar} ``POLAR: Fractional Spectral Link Recommendation via
Pole-Shared Krylov Sketching,'' provided manuscript.
\end{thebibliography}
\end{document}
